% !TeX program = lualatex
% One-source manuscript. The original Nous PNG follows the end of the document.
% Run LuaLaTeX three times in a writable directory. The release uses TeX Live 2026.
% For stable timestamps, set SOURCE_DATE_EPOCH=1795694400 and FORCE_SOURCE_DATE=1.
\documentclass[12pt,letterpaper]{article}
\usepackage{iftex}\RequireLuaTeX
% Suppress the path-dependent PDF trailer ID for reproducible release builds.
\pdfvariable suppressoptionalinfo 512\relax
\usepackage[hmargin=1.15in,vmargin=1.04in,headheight=16pt,headsep=14pt,footskip=28pt]{geometry}
\usepackage{fontspec}\setmainfont{Latin Modern Roman}\setsansfont{Latin Modern Sans}
\usepackage{amsmath,amssymb,mathtools,booktabs,tabularx,array,longtable,graphicx,microtype,xcolor,fancyhdr,enumitem,hyperref,bookmark}
\definecolor{ink}{HTML}{183831}\definecolor{accent}{HTML}{095D58}\definecolor{mute}{HTML}{455E55}
\hypersetup{colorlinks=true,urlcolor=accent,linkcolor=accent,pdftitle={Frobenius partners and spacing},pdfauthor={Justin Candler and Uriel}}
\setlength{\parindent}{0pt}\setlength{\parskip}{5pt plus .5pt}\setlength{\emergencystretch}{3em}
\linespread{1.14}\raggedbottom\pagestyle{fancy}\fancyhf{}\fancyhead[L]{\small\sffamily\color{mute}Nous Independent Research}\fancyhead[R]{\small\sffamily\color{mute}Frobenius partners and spacing}\fancyfoot[C]{\small\sffamily\thepage}
\renewcommand{\headrulewidth}{.3pt}\setcounter{tocdepth}{1}\setlist{itemsep=2pt,topsep=4pt}
\newcolumntype{Y}{>{\raggedright\arraybackslash}X}

\directlua{
 local source = assert(io.open(status.filename, "rb")); local rows={}; local take=false
 for line in source:lines() do
 if line:match("^BEGIN FIGURE nouslogo") then take=true
 elseif line:match("^END FIGURE") then take=false
 elseif take then table.insert(rows,line:match("^[0-9a-f]+")) end
 end
 source:close()
 function nousLogo()
 local hex=table.concat(rows); assert(string.len(hex)>100000)
 local bytes=hex:gsub("..",function(pair) return string.char(tonumber(pair,16)) end)
 local path=tex.jobname.."-embedded-nous.png"; local existing=io.open(path,"rb")
 if existing then local previous=existing:read("*a"); existing:close(); assert(previous==bytes)
 else local file=assert(io.open(path,"wb")); file:write(bytes); file:close()
 function cleanNousLogo() os.remove(path) end end
 node.write(img.node(img.scan{filename=path})) end
}



\DeclareRobustCommand{\id}[1]{\mbox{\small\ttfamily #1}}
\newcommand{\old}[3]{\href{https://github.com/nousentllc/riemann-hypothesis/blob/40e56bab026552e089d7afcfeff08d2defb9c005/src/#1\#L#2}{#3}}
\newcommand{\newsrc}[3]{\href{https://github.com/lsernau/grh-lean/blob/72fead102f663d5c7a407964a13dc0cf9d977207/grh-compact/src/#1\#L#2}{#3}}
\hypersetup{pdfsubject={A source-level comparison of retained orbits in two Lean formalizations}}
\begin{document}
\pagenumbering{roman}
\thispagestyle{empty}
\begin{center}
\vspace*{.52in}
{\sffamily\small\bfseries\color{accent}NOUS INDEPENDENT RESEARCH}\par\vspace{.65in}
{\fontsize{31}{37}\selectfont\bfseries\color{ink}Frobenius partners\\and spacing}\par\vspace{.27in}
{\Large What the original retained set can certify\\and what the compact Lean proof requires}\par\vspace{.35in}
\resizebox{!}{1.12in}{\directlua{nousLogo()}}\par\vspace{.47in}
{\large Justin Candler with Uriel, AI research assistant}\par
26 November 2026
\end{center}
\vfill
\begin{center}
\small\sffamily
Original Lean source \texttt{40e56bab}\quad\textbar\quad Compact Lean source \texttt{72fead10}
\end{center}
\clearpage
\section*{Abstract}

Can the original selection rules force the same unique short Frobenius partner as the compact Dirichlet GRH proof? At the fixed comparison value \(d=9239\), the original selector has no admissible scale-tied rank input. Its width formula implies \(d>94897\) under the literal \id{RankVal} and host hypotheses. This is a stronger obstruction than the earlier bound on the number of signed indices and applies at every old scale.

For a general original host, we derive two consequences on its given retained set. A good short window with cap \(K\) admits at most \(K-1\) odd offsets from a retained anchor, and a marked tile excludes an odd displacement with both endpoints inside it. These conditions leave room for multiple short partners. In a 16-cycle, the retained set \(\{2,6,10,15\}\) satisfies all three original local clauses with an active marked tile. From anchor 15 the odd exponents 3 and 7 reach different retained targets. Moving 15 to 14 preserves the tile, the size, and both occupancy caps, and removes all odd partners. An independent enumeration of 816 completions at order 32 locates the difference when the short cap is changed.

The compact proof excludes this collision by applying its three-term condition to mixed powers in a dyadic quotient and then ruling out a small order. It obtains at most one candidate in the small-prime branch, with a separate existence test. We replay the native original implications and exact controls in Lean and check the compact arithmetic guard and its algebraic use. The finite examples establish an implication boundary for the old local rules. The scope of the local replays is distinct from full builds of the upstream proofs.

\medskip
\textbf{Human--AI collaboration.} Justin Candler directed the research program, chose the fixed comparison and asked whether the same retained set forces uniqueness. Uriel/Codex located and compared the Lean declarations, developed the deductions and controls, ran the replays, and drafted this manuscript. The upstream formalizations and the external verification report are credited in the sources.
\clearpage
\begingroup\small\setlength{\parskip}{0pt}\tableofcontents\endgroup
\clearpage
\pagenumbering{arabic}
\section{Findings and scope}

At $q_0=18480$, $d=9239$, there is no original selector witness satisfying the literal rank input \id{RankVal}. The obstruction is independent of the scale and auxiliary prime. It occurs before the selector produces a retained set. At other parameters the original proof does have an existence theorem for its host and retained set. On any such host the old orbit conditions bound the number of nearby odd partners and exclude odd displacements internal to a marked tile. The compact proof's unique-eligible-partner condition is stronger. A finite mask passes every original local clause, including a complete marked tile, and admits two partners of the same retained anchor in one checked short window.

Our preceding study showed that the carrier of a hypothetical original host with $d=9239$ has at most 32 signed indices. That implication is compatible with the new finding. The present argument adds the identity required at the selector's entry point. That identity makes the input impossible at every old scale, so there is no selected original set to recover at the fixed parameters. The separate readable datum and its selected places remain the objects on that side of our earlier comparison.

This study distinguishes three kinds of evidence. Literal declarations from the original archive are replayed under its own Lean toolchain against prepared imports. They establish the fixed-prime obstruction and conditional theorems on the given retained set. Exact controls exercise the local finite predicate, including a counterexample and a one-coordinate control with the same tile, size, and caps. The compact source trace identifies the at-most-one mechanism in the small-prime branch. A second Lean replay checks its group-theoretic implication on abstract types. The verification ledger specifies where the source trace supplies hypotheses and where a new theorem consumes them.

\subsection{Notation and research questions}

$q_0$ in the first sentence is the modulus from our readable experiment. The original host has an auxiliary rational prime $H.q$ and a distinct $H.d$. Its eventual host theorem also uses the name \id{q0} for an auxiliary level prime. The new obstruction only uses $H.d$. We write $\mathcal S$ for the original scale record, $x$ for $\mathcal S.x$, and $s$ for the compact dyadic tower exponent. The rational prime whose local behavior we study is $p$.

\begin{center}
\begin{tabularx}{\linewidth}{@{}p{.22\linewidth}Y Y@{}}
\toprule
Question & Original selector & Compact small-prime packing\\
\midrule
What is fixed? & Host, Frobenius shift, and the retained finite set carried by that host. & Parameter record, prime, and the set \id{D.thick}.\\
What is tested? & Window counts, occurrence of marked tiles, weighted count of failed orbits. & Mixed powers for triples of selected places, with a dyadic order bound.\\
What is a partner? & A function of an anchor and exponent. A retained-target test can be imposed on its output. & A candidate short odd exponent with a target in the same selected set.\\
What follows locally? & Window occupancy bound and parity exclusion within a marked tile. & At most one eligible exponent at an anchor in the small-prime branch.\\
Where does existence enter? & A probabilistic selection theorem produces some retained set. & An existence case analysis precedes building a packing from an eligible exponent.\\
\bottomrule
\end{tabularx}
\end{center}

A native theorem $H.d=9239\Rightarrow\neg\id{RankVal}\,\mathcal S\,R$ decides the same-host hypothesis. A positive full tile count, legal short and long counts, and two distinct targets within one checked short window decide the local implication hypothesis. The mechanism hypothesis asks what can be carried between the proofs. Its positive evidence is a theorem on the given retained set. Its negative evidence is a concrete missing premise, or a micro-instance that fails when a condition is removed.

Fixing the set matters. Substituting a newly chosen mask for the host's field \id{TT} changes the question. Proving uniqueness after silently assuming an orbit exists can hide an empty prime filter. A finite counterexample to the old local clauses cannot establish that a global arithmetic witness exhibits the same collision. The experiments here are designed to expose exactly these distinctions.

\section{How the original witness reaches a local orbit}

\subsection{From parameters to a retained set}

The original host \id{Host R} is a dependent record over a frame and rank tuple. Its field \id{TT} is a finite subset of the unsigned group $R.A$. The top places are the images of retained indices under \id{R.tau1}. This explains why the previous study counted a signed carrier of size $2|R.A|$ while this comparison starts with an unsigned orbit. The sign still matters when interpreting the coordinate consumed by a local gain.

The theorem \id{topsChoiceH} asks for old ranks satisfying \id{RankVal}, a host meeting \id{HostConst}, and a paired arithmetic certificate involving \id{PairWin} and \id{PairRow}. It gives a scale threshold and a constant $C_b$, followed by the existence of a finite set $T$ at each admissible input. The conclusion uses the same $T$ in three conditions:
\begin{align*}
  \tfrac12\delta_1 R.n &\le |T|\le 2\delta_1 R.n,\\
  \id{BadOK}\,(\mathcal S,C_b,H[T])&,\qquad
  \id{Shares}\,(T,x^{-1+\varepsilon}).
\end{align*}
Here $H[T]$ is the old host with its retained-set field updated to $T$. The prime labels are tied to the rest of $H$. \id{hPH H p} takes the unsigned coordinate of a chosen arithmetic Frobenius element using a classical selection. The definition mentions the field, the Galois representation of the host, and the prime. Updating $T$ leaves \id{hPH} unchanged. Earlier weighted-orbit experiments used precisely this invariance to vary masks while holding the prime step fixed. Those alternative masks remain controls.

The base record \id{baseOf} passes $H.\mathrm{TT}$ to \id{Arith.Loc.Good}. The weighted bad-mass objective is summed over primes satisfying the range, inert-class, and long-order tests. For each such prime and retained point, the objective makes one unit of incidence if the orbit of that point fails its Good test. The summand is $\log p/p$ times the filtered cardinality. It counts every retained member of a failed orbit. A good orbit contributes zero, irrespective of its number of marked tiles or odd offsets. This last property is a structural reason the score cannot distinguish the local examples constructed below.

\subsection{Orienting the same Frobenius step}

The predicate \id{IsResidueClassZero} in the relevant \id{Arith.Base} is equality to \id{inertClass}. The signal predicate is equality to a distinct constructor. Writing $h=\id{hPH}\,H\,p$, we expand the two old declarations as
\[
 \id{stepP}(p)=
 \begin{cases}2h&\text{signal class},\\h&\text{otherwise},\end{cases}
 \qquad
 \id{partner}(t,e)=
  \bigl(\mathbf1_{\text{signal class and odd }e},\ t+eh\bigr).
\]
The indicator is Boolean, and false is the sign for tops. Thus on the inert class used by \id{BadOK}, the selector and consumer have the same unsigned step. The sign of the consumer's partner is false for every $e$. A retained-target query in this class is precisely $t+eh\in H.\mathrm{TT}$ on an already retained $t$. The theorem \id{inert\_step\_and\_partner} checks this on the original types. No assumption about the order of $h$ is necessary for that identity.

This equality concerns the old selector and old consumer. The compact place group and its class map are different objects. The original \id{hPH} supplies no evaluator for the compact place class. Comparing the cyclic actions requires choosing what counts as a forward step. Reversing a step sends $e$ to $n-e$ within a period of length $n$. Reversal can change which exponents lie in a short horizon. Thus orientation cannot be disregarded even when the underlying unlabelled orbit is the same.

\subsection{Good orbits and the local consumer}

Choose an orbit representative $b$. The source defines \id{topsPos} as the finite set of positions $i\in\operatorname{Fin}(\operatorname{addOrderOf}h)$ for which $b+ih\in T$. Its elements correspond to the same old retained set. The map is injective on reduced orbit positions because two nonnegative integers less than the additive order give different points. The representative \id{Arith.bsH} uses an ordering of the full frame. It is not an instruction to find the nearest retained point. Local parity is parity of a coordinate relative to that representative.

The local gain uses these data in more than one way. The Good test filters the orbit representatives in \id{goodReps}. \id{Cand} checks prime, inert class, scale, order, and an appropriate band exclusion. The representatives \id{VV} additionally require a bound on the weighted orbit contribution \id{Sig}. On a used orbit, the gain compares a direct alternative against an expression involving \id{gU} and \id{Pr}. The latter uses a sum over odd exponents and a minimum against a capacity. The direct branch uses an odd exponent test and a condition \id{Excl} that can discard a proposed interaction. Verifying a unique target for each given exponent is automatic because the map is a function. The compact theorem under comparison instead constrains the feasible exponents before a packing is chosen.

This trace explains why the old comparison predicate contains both anchor and target membership in $T$. It must also name a horizon. The old sums may range over infinitely many exponents with decaying coefficient $p^{-e}$. Membership repeats after an orbit period, while the coefficient changes. Multiple exponents can therefore occur even when they denote the same target. Our decisive fixture keeps the two odd exponents below one period and inside one checked window. The two positions are distinct. That makes this a geometric difference, without multiplicity from winding around the orbit.

\section{The fixed discriminant cannot enter the old selector}
\label{sec:rank}

\subsection{The input identity}

The preceding companion study used the fields in \id{Host}, including
\[
     4 R.t R.r+1<H.d,\qquad R.r+5\le R.t
\]
to bound the signed indices. The frame's dyadic restriction gave $o=1$, $|R.A|\le16$, and at most 32 signed slots at $d=9239$. A one-coset share check could therefore be automatic while faithful enumeration of the readable places was impossible. That conditional result did not use the scale-tied equality in \id{RankVal}. The present question is whether the old selector can ever take that host with that fixed $d$.

The last part of \id{RankVal} is
\begin{equation}
 R.r+R.t=8 R.f R.r+4+\operatorname{TOv}(\mathcal S)
                            \operatorname{nWv}(R,\mathcal S).
 \label{eq:rank-identity}
\end{equation}
Write the source terms as
\begin{align}
 L&=x^\varepsilon,\qquad
 \operatorname{TOv}(\mathcal S)=\lceil x^{377}\rceil,\nonumber\\
 W&=\operatorname{nWv}(R,\mathcal S)
  =8\left\lceil\frac{10r}{8}
  \left(f+\frac{f^2(k-1)^3L^2}{m}\right)\right\rceil .
 \label{eq:width-cost}
\end{align}
Ceilings in Lean return natural numbers. The term \(\operatorname{TOv} W\) is required by the rank formula before the retained set is supplied. In qualitative terms the ranks need room for a width term while the host insists the product $tr$ fit below the discriminant. This is a necessary resource bound on the host construction. It operates even at the weakest limits $x>1$ and $\varepsilon>0$.

\subsection{Native derivation}

Write $m=R.m$, $k=R.k$, $f=R.f$, $r=R.r$. The old fields give
\[
 m\ge2,\quad k\ge m+3,\quad f=2R.n\ge m,\quad r=k+m-1\ge6.
\]
In particular \(k-1\ge4\), \(k-1\ge m\), and \(L=x^\varepsilon\ge1\). Thus
\[
 \frac{f^2(k-1)^3 L^2}{m}
 \ge \frac{4\cdot16\cdot m}{m}=64.
\]
The parentheses in \eqref{eq:width-cost} are at least $66$. Multiplication by $10r$ and upward rounding give $W\ge660r$. Since $\operatorname{TOv}\ge1$, even the last summand of \eqref{eq:rank-identity} gives $t+r\ge660r$. In turn
\[
 t\ge659r\ge3954,\qquad 4tr+1\ge4\cdot3954\cdot6+1=94897 .
\]
The native field \id{Host.qd\_big} gives $94897\le4tr+1<d$. This contradicts $d=9239$. Every inequality makes the certificate a necessary condition. There is no implication that $d=94898$ yields a host. Several constraints that could raise the actual minimum were deliberately left unused.

The same calculation also bounds the other host prime when the first component of \id{qd\_big} is used. The executable theorem needs only the second. Changing the auxiliary prime leaves the $d$ contradiction in force. The readable modulus $q_0$ remains meaningful in its own proof and does not enter \eqref{eq:rank-identity}.

A stronger written bound can preserve terms discarded in the certificate. The identity implies $r+t\ge16r+4+660r=676r+4$, so \(t\ge675r+4\). Using $x>1$ gives $\lceil x^{377}\rceil\ge2$ and raises this to \(t\ge1335r+4\). The native result uses the coarser $94897$. A more optimized bound would be relevant for searching genuine original instances at larger $d$. For this fixed case the slack is already large.

\subsection{What the impossibility does and does not assert}

The theorem says every input \(\mathcal S,R,H\) with \(H.d=9239\) violates \id{RankVal}. It does not refute the old theorem that produces hosts under its actual linked, eventual hypotheses. Those hypotheses let $d$ vary. Nor does it assert that a bare \id{Host R} with $d=9239$ is uninhabited. It rules out the conjunction with \id{RankVal}, which is precisely the selector input.

At these parameters the original retained set we wanted to compare cannot come from that selector. The obstacle is a mathematical contradiction, so a bigger search budget cannot produce this witness. For general actual hosts, theorems in the next sections apply to the given \id{TT}. They make the needed hypotheses on an orbit explicit. Separate fixture sets test the strength of those hypotheses.

\section{What the old clauses prove on the retained set}
\label{sec:old}

\subsection{Two window lengths and a tile test}

Let $n=\operatorname{addOrderOf}h$, and let $U\subseteq\{0,\ldots,n-1\}$ be the orbit coordinates of $H.\mathrm{TT}$ relative to the representative $b$. For nonnegative integers $i,L$ put
\[
  W_n(i,L)=\{(i+j)\bmod n:0\le j<L\}.
\]
The original code stores $W_n(i,L)$ as a finite image. Repeated residues do not accumulate multiplicity when $L>n$. The short clause \id{CondA} requires \(|U\cap W_n(i,L_2)|\le K\) for each $i<n$. The long clause \id{CondIntervals} makes the analogous assertion with length $G'$ and cap $N_1$. In the old \id{Good} these are specific functions of the scale and rank. In particular \(K=\lceil 8(h_G+2)/\varepsilon\rceil\) with $h_G=2+8\varepsilon$, and \(r_G=K+1\). The short length $L_2$ is twice a supremum of a scale-dependent logarithmic length indexed by inert primes. It need not equal the compact $s$. Neither count singles out odd offsets.

For the middle clause let $u=\lceil1/\delta_1\rceil$ be the cell width and set $A=(6r+1)u$. A tile at index $b_0$ occupies $[b_0 A,(b_0+1)A)$. Its cells are indexed by $c=1,\ldots,6r+1$. Cell $c$ is
\[
  [b_0 A+(c-1)u,b_0 A+cu)\cap \mathbb N .
\]
A marked tile has no retained point in each odd-indexed cell and exactly one retained \emph{even absolute coordinate} in each even-indexed cell. Hence exactly $3r$ points of a marked tile are retained. The look-ahead clause requires a marked tile within a prescribed number of complete tiles forward, with cyclic indexing among complete tiles. Leftover positions beyond the final complete tile are outside the tile partition, although the two cyclic window tests still see them. This difference between the two grids is central to the small counterexample.

When $n<A$ there are no complete tiles. The middle clause then quantifies over no start. A counterexample with such a small orbit would leave unanswered whether the clause can matter when active. We require a complete marked tile and, in the smallest decisive example, zero look-ahead. The choice eliminates vacuity and the freedom to satisfy the tile condition using an arbitrarily remote tile.

\subsection{Short-window occupancy bounds the degree}

For a retained $t=b+ih$, define
\[
  J(T,t,h,\ell)=\{e\in\mathbb N:e<\ell,\ e\text{ odd},\ t+eh\in T\}.
\]
Use $0<\ell\le\min(n,L_2)$. The map $e\mapsto(i+e)\bmod n$ is injective on the union $\{0\}\cup J$. Indeed, equal residues give $e=f\bmod n$ by cancellation of $i$, and both offsets lie in $[0,n)$. The image sits in $U\cap W_n(i,L_2)$. Every offset $e\in J$ gives the required membership by construction, including when $i+e$ wraps. The image of zero lies in that intersection because $t$ is retained. Since zero is even, this adds a distinct point. Thus
\begin{equation}
       |J(T,t,h,\ell)|+1\le |U\cap W_n(i,L_2)|\le K .
       \label{eq:oldcap}
\end{equation}
This is a theorem for the actual original types. \id{old\_good\_orbit\_offset\_cap} takes the original base, its step $h$, the representative $b$, a proof of \id{Arith.Loc.Good B h b}, and an anchor with proof of membership in the base's \id{TT}. No second retained set appears. Substituting the base made from $H$ and the step \id{hPH H p} gives the condition on that host. The proof never asks an oracle to enumerate $H$ or $T$.

The source proof of \eqref{eq:oldcap} is worth reconstructing because errors in this comparison usually come from a silent change of representative or period. The Lean object bounding visits is a \id{Finset} of natural exponents. We insert $0$ before mapping it to the window. The membership of a positive image is read using \id{topsPos}, whose type is a finite subtype carrying the proof $i<n$. To pass from a wrapped index to the underlying group element we use $(i\bmod n)h=ih$. This follows by writing $i=(i\bmod n)+n\lfloor i/n\rfloor$ and using $nh=0$. The image of the exponent set under the residue map then has exactly the cardinality of the exponent set, by injectivity on the chosen domain. The final step uses the first conjunct of the native Good predicate.

If $K\le2$, \eqref{eq:oldcap} implies at most one odd offset. Even then a partner need not exist. In the source’s usual scale regime $K$ is much larger. Since $\varepsilon<1/52$ and $h_G>2$, the expression inside the ceiling is greater than $32\cdot52=1664$. Thus the statement using a small cap is a conditional comparison and a controlled ablation. It does not silently substitute a tuned threshold into the original proof. The main contrast is between a cap that is much larger than two and a group-theoretic collision rule that excludes the second target directly.

\subsection{Where the marked tile actually forces spacing}

The marked condition does give an exact zero statement. If $z$ is retained inside a marked tile and $e$ is odd, then $z+e$ cannot be retained while both lie in that same tile before reducing modulo the orbit length. In fact every retained coordinate in a marked tile is even. To derive this from the definition, write $j=\lfloor(z-b_0A)/u\rfloor$. The hypothesis $u>0$ and the tile endpoints put $z$ in the unique cell $j+1$. If that cell index were odd, the empty-cell condition would contradict retention. If the index is even, the singleton condition identifies $z$ with its even representative. Applying this argument to both endpoints makes an odd displacement impossible. Lean verifies both the cell location and the parity argument against the original \id{Marked}.

The conclusion extends to any portion of the ray that stays inside one marked tile. It gives no restriction when an odd step starts outside the tile. The location of a marked tile is existential for each start, so existence of a marked tile ahead is compatible with odd retained coordinates before that tile. On a complete even orbit, an odd offset flips parity even when the residue wraps. This provides another equivalent description of the fixture once $n$ is known even. The formal internal-tile lemma uses a nonwrapped integer comparison, so it also avoids making any argument that would accidentally require an even orbit.

The difference between this theorem and compact uniqueness is visible directly in their logical forms. The tile lemma has a condition on both integer endpoints inside a particular interval. Its exclusion follows from parity alone. The compact source has a collision condition on two possible targets indexed by powers of the same step, with a bound on their exponents. It requires no marked interval. Any transfer from the tile lemma to the compact claim would need a further theorem placing the relevant endpoints inside that parity-controlled portion. The supplied Good clauses have no such universal placement condition.

\subsection{At most one and a certified one}

An existence claim has an additional premise under either proof strategy. For the old rule define the degree of the anchor as $|J(T,t,h,\ell)|$. If the degree is zero the candidate query has no answer. If it is one then it gives a unique answer relative to $h$ and $\ell$. If it exceeds one the map of exponent to target remains injective below one order, and each target is distinct. Lean's count theorem gives an upper bound; it does not assert the degree is positive. A conditional statement that an old orbit agrees with the compact at-most-one guarantee needs neither prime reconstruction nor a choice of a target, but an assertion of the same \emph{unique} partner would need both nonemptiness and a verified map between the two labelled types.

At a tiny horizon the question is trivial. For $\ell\le3$ the only possible positive odd exponent is 1. At most one follows from arithmetic on exponents alone. This demonstrates why a comparison has to report the horizon. Reducing it to a single exponent would vacuously obtain the desired form of answer yet provide no evidence of the spacing property used by the compact small-prime packing.

\section{What the compact proof uses}
\label{sec:new}

\subsection{An algebraic spacing law}

Let $\Phi$ be the compact multiplicative group of places and $T=\id{D.thick}$. Let $\pi$ be the compact class projection into the units of \(\mathbb Z/2^{s+1}\mathbb Z\). The field \id{D.three\_term} says that for appropriate positive $i<j\le324s$ and any three elements $\sigma,\tau,\rho$ of $T$,
\[
     \pi(\tau/\sigma)^j=\pi(\rho/\sigma)^i
                  \quad\Longrightarrow\quad \pi(\tau/\sigma)=1 .
\]
The condition is on projected displacement, not merely on an arithmetic progression with consecutive exponents. The exponents can be any pair in range. This matters when the source tries to exclude two competing odd offsets, which can be separated by more than two. The projection is to a modular unit group. Collisions caused by torsion are handled by a further order estimate.

The constructor carries additional requirements on the thick sets. The congruence condition fixes their class to $-1$ modulo the nondyadic factor used in the data, and the three-term condition acts on the dyadic projection. The latter is applied while the former makes some primes or shifts unable to stay in the selected class. Thus the two fields do different work. Having enough places to feed the density coefficient of the final error bound does not establish either spacing property by itself.

\subsection{Trial of two odd targets}

Fix a prime $p$ in the compact small-prime branch, $p\le D.k$. Write $g$ for the place class of $-p$ where it is available as a unit. The local candidate is called \id{Q} inside \id{Data.exists\_isFrobPacking}. For anchor $\sigma$ and exponent $e$ it requires
\[
  e<s,\quad e\text{ odd},\quad p\ne P.d,\quad
  \sigma\in T,\quad g^e\sigma\in T .
\]
The source writes the last term using \id{D.placeOfInt} and the compatibility of powers. For odd $e$, \( (-p)^e=-p^e\), which is the sign convention used by the packing. At even $e$ those two integers differ. The identity is invoked only after the parity test.

Suppose both $i$ and $j$ are candidates with $i<j$. Since they are odd they are positive. Let $\tau=g^i\sigma$ and $\rho=g^j\sigma$. Their projected displacements satisfy
\[
 \pi(\tau/\sigma)^j=(\pi(g)^i)^j=\pi(g)^{ij}
                              =(\pi(g)^j)^i=\pi(\rho/\sigma)^i .
\]
The equality is automatic. All information lies in the fact that both targets were selected and that \id{three\_term} applies to this triple. It forces $\pi(g)^i=1$. That equation conflicts with the small-prime size bound used in the source. This is where the compact proof excludes the second candidate.

A generic Lean lemma \id{two\_targets\_conflict} formalizes this implication over a multiplicative homomorphism into a commutative group. It keeps the hypotheses separate. The mixed-power identity is also checked as an unconditional theorem for a monoid. This factorization helps interpret the old counterexample. If two exponents of a common old displacement reach retained points, they also satisfy the mixed-power identity. The difference is which conclusions may be drawn about the projected order.

\subsection{The order bound and its exact role}

The compact arithmetic lemma \id{two\_pow\_le\_of\_pow\_eq\_one} states that if an integer $a$ has $|a|>1$, $i>0$, and \(a^i=1\) modulo \(2^{s+1}\), then
\[
                   2^s\le (|a|+1)i.
\]
At $a=-p$ the right side is $(p+1)i$. This is substantially stronger than estimating the integer magnitude of $(-p)^i-1$ directly. The latter would grow exponentially with $i$. The proof uses divisibility at 2 to get a bound linear in $i$. The compact source then takes $p\le D.k$ and a short bound on $i$, invoking the record's parameter inequalities to contradict it. No enumeration of actual place labels enters. The law excludes the collision uniformly on the supplied set.

It is important that the source checks both conditions. A three-term implication on a quotient alone can be true because the quotient collapses every displacement to identity. In that case it cannot distinguish two different points of the full orbit. Conversely a projection with large order gives no uniqueness if both targets are allowed without a spacing condition. These are distinct failure modes, illustrated later by exact finite controls. They also explain why a weaker argument based solely on the size of the old cyclic group does not recreate the compact proof.

For the $p>D.k$ branch the compact source constructs the packing with exponent 1. That branch follows a separate quotient and valuation argument. A constant exponent in the packing does not say that no other odd offsets connect members of the retained set. The stronger candidate-count result belongs to the $p\le D.k$ branch. Both branches enforce the capacities of their packing. Keeping their statements separate prevents taking a property of the witness selected by one branch as a structural property of all possible edges.

\subsection{From uniqueness to a packing}

Inside \id{exists\_isFrobPacking} the source proves the key implication in a local term called \id{N}. It sets up the candidate predicate \id{Q}. It then splits on existence of an $e$ satisfying \id{Q} at each anchor. Where a candidate exists the chosen exponent is unique and allows an amount bounded by \(n_\sigma/p^e\) using natural division. The source still checks a capacity at a target when it combines outgoing and incoming contributions. It also needs to rule out conflicts involving shorter odd arrivals. The same spacing mechanism is used with another anchor after rewriting composition of Frobenius offsets. This is why the proof needs a condition stable under choices of anchor in the set.

The theorem's external interface packages existence of numbers \(i,y\), the abstract \id{IsFrobPacking} conditions, and a bound on divisibility loss. It does not export a list of actual compact partners. Our compact-side Lean lemma checks the structural implication behind the local term. The fact that the pinned compact proof uses it at the stated parameters is a declaration trace of the existing source. A complete verification of the original local module from scratch would be a broader replay than the separate abstract lemma.

\subsection{Two distinct notions of collisions}

The first notion concerns a fixed source with two different admissible outgoing steps. It is controlled by the comparison just discussed. The second concerns two sources whose steps arrive at a common target. If their exponents agree, multiplying by a common group element is injective and the two sources agree. If the exponents differ, the source considers a shorter arrival relative to the chosen outgoing exponent in its proof of \id{IsFrobPacking}. The formulas treat that case using more than one membership premise and rebase at a source. A unique outgoing edge everywhere does not by itself imply a unique incoming edge. This matters for comparing local products and capacities across editions. Our finite question isolates the outgoing predicate because one anchor with two qualifying offsets is sufficient to refute an implication from the old local clauses.

The proof relationship can be summarized as an implication chain:
\[
 \text{two targets}\ \Longrightarrow\
 \text{mixed powers agree}\ \Longrightarrow\
 \text{small dyadic order}\ \Longrightarrow\
 \text{scale contradiction}.
\]
The first arrow is general algebra. The second is the distinctive meaning of \id{D.three\_term}. The third comes from the compact parameter bounds and the arithmetic lemma. Checking each arrow separately is more informative than counting theorems or concluding from the labels of the definitions that the two systems do the same thing.

\section{Exact controls with active marked tiles}
\label{sec:controls}

\subsection{The 16-cycle certificate}

The two masks use the same cyclic group, step $+1$, and orbit representative $0$. The literal input to \id{GoodOrbitIntervalTileCondition} uses
\[
 n=16,\quad L_2=8,\quad K=3,\quad r=1,\quad u=2,\quad
 M_0=0,\quad G'=16,\quad N_1=4 .
\]
The tile length is 14. Put $T=\{2,6,10,15\}$. The complete tile is $[0,14)$, containing seven cells of width two. Its three even cells contain the singletons $\{2\},\{6\},\{10\}$. Its odd cells are empty. The look-ahead clause can only use that tile and passes with displacement zero. The last two coordinates, 14 and 15, are outside the complete tile.

For the short windows, an occupied arc is read by consecutive cyclic coordinates. The following row gives the exact number of retained points in $W_{16}(i,8)$ at each start. The second mask $Z=\{2,6,10,14\}$ changes only the final point.
\begin{center}
\small
\begin{tabular}{@{}l*{16}{r}@{}}
\toprule
$i$&0&1&2&3&4&5&6&7&8&9&10&11&12&13&14&15\\
\midrule
$T$&2&2&2&2&2&2&2&1&2&2&2&2&2&2&2&3\\
$Z$&2&2&2&2&2&2&2&2&2&2&2&2&2&2&2&2\\
\bottomrule
\end{tabular}
\end{center}
The short cap is 3 and holds for both masks. Every long window is the whole cycle, so its occupancy is exactly 4 for both. All three source clauses pass. In $T$, the retained anchor 15 reaches retained 2 at $e=3$ and retained 6 at $e=7$. Both are positive odd offsets below 8, and both are within a single short window starting at 15. The bound \eqref{eq:oldcap} is attained. It counts the anchor itself plus two partners.

A hand calculation of the short table can use the sliding recurrence
\[
 C_{i+1}-C_i=\mathbf1_{i+8\bmod16\in T}-\mathbf1_{i\in T}.
\]
The window starting at 0 sees 2 and 6. Substituting the four retained residues gives the entire row. The rotation changes the count only when an old endpoint is dropped or a new endpoint is added. This is also a check on the half-open window convention: a window of length eight covers offsets 0 through 7. In particular, the exponent 7 in the certificate is included. An off-by-one horizon of seven would omit that edge.

The control $Z$ has only even positions. Since the cycle has even length, all odd shifts lead to odd positions at every anchor. Its odd partner count is zero throughout. The marked tile is preserved exactly. The short windows and long windows still meet the same caps. The retained cardinality is four in both masks. The control therefore isolates geometric information that is absent from the cap inequalities.

There is also an exactly-one control inside the first mask. Over a full period, each of its even anchors 2, 6, and 10 has exactly one odd-offset partner, namely 15. The odd anchor 15 has three possible even targets. Accordingly a claim about one selected source cannot substitute for the universal statement over all selected sources. The short-window filter changes this picture further. At anchors 2 and 6 the sole odd target falls beyond the short horizon, while at 10 it occurs after five steps. The table below gives the reduced full-period and short exponents.
\begin{center}
\begin{tabular}{@{}ccc@{}}
\toprule
Anchor in $T$ & Eligible offsets $e<8$ & Eligible offsets $e<16$\\
\midrule
2 & $\varnothing$ & 13\\
6 & $\varnothing$ & 9\\
10 & 5 & 5\\
15 & 3, 7 & 3, 7, 11\\
\bottomrule
\end{tabular}
\end{center}

The old local predicate is checked using the exact imported definitions. The finite Lean proof converts each universal test on starts into a decision over \id{Fin 16}. For the tile clause it proves that the only complete tile is zero and checks its seven cells. The long window assertion is bounded by the cardinality of the mask, a proof that also works when the long window is enlarged. These kernel-checks are independent of the Python counts. In particular the central negative finding does not depend on a long search having sampled sufficiently many masks.

\subsection{A modular check on the missing implication}

The counterexample is combinatorial, but the failure of the three-term rule can be seen in a faithful dyadic modular control. Take the unit $g=51=-13$ modulo 64. Repeated squaring gives
\[
 g^2=41,\quad g^4=17,\quad g^8=33,\quad g^{16}=1\pmod{64}.
\]
Since any proper divisor of 16 divides 8, $g$ has order 16. Also $33$ is neither $1$ nor $-1$ modulo 64, so its class modulo the sign subgroup still has order 16. Label the positions of the cycle by powers of $g$. From the anchor 15, the relative displacements to 2 and 6 are $g^3=43$ and $g^7=27$, respectively. These have matching mixed powers because
\[
       (g^3)^7=g^{21}=(g^7)^3 .
\]
Yet $43\ne1$ modulo 64. Thus the old window and tile predicates, even on a cycle realized by units of a dyadic ring, allow a triple violating the mixed-power conclusion. This is an explicit failure of the hypothesis \id{three\_term} in the finite model. It is not an identification of this numerical unit with \id{hPH} for a full original host.

The order computation provides a clean ablation. A projection with trivial image would satisfy the mixed-power implication for every set because every displacement maps to 1. Both offsets would still be present in the underlying orbit. In the faithful projection above, the first displacement is nonidentity and the spacing conclusion fails. Hence the missing ingredient cannot be replaced by quotienting away the information needed to distinguish the targets. Lean verifies the congruences and inequality directly in \id{ZMod 64}.

A separate Lean lemma checks a simple order observation specific to odd exponents. If $\operatorname{ord}(g)$ divides a power of two and $i$ is odd, then $g^i=1$ implies $g=1$. The order divides $i$ as well. An odd integer is coprime to every power of two, so the common divisor is 1. This verifies the role of the odd parity for many modular examples without evaluating long powers. The compact source uses the quantitative dyadic bound in Section~\ref{sec:new}; that bound also controls more general exponents needed by related spacing arguments.

\subsection{Enumerating all completions of a fixed tile at order 32}

The 32-cycle experiment changes the short cap and holds the rest of the fixture design fixed. The complete tile at index 0 has width 14 and selected positions $\{2,6,10\}$. The second complete tile starts at 14. The look-ahead bound is 1, so every start among the two complete tiles can reach index 0. Each mask adds exactly three positions among $\{14,\ldots,31\}$. There are
\[
                         \binom{18}{3}=816
\]
such masks. The prescribed long window is the whole cycle, with cap six. Exact enumeration therefore changes selection only through the short cap. Every trial still calls all three checks, so a future change to the fixture definitions cannot silently bypass the marked clause.

A mask is called multiple if at least one retained anchor has two or more eligible odd offsets below eight. For each cap we also record the largest degree over anchors. The counts are:
\begin{center}
\begin{tabular}{@{}rrrrrrr@{}}
\toprule
Short cap & Passes & Rejected & Max 0 & Max 1 & Max 2 & Max 3\\
\midrule
2 & 335 & 481 & 57 & 278 & 0 & 0\\
3 & 802 & 14 & 92 & 462 & 248 & 0\\
4 & 816 & 0 & 92 & 464 & 259 & 1\\
\bottomrule
\end{tabular}
\end{center}
Every row sums. For cap two, \eqref{eq:oldcap} gives the zero multiplicity result independently of the enumeration. The surviving masks are not all unique at every anchor, because a maximum of one still permits zeros at other anchors. For cap three, $248/802$ of the accepted completions exhibit multiplicity. The fraction describes the announced finite enumeration. It is neither a probability under the original host distribution nor a sample of actual arithmetic witnesses.

At cap four the unique mask in this panel with a maximum of three is $\{2,6,10,27,28,30\}$. At anchor 27 the offsets are 1, 3, and 7, reaching 28, 30, and 2. It saturates the count bound again. Since the anchor is near the end of the residue interval, the final odd target is across the modular boundary. The example distinguishes two features of the proof. The cyclic window must wrap around the orbit. The tile construction applies only to complete contiguous pieces of the coordinate range. Ignoring either feature would misread the admissible partners.

\subsection{Why the fixture is a logical test}

The two control masks fit the same local certificate even if the selected set and its competitor are evaluated under the same step and unit weight. In an incidence objective that charges only failed orbits, the contribution is zero for both. If there is just one coset in a finite fixture, matching retained cardinality also makes their normalized share identical. Consequently these observations cannot determine the odd degree in that fixture. The negative result is an implication counterexample at the level of the local clauses and the readout of the incidence objective. A complete \id{Host} would have extra conditions on the field, prime choices, ranks, and asymptotic scale coupling. No synthetic record used here is asserted to satisfy them.

The fair comparison keeps the active tile rather than deleting it in the negative arm. It also makes the two offsets fall below the actual eight-step cap being tested. Requiring only that some tile exist anywhere on the orbit could make the result uninformative at shorter horizons. Finally, the cap change is treated as an ablation. The original source sets its numerical cap through $\mathcal S$. This experiment shows what the predicate can imply when the cap is strengthened, without claiming that the parameter change propagates through the selector's existence theorem.

\section{What the comparison gives the architecture paper}

\subsection{A scale-dependent necessary budget}

The fixed numerical contradiction in Section~\ref{sec:rank} uses only a small part of the rank restrictions. An additional written derivation shows how rapidly a prospective original discriminant must grow with its scale. This helps explain why varying a fixed integer, while a reasonable initial test at the level of prime congruences, cannot find an eventual old witness with a fixed $d$.

Use the first three comparisons from \id{RankVal}:
\[
 R.n\ge x^{2+8\varepsilon},\qquad
 m\le 2\sqrt{x\delta_1},\qquad
 k\ge x^{9/2+10\varepsilon}.
\]
Because $0<\delta_1=x^{-\varepsilon/8}\le1$, the middle inequality gives $m\le2x^{1/2}$. The rank fields give $r=k+m-1\ge k$ and $k-1\ge k/2$. Then the width cost satisfies the sharper scale bound
\begin{align*}
 W&\ge 10r\,\frac{f^2(k-1)^3 L^2}{m}\\
  &\ge 10r\,\frac{4x^{4+16\varepsilon}(k/2)^3 x^{2\varepsilon}}
                       {2x^{1/2}}
  =\frac52 r k^3 x^{7/2+18\varepsilon}.
\end{align*}
The term $8fr+4-r$ in \eqref{eq:rank-identity} is positive, so $t\ge \operatorname{TOv} W\ge x^{377}W$. Multiplying by $4r$ and using $r^2k^3\ge k^5$ yields
\begin{equation}
              d > 4tr
              \ \ge 10r^2 k^3 x^{761/2+18\varepsilon}
              \ \ge 10 x^{403+68\varepsilon}.
              \label{eq:asymptotic-budget}
\end{equation}
The terms \(761/2=377+7/2\) and \(403=761/2+5\cdot9/2\) track contributions from scale amplification, width, and the rank powers. This derivation uses real inequalities on the concrete source formula. The local kernel report establishes the fixed lower bound of Section~\ref{sec:rank}; it does not certify \eqref{eq:asymptotic-budget} as a new Lean declaration.

This sharpens the interpretation of the old theorem's quantifier order. The theorem can provide a host at each eligible large scale, with its primes varying in a family. It gives no promise that the same discriminant can support all those scales. The necessary inequality \eqref{eq:asymptotic-budget} already rules out any such fixed discriminant for arbitrary large $x$. It also shows that adjusting the number of retained places cannot compensate for a too-small prime. The width cost is imposed at the level of the full rank data upstream of selection. The rank constraint and spreading of the retained points therefore should appear on different lines of an architecture diagram.

For an admissible search at larger parameters, \eqref{eq:asymptotic-budget} is an inexpensive filter. For example, even discarding the positive $\varepsilon$ correction, $x=2$ requires $d>10\cdot 2^{403}$. This is a necessary condition from ranks alone. It carries no claim that the relevant splitting conditions, pair windows, primitive spacings, and existence estimates can all be met at that numerical value. It is nonetheless a strong reason to distinguish experiments on literal small parameters from symbolic checks of implications valid for all sufficiently large scales.

\subsection{An exact invariant of the old odd-step relation}

The occupancy bound uses a short window. On a full orbit there is a separate elementary description. Let the orbit have even length $n$. Write $E$ and $O$ for its retained even and odd positions. For a retained anchor of even parity, the reduced odd offsets visit each odd position exactly once. This follows because $e\mapsto (i+e)\bmod n$ is a permutation of all residues and rounds an odd integer to odd modulo an even $n$. Hence every even anchor has $|O|$ possible full-period partners. Every odd anchor has $|E|$. Summing over both classes gives
\[
          \#\{\text{directed retained odd edges in a full period}\}
                         =2|E||O|.
\]
This is an exact invariant of the mask and the step orientation over the full horizon. Changing the representative rotates the coordinate labels and can swap $E$ and $O$, while it preserves their unordered pair and the total count. Neither $|T|$ alone nor an old one-coset share determines the product. In the 16-cycle fixture $T$ has $E=3$, $O=1$, giving six directed edges. The control $Z$ has $E=4$, $O=0$, giving zero. The assertion can be checked directly against the full-period exponents listed above.

Two corollaries are helpful for reading claims about uniqueness. A nonempty parity-pure mask has no odd partners. If both parities occur, the full-period at-most-one statement at every retained anchor holds precisely when each parity has exactly one member. A short horizon can give at-most-one even for a mask with many edges in a full period. The full-period calculation is therefore useful for understanding potential multiplicity, but it gives no direction-free guarantee about the compact short horizon. The executable finite experiment names its horizon explicitly for this reason.

For a fixed even orbit and a shorter horizon, replacing the point count by the full-period count can overstate the contribution substantially. Let $a_q$ count selected anchors whose $q$-th successor is retained. The short total is
\[
   \sum_{1\le q<\ell,\ q\ {\rm odd}} a_q .
\]
Each $a_q$ is a cyclic autocorrelation of the retention indicator. The original short cap bounds degrees using unweighted window counts. It does not fix these individual correlations. Our previous masked-orbit work investigated adjacency at a first step. The current comparison shows where adjacency belongs in the more general short-ray problem: it is the first term of the odd correlation sum, while several farther offsets can contribute from the same retained source. A failure to meet a retained neighbor at the first step does not imply no retained partner at a longer horizon.

\subsection{From geometry to a gain requires a budget}

The old, direct quantity and the upper branch have different formulas even when they use the same target map. Retention controls whether an endpoint is available as a top place. The lift, capacities, weights, and exceptional conditions control how much an available endpoint contributes. For example an odd offset $e$ enters \id{gU} through a minimum involving the term $\id{al}(t)/p^e$. Even a certified second endpoint might make no numerical change when its capacity cancels the proposed charge. Conversely a zero second term in a gain calculation tells us nothing about whether a second retained endpoint exists.

The decay of proposed odd-step weights quantifies one distinction. If $p\ge2$ and an exponent $e_0$ is the first considered odd exponent, the sum of bare weights over all later odd exponents is
\begin{equation}
 \sum_{j\ge1}p^{-(e_0+2j)}
              =\frac{p^{-e_0}}{p^2-1}.
 \label{eq:tail}
\end{equation}
For the first exponent 1 and $p=13$, the tail is exactly $1/(13\cdot168)$. It is about $0.595\%$ of the first bare term. When the first retained partner is instead at exponent 3, this same ratio compares exponents 5 and later to the exponent 3 bare term. The individual exponent 7 in the 16-cycle fixture has bare weight $13^{-4}$ times that of exponent 3. These estimates concern the proposed weights only. A small cap on the first actual contribution can invalidate an inference about the ratio of the realized contributions. This is why the combinatorial verdict does not revise our previous percentages for recovery in \(g_0\) without a separate accounting of the canonical capacities.

The small-prime packing in the compact proof does more than identify an isolated edge. It enforces simultaneous constraints such as \(p^{i_t}y_t+\sum_{\sigma^\ast=t}y_\sigma\le n_t\) at every target. The original gains and exceptional branches protect a different interface. Removing duplicate or long edges from an old orbit would generally change the mass imbalance, the saturation of capacities, and its comparison with the fallback. The right follow-on mathematical question is which weighted bound can be shared under weaker assumptions, once the literal partner structures have been separated. Equating a Boolean existence flag, a degree, and a divisibility reward would erase that distinction.

\subsection{A diagram of the two uses of selection}

We can now place the spacing laws in the program's larger comparison. In the original pipeline a scale-tied host first satisfies arithmetic and rank premises. This is where the fixed test fails. On an admissible host, the selector certifies density, shares, and weighted bad mass for one set. The good representatives feed several local tests and gain branches. Spacing is expressed through marked intervals and occupancy caps. In the compact pipeline the data already contain the finite set and its three-term spacing law. A small-prime projected order estimate makes possible a local uniqueness decision. The proof then uses the resulting choices in its packing capacity and local loss inequalities.

\begin{center}
\small
\begin{tabularx}{\linewidth}{@{}Y Y@{}}
\toprule
Original path & Compact path\\
\midrule
Ranks and full host \(\longrightarrow\) selector input
 & Data fields and a particular \(p\le D.k\) \(\longrightarrow\) candidate test\\
Density, shares, weighted bad incidence
 & Three-term spacing and exclusion of small dyadic order\\
Good orbit, short bound, marked-tile exclusions
 & At most one short odd exponent prior to packing\\
Several odd terms and alternative gain bounds
 & Optional exponent and simultaneous packing capacities\\
\bottomrule
\end{tabularx}
\end{center}

Both paths are consistent with a proof that must assemble local information into global height and prime-sum bounds. The matching endpoints do not require every local witness to agree point for point. The modern proof sometimes encodes stronger local structure in the hypotheses of a data record, while the original pays for more complicated behavior with checks of local mass and fallback constructions. The present calculation is one example with a proof interface attached. It should not be generalized to all changes between editions without checking the declarations that carry the corresponding invariants.

\section{Falsification and evidence ledger}

\subsection{What would change the conclusion?}

For the fixed numerical parameters, a purported old witness would have to supply a concrete $R$ and $H$ satisfying the literal original rank equation and the old inequality on $d$. The checks in Section~\ref{sec:rank} apply to those fields. A proposed witness can only pass if it changes an assumption, such as using a branch with a different rank formula, choosing a different $d$, or taking a record that is not the selector input. The variant has to be named and checked separately. The earlier no-injection theorem would then remain relevant in any setting where its conditional host assumptions still hold. Spare capacity from a larger $d$ would remove one necessary obstruction.

For an admissible general original host, a proof of universal short-ray at-most-one would need stronger information than the original local predicate. The finite fixture establishes this necessity by exhibiting both competing offsets in the same checked short window. One possible source of extra information is a constraint relating the host's prime labels to the retained set in a way that forbids the one-odd-to-two-even pattern. Another is a theorem that the relevant queries can be taken entirely inside a particular marked tile. Those statements are different from density and window occupancy, and require independent verification. At present we have no witness of a collision in an actually instantiated, globally admissible original host. The bare local counterexample does not decide that stronger claim.

There is also a positive direction for future testing on an actual set. An explicit certificate can list a prime, an anchor, the canonical representative, the reduced exponent range, and the membership of each resulting target. Such a certificate is local and can be checked once a valid original host is supplied. It establishes an answer for that prime and anchor without replacing the set or reading a score as a proxy. Its output has three possible statuses: no candidate, a unique candidate, or multiple candidates. Showing the first two at some anchors would be compatible with the kernel-checked upper bound.

\subsection{Separate decisions for source and finite arithmetic}

The old-section kernel replays import literal original definitions and use typed hypotheses on actual host and scale records. They check a universal implication from the rank fields, the exact inert-class partner map, the same-set occupancy implication, and marked-tile exclusion. A successful Lean run with these types verifies the stated implications. The existence of a supplied premise remains a mathematical question, especially when the theorem asks for a good orbit in a globally selected set.

The local finite certificate imports the source predicate for a cycle and its three clauses. It uses the original Lean environment and ordinary decision procedures on small bounded sets, together with explicit integer reasoning for the tile. The standalone Python code checks the same fixture and independently enumerates the 32-cycle panel. The resulting count table is an exact computation according to a stated finite test. It is not an asymptotic statement, a statistical estimate, or a theorem about the distribution of old Frobenius classes.

On the compact side the source trace identifies the local proof term, the \id{three\_term} projection and the numerical guard. The standalone Lean proof checks the algebra for a supplied projection. The additional import-level check calls the compact project's dyadic order estimate from its native module. It verifies the contradiction from the stated spacing and scale hypotheses. It does not instantiate a full value of \id{Data} or replay all of \id{DepthPairs}. We began an import-closure rebuild for that larger module and switched to the smaller dependency cut once the scope became clear. The proof report records which module was checked. The earlier full compact verification by Sutherland is independent and remains separate from this local work.

\subsection{Cases where the experiment could give false confidence}

There are several ways the right-looking result could come from the wrong experiment. A test that uses an empty or all-even mask will always have no odd pairs. A horizon that only includes exponent 1 makes at-most-one automatic. A disconnected probe that reads the cardinality or normalized shares but cannot see the actual coordinates cannot decide multiplicity. A probe with a one-point orbit trivially removes all positive reduced exponents. These four situations are legitimate controls when labelled as such, and inadequate substitutes for the active marked-tile case.

Another failure mode is changing sign conventions. The original signal prime uses step $2h$ for representatives and retains an odd sign bit in its partner. The old selector's weighted test concerns inert primes in the declaration under review. The proof that the two old maps use the same unsigned step relies on this class restriction. Extending it to the signal branch would be false in general. The compact expression also has \(-p^e\) where one might mechanically read \((-p)^e\). They agree at odd $e$, which is why the local proof first invokes parity. Mixing sign and exponent conventions could turn the two-exponent control into a different claim without changing its appearance in a table.

Finally, a whole proof build, a kernel replay of newly compiled modules against prepared imports, and an executable finite scan answer different questions. The experiments can check missing algebra and produce exact examples while preserving the certification boundary of the two full upstream projects. They do not establish RH through numerical testing. The deliverable includes the pins, local source hashes, explicit theorem statements, and separate numerical results so that a reader can repeat only the part they need.

\section{The occupancy bound is sharp as a local theorem}
\label{sec:sharp}

The examples above show that two or three short odd partners are possible. A general construction shows that the loss between the old cap and compact uniqueness is as large as the cap allows. This is a statement about the literal finite interface, with parameters declared. It explains the limit of any proof that tries to derive uniqueness from those three local predicates alone.

\subsection{A family with an active tile and an exactly saturated short cap}

Fix integers \(r\ge1\) and \(q\ge1\). The integer \(q\) in this section is a prescribed degree, independent of both host primes. Use cell width \(u=2\), tile length \(A=(6r+1)u=12r+2\), and short length \(L=2q\). Choose a power of two \(n\ge A+3L\). This allows the entire example to sit on a finite dyadic orbit. Take the marked coordinates
\[
                S=\{\,2+4j:0\le j<3r\,\}\subset[0,A).
\]
Each lies at the beginning of an even-indexed cell, and no odd-indexed cell is occupied. Define the retained set
\begin{align*}
 a&=A+L,\\
 U&=\{a\}\cup\{a+1,a+3,\ldots,a+2q-1\},\\
 T&=S\cup U.
\end{align*}
The anchor \(a\) is even. All \(q\) successors in \(U\setminus\{a\}\) have positive odd offsets less than \(L\). The points are distinct modulo \(n\) because \(a+L-1<n\). The entire cluster fits within \(W_n(a,L)\), whose occupancy is \(q+1\).

We check the short cap at every start. The final point in \(S\) is \(A-4\), and the first point in \(U\) is \(A+L\), so their forward gap is \(L+4\). The final point in \(U\) is \(A+2L-1\). The gap from there forward through the cyclic cut to the first point 2 of \(S\) is at least
\[
         n-(A+2L-1)+2\ge L+3 .
\]
No arc of length \(L\) reaches both clusters across either gap. Within \(S\) the distance between consecutive points is four, so any such arc sees at most \(\lceil L/4\rceil=\lceil q/2\rceil\) selected points. Within \(U\) it sees at most \(q+1\) because that is the cluster size. Thus all short windows pass with \(K=q+1\).

There are \(m_0=\lfloor n/A\rfloor\) complete tiles. The marked tile at 0 is unchanged by the second cluster. Choose \(M_0=m_0-1\), the permitted look-ahead needed to allow each starting tile to reach tile 0. Its successor is interpreted modulo \(m_0\), so there is a suitable distance in \([0,m_0-1]\) for every start. This uses a positive complete-tile count. For the long clause use length \(n\) and cap \(3r+q+1\), exactly the cardinality of \(T\). All clauses hold with the stated parameters. The anchor has degree \(q\), attaining \(K-1\).

The 16-cycle counterexample remains a stricter small control in one respect. Its look-ahead is zero, and its collision uses wrap-around with only two leftover coordinates. The general family instead uses padding between separated clusters so the occupancy proof is uniform. The two constructions test different potential explanations. A local failure of uniqueness is possible without extended look-ahead, and arbitrarily large local degrees are possible with a positively enforced marked clause.

\subsection{What can be verified by hand and by computation}

The proof uses four comparisons that are easy to recompute at any $r,q$. A valid numerical certificate lists \(A,L,n,a\). One verifies the tile cell positions and the two gaps, counts the long mask, and lists the odd increments less than \(L\). For example with \(r=2,q=4\), use \(A=26,L=8,n=64,a=34\). The marked tile has positions \(2,6,10,14,18,22\). The cluster has positions \(34,35,37,39,41\). Its shortest distance back to the marked tile across the wrap is \(64-41+2=25\). The gap in the forward direction from the final marked point is \(34-22=12\). Both exceed the short horizon. The long cap is 11, the short cap is 5, and the four offsets from 34 are \(1,3,5,7\).

The same construction at \(r=1,q=1\) gives a case with an exact unique partner at the distinguished anchor while respecting an active marked tile and all caps. The finite enumeration at cap 2 independently confirms that no degree can exceed one if the anchor is included in the checked window. This positive arm is useful to distinguish conditional at-most-one in the old proof from a global claim that it never obtains a unique partner. There are old admissible local patterns of every cardinality of the candidate set from zero up to their design cap.

The independent program verifies 21 points of the family, from three tile sizes and seven prescribed degrees. It checks every old clause across the full cycle and compares the computed offsets to the closed form. The family proof and its tested instances have different roles. Exhausting a finite table validates particular inputs and guards against indexing mistakes. The symbolic argument gives all \(r,q\) with the stated formula. The new Lean theorem bounds the degree for the actual original types. We use these three pieces together without using the table as evidence for a full host.

\subsection{Why stronger old source assumptions remain a separate question}

The family freely chooses a length, cap, and density through a finite construction. In a complete \id{Host} witness those quantities are coupled to the scales, prime labels, and independent conditions of the original arithmetic. For example the cap is \id{G.KG}, and \(r=\id{G.rG}=K+1\) in the native predicate used by the host. The window length depends on \id{B.Ltwo}. In the finite construction one can take \(r=K+1\), and the derivation still works because it covers every positive \(r\). The relationship of all parameters to \(\varepsilon,x\) and the weighted selector is an additional requirement. We did not construct an instance satisfying that entire system.

The important logical conclusion is narrower and stronger than an analogy. Any proof of the compact uniqueness property on an old actual set must use at least one hypothesis absent from the three local clauses. The sharp family shows the local occupancy bound cannot be improved as a function of \(K\) alone. The marked-tile assertion changes the possible geometry within its domain. Its scope does not constrain arbitrary short windows outside a marked tile. Hence an improvement valid for the literal original selector would have to exploit where the tested windows are located, the distribution of prime steps, or another constraint on the selected set. Retuning a bound in the proof of probabilistic existence would be a different project, since it might also affect how often a valid good orbit occurs.

\section{Choice of partner, incoming conflicts and weight}
\label{sec:incoming}

\subsection{Outgoing uniqueness alone is weaker than a packing}

The compact proof's local use of \id{N} needs to be read alongside the definition of \id{IsFrobPacking}. A condition that every anchor has at most one outgoing candidate still permits two different anchors to send to the same target. The finite control \(n=16\), \(T=\{2,4,5\}\), and horizon \(\ell=4\) demonstrates this. The allowed odd exponents are 1 and 3. Source 2 reaches 5 at offset 3. Source 4 reaches 5 at offset 1. Source 5 has no outgoing odd partner. Every source has degree at most one, and target 5 has incoming degree two. The example illustrates a logical distinction and does not assert the compact spacing field on this mask.

If a target has two incoming candidates with different exponents, one could try to rebase them to the same source. Their displacement from one another is governed by the difference of the exponents. When both are odd the difference is even. The outgoing candidate predicate tests odd exponents. It therefore does not automatically rule out this new displacement. The compact source considers shorter competing arrivals explicitly and tests the relevant triples at a chosen source. The \id{IsFrobPacking} constraints also budget the aggregate amount received. Those constraints are necessary even when one can rule out certain collisions combinatorially.

This distinction is directly relevant to comparisons with old gain terms. A candidate set is a graph determined by the mask, oriented step, and exponent horizon. Picking one candidate per source is a choice function. A packing is a feasible allocation of weighted amounts constrained at both sources and targets. Old local inequalities can rely on bounds for the graph or on estimates for all weights incident at a vertex. The existence of a choice function with a desired property alone does not establish such a weighted theorem.

\subsection{When incidence objectives can and cannot see the graph}

Let \(\Gamma_p(T)\) denote the directed odd-offset graph for a chosen prime and horizon. The source's bad-mass functional instead uses indicators of failed orbits. For fixed labels and fixed prime, if \(b_p(t)\) is the representative of the orbit containing \(t\), the contribution is
\[
   \omega_p\,\sum_{t\in T}
                \mathbf1_{\neg\mathrm{Good}(T,p,b_p(t))},
                \qquad \omega_p=\log p/p .
\]
Goodness marks an entire orbit at once. The function only knows how many retained points lie on each failed orbit. It has no separate coordinate for the degree of a good anchor. If every considered orbit is good, the value is zero even for a graph containing many edges. Replacing this statistic by the number of new partner conflicts would require a probabilistic expectation bound for conflict events in the selector, not just another way to interpret its existing inequality.

The two masks in the 16-cycle example illustrate exactly this loss of information. Either local graph can have the same all-good status, the same retained count, and the same one-coset share. The graph with a single odd site has six full-period directed edges. The parity-pure graph has none. On the short horizon one of the old occupancy bounds is saturated and the other mask has room, yet both still satisfy the inequality. Knowing the actual occupancy at every window would be more informative than knowing it lies below the cap. Even that would not by itself certify an oriented edge without specifying the individual bit pattern or sufficient correlations. The original score only requires the cap.

This is why the experiment's two outputs should enter the architecture paper separately. The rank calculation closes the attempt to retrieve a fixed original retained set. The local geometry gives a positive theorem on an arbitrary retained set with a good orbit and a precise nonimplication for the cap-only route. The compact source supplies a distinct algebraic mechanism for excluding a pair of candidates. Putting these together identifies a usable comparison at the boundary of the two proofs without treating any one proxy statistic as the whole witness.

\section{How much of the compact short-ray cutoff is used?}
\label{sec:newbudget}

The compact proof uses a short tower cutoff \(e<s\) for its candidate. The field \id{three\_term} is available up to \(324s\). Those are distinct constants. One can isolate the role of the arithmetic guard by deriving a numerical budget directly from the magnitude fields of a compact datum. This gives a stronger geometry-only implication without changing the source's packing theorem.

\subsection{An explicit dyadic calculation}

To avoid overloading the exponent of a candidate, write \(m=2^v\) with \(v\ge1\). The compact \id{m\_pow} gives \(s=22v\), so \(2^s=m^{22}\). The magnitude fields give \(m\ge\sqrt x\) and \(r\le3x^5\). The definition \(k=r+1-m\) uses natural subtraction. Since \(m\ge2\) and \(r\ge1\), even when that subtraction truncates it gives \(k+1\le r\). Therefore in the small-prime branch
\[
             p+1\le k+1\le r\le3x^5\le3m^{10}.
\]
Using the entire source candidate horizon gives a uniform estimate
\[
     (p+1)s \le 66v\,m^{10}, \qquad
     \frac{2^s}{66v\,m^{10}}=\frac{2^{12v}}{66v}.
\]
For \(v=1\) the ratio is \(4096/66>62\). It increases with \(v\), since the ratio between successive values is \(4096v/(v+1)>1\). Thus the arithmetic guard is comfortably strict at the source horizon even in this very coarse estimate. This explains why the automated numerical step in the source can succeed while using only the supplied scales. No estimate for cancellation of prime sums appears here.

The same calculation yields a conditional extension of the uniqueness mechanism. Suppose \(v\ge2\) and the chosen prime has the needed unit label, as in the source comparison. For every positive exponent \(i\le324s\),
\[
        (p+1)i\le324\cdot66\,v\,m^{10}<m^{22}=2^s.
\]
The strict comparison starts at \(v=2\), where \(2^{24}=16{,}777{,}216>42{,}768\), and is preserved as \(v\) grows by the same ratio test. If two such shifts of a selected anchor gave selected targets with \(0<i<j\le324s\), the field \id{three\_term} would be applicable and the guard would exclude its conclusion. Thus at this level of abstraction the selected triple test and the size inequalities together can forbid two retained targets through the whole \(324s\) range, when the label is valid and the exponent convention is chosen consistently. At \(v=1\) the same estimate suffices through \(62s\) because \(62\cdot66=4092<4096\).

\subsection{Why the cutoff still matters to the valuation argument}

The extended range above is a written consequence of the spacing and magnitude conditions, using the native order estimate already replayed. It is separate from a Lean instantiation with a full \id{Data} witness. It does not say that the compact proof could change the horizon of \id{Q} freely. Its packing uses the horizon again in the valuation inequalities, where splitting off exponents beyond the horizon controls the total proposed cost and interacts with the lower bound on local depth. A proof for a much larger horizon would have to preserve these estimates, not just uniqueness.

This identifies two independent parts of the design. The finite set's algebraic geometry makes certain collisions impossible. The divisibility argument chooses a horizon at which a specific capacity model is useful. The geometry can have slack even when the analytic or valuation choice has little. On the original side, the cap is a bound on counts over one window. To increase the horizon without changing the count bound, a proof would need another occupancy control over that larger window. The old and compact conditions therefore have different monotonicity mechanisms.

The inequality also shows why a comparison of the numbers of positions in the old and compact groups is insufficient. The old host at a fixed discriminant is bounded by a product constraint on ranks. The compact deferred selection can reserve powers of two by scaling the tower with \(m\). The relevant compact reserve is \(2^s/(p+1)\), compared to candidate exponents, while the old local statement puts a finite budget directly on points in a window. Even an exact dictionary between residue names would not furnish this reserve or its associated three-term implication on the old set.

\subsection{An instructive limiting case}

One might try to weaken the compact hypotheses to a small window count of three. The 16-cycle mask makes three the first cap at which this idea can fail. The dyadic modular check makes the failure visible in the same form as the compact proof. There is a common mixed power and a nonidentity relative displacement. A stronger order bound would only confirm its nonidentity and would still leave the two points in the set. The improvement has to come from the selector excluding the common-power triple, or from changing the count cap to two. Those changes act at different levels. The latter excludes any three occupied positions in the window, even three which could not produce two odd partners. The former can leave such a window populated while excluding a specific relation. This distinction is the central reason to compare declarations and counterexamples together.

\section{Reproduction and proof audit}

\subsection{Declaration map}

This table lists the native declarations that anchor the logical handoff. The bundled JSON stores the Git blob identifier, exact source digest, line and permalink for each of the audited files. The original archive matches commit \id{40e56bab}. The compact source matches release \id{72fead10}.

\begin{longtable}{@{}p{.33\linewidth}p{.60\linewidth}@{}}
\toprule
Declaration & What the comparison needs\\
\midrule
\endfirsthead
\toprule
Declaration & What the comparison needs\\
\midrule
\endhead
\id{Ranks} and \id{Host.qd\_big}
 & Native rank lower bounds, discriminant product inequality, the finite retained field.\\
\id{RankVal}, \id{TOv}, \id{nWv}
 & Scale-derived rank growth and the literal identity used before selection.\\
\id{topsChoiceH} and \id{topsBase}
 & Required inputs and existential shape of the old output. The latter uses a representative for each orbit.\\
\id{hPH}, \id{baseOf}, \id{BadOK}
 & The Frobenius step is chosen from the host, and the score is computed on the retained field of that host.\\
\id{Arith.Loc.partner} and \id{stepP}
 & The old signed target map and the separate signal rule. They specialize to the same unsigned step on inert primes.\\
\id{topsPos}, \id{Good}, \id{Cand}, \id{VV}
 & The reduced coordinate set and its use in choosing the local consumer's representatives.\\
\id{CondA}, \id{CondTiles}, \id{CondIntervals}
 & The three original local constraints used in the finite model.\\
\id{D.three\_term}
 & The compact spacing premise on selected triples at bounded exponents.\\
\id{two\_pow\_le\_of\_pow\_eq\_one}
 & A native compact arithmetic bound used to exclude the mixed-power consequence.\\
\id{Data.exists\_isFrobPacking}
 & The source location of the short-ray uniqueness implication and the separate large-prime construction.\\
\bottomrule
\end{longtable}

\subsection{Why checking the exact input changes the fixed example}

A weaker argument could already diagnose the late fixed case using the old index ceiling. In the previous analysis a host with $d=9239$ has $R.n\le16$. If the first component of \id{RankVal} also holds, then $x^2\le x^{2+8\varepsilon}\le R.n\le16$, hence $x\le4$. The archived proof of \id{topsBase} takes a threshold at least 16 by replacing its previous threshold by the maximum with 16. Thus the weak implication would already exclude using this proof directly at its advertised large scales. It would leave a reader asking whether a very small scale is a viable exception, whether the numerical threshold is merely an artifact of one proof, or whether the old host could be prepared through a separate arithmetic argument. The width equality settles this part of the question for every scale. Its contradiction does not use a late-tail threshold, the upper bound on the frame size, or a result about where arithmetic Frobenius primes can be found.

This comparison is a useful instance of witness verification. Checking \id{Host} alone exposes some structural necessities. It does not show that the record can be used as a valid input at the point where one wants it. A complete typed call to \id{topsChoiceH} requires \id{RankVal} as a separate premise. When a prospective witness is artificial, auditing that boundary can prove impossibility much sooner than expanding more of the downstream consumer. This is more informative than reporting that the available small primes did not produce a host.

\subsection{Audit modules and their exact roles}

\id{OriginalSpacing.lean}, under Lean 4.33.1, imports the original module for \id{ThinScalesIffHostBound}. The rank theorem uses natural and real inequalities with explicit casts and the natural ceiling. The prime-class theorem rewrites the Boolean in the partner map using the inert constructor. The cap theorem is first proved for an original base record. Its host specialization uses the canonical representative \id{bsH} and obtains an exponent for the given anchor by unfolding the representative definition. Reducing that exponent modulo the orbit order leaves the original anchor unchanged. The theorem then uses the same retained field \id{H.TT} in the conclusion. It does not bypass the proof obligation that this original orbit is good.

\id{OrbitControls.lean} imports the original interval-tile module. Its cell theorem works for arbitrary positive cell width. Its two finite declarations use the exact original local predicate. They make no existence assertion for a complete old field or for a globally selected host. The ordinary finite decisions in this file are checked by Lean's kernel. The decimal and rational values appearing in the paper can be recovered from the integer results, without tolerance or floating-point comparison.

\id{SpacingLogic.lean} is checked under Lean 4.35.0. It separates the general mixed-power identity and contradiction from the source's implementation of the place group. It also checks the dyadic modular certificate and an odd-power order lemma. \id{CompactOrder.lean} imports the native compact module containing the arithmetic guard and replays the contradiction with an explicit projection premise. The targeted compact import closure is built under the release's Lean options. The order module itself is separately replayed. The two native project environments stay distinct. Names such as $p$, $s$, $h$ and $T$ in the prose denote the comparisons already specified, rather than an equality of objects across toolchains.

\subsection{Running the package}

The archive includes the single editable LuaLaTeX source, PDF, four Lean modules, source trace, results, receipts, and a replay runner. Use the bundled standard-library program to recompute the exact finite controls. Compare its output to \id{results.json}. The SHA-256 manifest verifies byte equality after extracting the archive. The TeX embeds the original Nous PNG. Compile its single source in a writable directory with LuaLaTeX three times to settle links and the table of contents.

To reproduce the old Lean work, point \id{replay.py} to the pinned original source directory, the Lean 4.33.1 executable directory's parent, the matching Mathlib checkout, and prepared native objects. For the compact check point to \id{grh-compact/src} in the pinned release and the Lean 4.35.0 and Mathlib pins in the receipts. The runner can build the small compact order closure if a prepared object directory is absent. It verifies the referenced source hashes and toolchain versions before elaborating the audit declarations and invokes \id{leanchecker} separately. Its purpose is to replay the specified proof segment. Both full proof projects have broader import graphs.

The receipt records executable output and object hashes. Path-sensitive build information can cause a recompiled object hash to differ on another machine or directory. The source and toolchain pins and accepted verification status are the portable requirements. When the old project imports are supplied from an existing build, the runner checks the anchored source files and assumes that the prepared imports were built from the pinned original. The replay's scope is consequently distinct from a full independent build of the archive.

\section{Sources}

\begin{enumerate}[leftmargin=*]
\item Alp\"oge and collaborators with Claude. Original Lean formalization, pin \id{40e56bab}. \old{Development/AutomorphicArithmetic/ThickFrames/topsChoiceH.lean}{30}{Selector and hypotheses}. \old{Development/AbelianVarietyTheory/Definitions/GlobalTwisting/PairRowIffQdRows.lean}{17}{Width and rank identity}. \old{Formalism/LocalRingCardinality/DataHost.lean}{39}{Host structure}.
\item Original local orbit source. \old{Formalism/SubspaceCounting/GoodOrbitIntervalTileCondition.lean}{9}{Tiles and windows}. \old{Development/AbelianVarietyTheory/Definitions/GlobalTwisting/ThinScalesIffHostBound.lean}{19}{Host Frobenius and weighted bad mass}. \old{Formalism/LocalToGlobal/GlobalTheory/IsResidueClassZero.lean}{77}{Partner map and local consumer}.
\item Sernau and collaborators with Claude. Compact release, pin \id{72fead10}. \newsrc{New/Data/AdmissibleData.lean}{35}{Selected set and three-term condition}. \newsrc{New/Curve/DepthPairs.lean}{59}{Short-prime branch}. \newsrc{Lib/General/NumberTheory/QuadFieldSplitting.lean}{83}{Dyadic order inequality}.
\item Andrew V. Sutherland. \href{https://math.mit.edu/~drew/shortgrh/}{Verification files for the shorter Lean proof of RH and Dirichlet GRH}. Independent build and checker evidence on the pinned compact release.
\item Nous Independent Research. \href{https://github.com/nousentllc/riemann-hypothesis/tree/43141b319692e076cd6f6132534a3bec9f86fa39/research/uriel/experiments/selected-place-transport-expanded}{Selected Places and Signed Indices, expanded edition}. Earlier conditional capacity result and the fixed readable datum.
\end{enumerate}

\end{document}
BEGIN FIGURE nouslogo
89504e470d0a1a0a0000000d4948445200000296000004960806000001e725dca6000000097048597300002e2300002e230178a53f76000020004944
4154789cecbd3f8c16c913f737fdd3459665c0720e64ce00a796059b387002a44e0ed23739909c1fc889250707927320b0e4c0d241e0c801acf48696
60035b7a238ed8b2616559b21cb8ad5e6a96d9e93fd37faabbab67be1f6975dcccf3cc9f7e7abe53dd555da5b4d6532a4aa9575aeb67c95f148c52ea
add6fa49c9156635e6e59795325f3ed35adfb5760e8069c0699a7ed75a2b8eab2d6acccb83fc6cd40bb82eac1635aff55fd6960ccc45cd17662e76f1
f7b647832d3192b4bca66975bdace7e2e899d64117bfbe8313adf5277b33cb798ddc7cb67610b59f9a2a8d79e504e186f5f1729aa675833f98a6e94f
cfe7bdb4949daa8da994faa1b5bebef8ffbabf1cb16e4073de168d5aad31636f80a1816f6badffb1b63ace53bb41595e40252c5e064fe77fbb5e10ae
7df4efcd86243e585b98a9f602cae905ebeff9cc18aee373d3bd67ee89e2c6544a55317346a4a831e931bcef7a89d436d895520fac8d69dfbf62c8b3
5c13c3d8fc0297d6a5e893eb3b358fef3b76096cc3496b6382d993db3b0a8f7fdbda524876632aa51e591b09572ff291d343388e9f60524593fd98fb
2ed2f79908ee69adbff8beef3a47caf137be7faab52ed2e089a3315d17e9fa5c08df31b87eb0c077cd8f7727f49914aadb997491f7ac1d3f79577a13
f3e8c9daf193e7a1e3734f6a57ef99b9b87a5c953948c6fbc8ea990b1bef86b5f3c0e43ee61fa79fbfe60f6b0f03ae5e3905b697b0b03b8b0719b9de
c9e68ff88ce447bde40574666d6120d49053c309e61cb27b668f5e3953ab7776eb99dc9304335b3754f3312f25ab31d7ae5deb03e5bcf61ce1b1b5a5
80458738e1f891588310387b8deb47aa70fcf3a5c3aff8989c6e0b6e2dad314d56d312611d4e9a0b3406bd52eabdb55300f38f5dede5d9dbd51b82ab
677244b8459da74144c7f5dc911247632aa56ed598bb7451bd314b1829ba6e82ab9717342623684c46d0988ca0311941633252ad319552ff9bb571e7
541d014de52317d899ae199f23c0de98eb863c52c3e205c4086b63fa7aa16ffbde606b4ca5d40b6be3c1607b9bd7f02a1ef26d2efd316eb586937d55
6f08f4cc0d527ae5003dd838b1aeaffe5fc7ae28295d6dd16509b4ef47312e0a6be3f677d6ab2ebe98359fabcf3b57945898c73cf7cf7c3df52fe55c
cb63bbb6fb3eef3b966bbbeb5c8bbff7abff7fe2faeee5315c1b536f34f1efbdeb78aebf8cc6fc14d8676d731dd3f5b9d0bee51fcb028154625f26be
17906ffb149860d9723dc71cd3b56f496ee470ef17c966cc912b10a276845d51b06b2ea53d73b5efdc9720c0d59b7de70eedf71d734d72cf1468de5c
5bfdff37eb138d486acca50d5602c74a605f0fd15a7bcda3da243de69cbdd2d718bef3b93eef7b34e947ff2ef63137f981ac8d9df1dd586a6cd3b2b1
1c93db3f56ffef4ddb16dd33b9b5d2d710be737245d4397ab1f3be283cd258040f5dfbac2fc436a652ca4491ddb47614e0bb20df8dd66a4c4e621bb3
f94280f5795b34a6ef3e43df59b2a999be13ec8df50fe77a791537e6c11af27cb5c03fa941838d7934fff7d6ca8b2d8b26a899351b53e2dbdc778e62
3b93dee020016f63729b423d692557cec63c980ffc7cfe476874138353335bfc9221ed715d075334dd07adb5954228770e608db36736e0a4c7495d0d
49bc9bffe1189bcfffef4bdaf2ebb3eb9e29a557aeafa56051d545faf0ad639067f3abb563e37b578e7180c64c3ec6fcced05a27bd3bae3466abb79e
c4c6a4746c7ffb3e13b316f4376bcb8e588e587c0d11ea40a17d2e2e7ba6e45e99f23ddf3142df5fce5bba3e27fd6d9ecae69b7483d3f0ee2938e68e
6588c65c672f8c61d52b8b3310c670d198d21ef151f997c4a163e98fdb2baef35fa9b6542eb9375512201b48055905719a59d22bd741125aebd8f851
af6b38e57afe35b57914b29225675cd7f79cef2e5f701b3ef4e08bac899dd9c2b62cd549dffd93ff3ceabafeb5f58152121af2caa3d6b2217ddf736d
a35a454e7e5b7fd9f70bd5842a4cada3d936e11825c57cdfb77d8df502e2eca19113c04fd6a5ba52278e63bf531b6b0ace77b1a9e434484ca338bef3
8d338cb0e447f2ce1a153ef2c1acafbe8958df85fbaec3f7f95c5ce7498951f2f6ccad1b09e13b31931edfa891387a716d57ca88a534a6a5996be820
efac1d01d4d5da93dca562bed3f1aa24ca2bc91b971a39fc23e7ad5b1b96acac577f6ca717738ba4e1a489c5a9997fd2d870da2e3cb7e9c964eef986
8739c7ac9e3dc647a11ffcfa72e8b886c35f9473cc1a094fd68b3897dce0e8d9e605143a4ee0fc4142c78ca1468e0e4b53178f2dfb5bd8d700f4985a
abd47c78825dcf16c7db9caae45abcef7cec5c37591b57af8c1d40b83eb738de66c1e6e2c6745dfcd47978e7baa6d891952f7eb3ba66fa2ebaf73899
ce7fa5bc8eeb5a5d786ce3283b9b7589b44487d9ea3a838faa6b989b3415c8b5aa57aae771ade735af336755af15102ad985bbb6205c1d818be49e29
711e318616d79dba447ac8869c1a5d6bcaaade611b72c165cc518dc73d7b55efa8a12e35a33da27ae65e1a92b8f4e173f7cee455bd7b08beaa753f87
5fd5cb9577644af501ed2924b0c67d61893423c1c63cda12e952828db964cf51bf5cba79e4c77ce9a8b326b67308be80f6faf299e1be3fbc8018f136
a6526ab95cc41b9308166de67bccf7fe884f78ccf9a0005bde631eb567d698bcc10be8275b6b2ba30835e6615e3a5c6b2bbd8ff974f55138d35ab36b
4c4f6ac8d8a18d766e428f3948048dc9c856635ea6e9e69c91de2b41cd9ca09b49e0316724a6312fed4dccbc8739a4abb7167ba86d2186a8c6446f8c
23a567423b3728291472be953dfa681405bbe2f1bf4a8e9df97afe071ef7abb02c10400ffd49d608c8b7f0e8e8640f275d0d5ac34935125596fb1df5
b12f9ee870351cf5d22e757c7bc2592cde9b72c2d5e07b843d1382ebb13f4aa3b2cf670656f53eb7b6ec8c6a393a6652d66b8f4ef5c63c12d98fb971
b085f47104c8ea60cb94cd61672e6b52de2ec958559bda6bcf394d239361e0fe729b88348cf6d3536deab08a663a6ee0929a0d1caadc2a6a89740a0b
f3e8c3fa6bae847babbf4f2669deac65e6bf26fb026d777d7e999c64dd90ff6fc05463c79b3f93891fcb1b09f5d805f767b9502aab0dd6291b6fb5d2
f1aab9e0227209593a9bc8694c6ca519eab670b154d3cc9cccffae69bdf9dfeb1ece717c6e101ec3488d2c85de244c1b9cc6686a61b986aa838ca2c7
dcccacaf6bf5e48ec55ddf733de6aecfe51e9f9bd2c6b42ed0b5adf4585385c6acd1b8256373dfaf708ff6b394d272f0928e9f33a60ea63c2f8525b1
9eef2d9c50d7c2f979df397c9ff7e1fabcefd82570bc80bcc5e26864e2cde64a239ba2174aa874b8f143b59cd92aee99be5f35f5265cc709f51ecee3
bbf6e550cdce4cb9c09c9ba97dfc1cb2c6e68b5ef1cddae9b8095f2f0adda4ef3baeeffb3e1b3a7e0d8a628d2a4fa7358967e2bc1791c349574f736d
9306c6e68cb0056e7111ea81a17d055813d8b9e4440e5f84c1d4684c5ff2fc2595cecb323d971d865deba6ac8d0ed8bd8a4cf794fc986f992325c4dc
4cad86643916f7ac11cb45796e506a8f9c297a9bd345444dea661c77735b2eb33793db73c91984a0b91dfca1b179c9316b59226caedea5960aad7161
66973ed5bc3676bff9a2511f69ada38b1bd5a4d50f5c73d6a87b439a1f746a38e1213a3e73b4a585189b3382c664048dc9081a9311342623684c46aa
35e611974d5769cca3ae3faffa981fad516bc4671e76c95b8d2ad29bdbf60adee68cb03566a80786f6ed0996c644f6989f347bcce7b9c53d53ad58bc
8bdc38f4a9703ed30437ac0b20d7a0b458fc2889f582512213c5e02f7f3cb34c505d2d18bf8de999b97fe6eb897f8f52ceb53cf66afb5bdfe753b6bb
eec1b58dfe5eb9bebffc63adbc1ff9e3c557b2f73ce63e0759aebbd9732fdfd6ab845de75c527b556f357c0deacbad14f8bc856f55c616599a29dd14
2ab93e47839f581ff2d07c046496ab581bf38f25ea47ed910dfb2f6b0b138e5ed59421c7e6318da6947a616dac4c52630e326c9c97d3fcb9dea194aa
1a6512dd98d28683bedea9b5be656dfcc5436b0b23293df36f6b4b67020daa626bc0ade5c0b1d6f3e3625ff8c97459f28e51c227c78820fbcf750ed7
9f6f0494f3e73b8ee71e9e05f659c798ff920b8530716f9d41c1856f0494438b55759b8fb9a3db7320723ace9578ca4c78581ff410a399cee159216c
990139301d66d97357bdf76bec93196ccc8aa650f4afcd4184cd79d161968de858dfbe2d4b21cdacd898b199b2b82687bdc709ed4bbd066fcfaca495
33b9b98f4a61a9c9ebc3db336b8e76521395a47cc745e84d5ebd671664cdda23d1bdd9d9336b8fc10fd33339131a4b20255955a983d0ea992d66865a
f6cc9863f83eb32c31e1fbee12b13ea0964f8869a8b9415d9d29f687bcf298bb0ed4918f9ca7de6a9079e433cf38adfe3f0a4b336b93f9b8464da7ad
49e91c6a5176476bbd9e967b1f95db73d6cc56bd52925ebaee39e4e6751d63c921e2337d8db0f51887f6b9b868ccdabe9105cdcad8b87a5df5735ecc
100b7ec453bee73bc6d6f7439f0ded5b23fe31afdd909cfc4b5aaf2c8533622415454ea2eab47ac4737a65e83ba17d6bfed5a8c7dcb6b644d02bd238
97269a195ba8a334fff0e27fcfad0f941f73938bc6acf92be61e3b372896be9b3afb73a510be2ffc9ab63fb1be4d2c7be6636b6f435c171f83a32173
f2175b0eb79c6c5c57a6e05aa4217391db20ad52e9c6724533992f266aba3fa7411c8fe1a984543dd67ce6726eaf8454572ef1cefa50f8f3ecbdb1a4
b73bdfe6a51718337c7334cc3badb525ee731500d74dd6c85ce8da6e7dd8872b9a6bb9dec6150516fa731d6723e28cedf8dc51721413151fb5e7da68
7dc87153bdff5cd759d0902f5c8db9b8f717aeefadff9c8ff91aee478a83853dc85194fe324ec0232796e9e4c4d5c211bfe4930a3dd378026f39ce15
1b68fbc575ada98f7989ac58ae5ea65ef360cb21c61588d5e2b8d1762f7763fa2e28e5a238ce57e0377ae48adf8fb27fb91a73e9b05f51add0f0eafc
d68d248ea46eacd7a4a7ce4271a695b01a92c4bbc99a74bad9d78eeb8ac5b5263d2904b14a2684de6f7ed71bd9fa90e3f32513c35369634a6cc89994
0675dd47ec779770964a7c27c9160df532d7673dd77ec3b3dd4951a52adf854ba2e5b5162f91969ea5dad143abbd104b57f57a6b4e4a62d5a0aeb736
0b29ab7aadf9c998257b1271740a16a23573b4c7db456d577054cfdc43431297c94b6af4ce1cd3283a9b8a34b4d65597e4c4acea5df7cad1d7085dfa
99b87be7a666f60e39a941ad7b0af6cc9895ac20b231a769bab3f8f7536befb82c1f7536d90a3ee67b7cc4676adc1b123833e26dcc75f889f50160e1
7dccf7fc884f78cce583c664048dc9081a9311342623684c460ed998b5b2be1ed2cea48ad29739d8b9eeefb0467b0d428ff9376b0b08e2ed99137a67
32a19e0912416332126ccc50261560136c4c90464c635eaedd46ef0c137c9bcfe0ad1e47ec638ede194156e0167aa79be81710deecdb70c6b41f9ea4
c6748434a3411724f74c34a89facc71c0deaa66404746571001ab4a0316971001a7401cbaa5e57231ed116e55c22ed3cd0911a9575f1beaf41a78334
2aeb145c6041e721a8329fb968d4d7d6ce1d5325e1490892820f5a6b9145e98e08a533fa5d92aa34ef98574efeb3bacc1f8b4d4d72a31c155a89b45c
40f36da31a6f37ba76cc35eb08a01551f5008f0ed582f096771ec5d612e5e7354fefaa08cf7249dc7d65d77d5dfefda8992f420a46f51cf7becc38bd
ec94af97ed39d400c095694ff21f6376c4b7ae8c882dffa669320fd22bc7b565ddcf68bf65e84fd4ab3c1552c8ef946f7a6dafee952b03475a90fc4a
6bddaa9c531386ec98f4ca3ad35adfb5767ad8b05fa5713bb612c8a24d4c5b7cdecb7cdd50b184264395e994642f4577cac961bffafe165f39d9cabe
1f8159537eb29e8bf49d7bf197d429e9985fe64a05f4100ecd308a392b5e0b45f02525757d668def3b31c7e4a2e5b96a319262367f0d732594514a75
992bec5993af942114b375786dccf9521453e2f54ba79b623ae6e0926c4650d4f63f7c458fa46095fa6a81a7c4e0e7e96755d620f3e027f499126aff
58bdaf7fbdbff6f5e4d2e555ee6abcd82a09b19fe7b8aed0395cf710fa8e94ebf75c77d2d45b0b8619fcb822e73d8d9c8ce7582f190eddf3fa9de5a5
5d9d7515d82182e61dd3d180bec6b2707dcee12b8eb986b7a1ef24559d4b20e2fa7fc41c6ded2ff75c7f682ef3b9b54518bde23197984083e4690dc7
7138b8175b52c2777e57e773e1fb7e2156a9321faef34bb2359b2aa6a7b648d65cdbc29b525a86fddbc2e312db298be73717d75f6a32bc5b5c7f54a7
9c0698466aaa982dcb8bd0e4b20976b86f8cfb699ade73bc9e0305202f60cb22f733d5df03ba7ee3dafcc46d5e385453ccca82dd76cc1a387e481fc9
411852ee49ca6fd2ec557ea04e69f85a2bb969054a4da12a3453ccd50ffb546bcd51cfbd0905039521968348148d261d938218beceff3f925a728c9e
47b85f69fef556aff2c3764acee3b4825603f4bd86468ab93ec9302b1e1dd79ecc208a7965b6a1f73537514cc72abde58a47d1318325ab0ba5af4c34
4b7d17dea38b4e29e59a9bbb24579dd42c5ff86bd149452ebe9fd215e44c6a8734915daba5beef1cc2d11d69090f8608700dbdde4586900d38f89496
f0e0ca93eb0b52e8cd2803b845fb5d1f2de981e8b0b755277d7ff4149c312ca39452630024316cc283de91d7be87a4c735498d422f61d81a35ab9205
62074db558ce35ee3129e9d029627a821cca7541552f2012744c2012744c2012744c2012744c2012744c2092e13a26721d1d83613aa6c36ffe992a34
801d325c1a42075d6a03f9ae0913ec3c88574c5f0758702d36b50a1807b11dd3a4b64e8826bac695fd17c840e4abbc20bced71abb222be6bc4ab9c07
711dd3f783c7d2aa63f8ae131d930731af72534ccaf763271e07e1523b4044c724fbf0bbb523ff78e89c83d3bd635227ba6fed283feeb0a54472d9d3
03d9b563566ec8bfac2d1df114446027359adf513dc4fa5b7cf6966bffea8f657d51d7722ad6c601cf91c01bfa01a39c01ab6404b13c5f2ed38dc4e4
6aff10f8e8c9e23a628e7d8765ac20a56a452d6a8d927df7e03b5f6a46b5dc0c6cf4bdac04acbe7bdaba06dff77c9f8fa175aaeb5bbe9ba8784e9176
57c4753d4dfcfc9287d6960222d6a39f5a5b0a6999b8f56dc66b66d7843a9b2b7f68e8f38ecf72a5ac76966559c19ea4b64965b4deaa2561ddf55cb2
d9da9148c2bdfc1d53692ee2babba4ecaeae989853bcc2ebf58650fb1cd98b54b563861abd3512922250e998736b470057e74c28b435440e5217d53a
a640a51461df52ece83b6b470057e78ce463ddbba947958e89d77718adf593d4ce469f3fb376ec14f68e894e194f6a5a405321973b9560a826e57aff
325fd2b29ea54f99536b7c5ef92ee704bbf44ec9fd835a1b050c5846afa734c3a698be1f4a12adfcd57b62991a7bf567cdb372c2a29803bdbed9aa65
f8ee7974c5a48a6e7f5a3b22e0bcf762c5f4fd40602c16199b979df2de3245f62a5db63510e34c4d5ee4f941a79445eeefe1fade96fa998198efbb1c
9eb66cc5745d10188f5a6bac4a8f9bd53107ee94a354c4e5a2b4487f3792073f232ba5efe9cec1d70ebd063f942ee74ab85bccb570dc87c901304dd3
cd9263ac49524cdf4d0011b0c66026c21e8114ad987be8943b574ceb7a12a3de2d4abedf64f0e33af1802445f58c84ebf74974753a3feb3a6ecc677c
c74b61b363ba4e3c223d32c28dc4d6dc24fd3da0bf2f9e39cba75c6f8ce03ca6e3c4e057820611787ea3e7d69608e6b9c9e9d75298df57df72056b3c
752d0329c56b63fa465a83c29a6ccbd3192e686963d2ba9ebf7b5e432d9c1d93821dde583b0685fb87f275ccd61dc2751d89d33c3fe602fa3e7cc773
9d3bf73a5cf86cccdd74cac8557ed15090437738cc2cb2bb4fac1d71df553527f02dc5dc9b5dd94a2d6b9ccb87ef1a4ace4feb832c1b3277929e5531
5d2718990a9dd25b2903793979b9ec98be1b1e987b152efdb3b5a521bedf688fcb7c2f3ae61e53f669ad59b28ecdf83ac5d4a863f8cebfd7b5e7b362
8a4ad9574a85577828ddca0d6b0b33474cab3d77ccac9199442afd58d65c21616a0c552be56226f28f9aebfda2636aad3f8d1cbb375341298379e16b
ba39039996efd5e894d40792985d93dcd732ada78b288dca9019d92a744a2bbeb1e6f916e7ed66cb86ce9d43d1f495c7f3636f144ccbb9caa9520709
9db3b19bd37b1d6be6eb8af94eea3d383d3f03d92f1f98632c9f6d34f209f3f91e78a2742ee0ceba11836335e44bfa3b71ecbb729da1bfd4eb702ae6
950f781aad37cc1de4eed61c25d7f928077bb0740c8a58791473498fa7760bc64ef2831ebc50a72c56c9f93c742e67a77429d191895ef323a4d16e30
74924f8b4e128aac39a37b4e1aadba52aaf8ce23b9333a52c26cfdf1c6a89a5779ce9ff96acbbf82ebfcaf12aff381eb38ab63be777c6ff3cf752c69
7fd334bd75dcc7dd9836e0bc17f1b98b1216441917e41d6b8700467a3d17967d39c9990f75115c5a11cbf2e2b93be91e824b02f770a3a6e7281557ac
e956ba172a7a603aa3990337a1732c1db37a01aabd45c357e6bcf7a2b9c043b4a47a5df8aef5ca7dc1a90d314b7a6fd5522d06f382edd59842c65b8a
fd81125748df47c92b3df42aea41a6eb37ab0c5f09f48a76f9eb7df0e51f95dc31733ae3a8f38089ab52bbd9a65ba619db1cb3b48e99da19f73a211d
db0eb5d2de1424e47aa9b5b60651a9742fa43f2dc2a73c37bae6f111bc24b1f7c899c57775dc9857b2957687a3534e02063fb127af3e0a1c8198198e
c29592aedfe39bd6fad6ea734e7f3fab7a4bae570ebfb19f8d36b43a13d3717db0dbbc4d3b66e44ddfe35e48b667426d3af283ddc4c65c44f17859d8
54e8940950e77326d1aa657fb6a085e7277802bcaef908b5f568ed5cb5fa6ea8a14c055a744a5eb6d4d3da28982a8a0995ec4fe037e8e2e64ca575f5
dd2a4b4f814da09d3f8e5053934d315bb9aa401a01a1103dfbc115288c4e29185fe794fcbb14774c9f1760069d5206a3754e0e1b139d720c9c99957d
1db637458a19ba29744a79f87e2f89bf5549f55d6f50053aa54c7cbf8bafc3f6245b310337735b6bcd5e5b10f0e0fbdd7c9db617b965a1bd13b4e894
e239755d60e837ed4196621e3599e85e18e1f71311c10ec09ae48ee95a140f0037c9af72bcc6c707af72003241c73c18a1f96749e0557e307cbfdfb4
8357f999b505ec01a72fbd17391dd3b9107ea37a1810c0865a8a728c2477ccc0fa615ff530208455768f6f8bab72ae13ea093c3f402459a3725f0754
4abdb5360290017b7491afd3029042f63ca6af03fa3a2c004d3a66088cd041291c8bd19c07f0292a0031142ba6af03fa3a2c004d3ae684ce092ac066
63a273024e58073fa1ce296d4d09904df36c6fbece0bc0922ad345d4f9ac8a06d32ff5bc6bed0060d94f7a6614867a021f4d8a03843ae7840e0a1c88
aa5a810e0a669aaef9a18e77c3da418c5c6501f0d27c319a0934a60e7acfda49cc1d149df4b8742f72ba95f875095ef5c7a1fbf2dd594163abbd2efe
5e591f00bb4164bdf214155dd1bcd83ca883e842fa33a8617e3c86e8986b421d159d721f0c992266b90c75619f3eb53e0886a599621aff382aebcac6
d8f681bc01cd68a696a6432aa59ed188fabaf501d08d39245142879c7ad9962bdb1045043ab0fc0d765110a094d55ce5d7c57c24a6772a61cca7b517
2d76ceb8355d073b8e46f97bd17043e46d948cc984b2e8849f17977a7bd7752339094df98c5253bb274aa95be6cd13b884214c25a91e9e988b3abc2d
9ae009bb2165101383c8b94ac76bddc5d7957f7cf7153316b317f32b39d8211773b9c374c849fa04faa273beb676dafcb9eaa443077298876c7d3fd3
34fd657dd0e654ea002696e1dc8d4cf1972621e90bad75d7948834db6014fe8eb533913db95987f481cf340a10363515e701d6e5406b1e7429a59629
bdcdbf8d9d675677deb78ec4cfbd3d7ac986ef94b3421c2882fdcaa065d9067b61d8ba3deb1f6315a0f1d2fac2b87c5805a05c19b4986d7b73dd0ed9
29b7d4416bfd62d549c525b20ff07ad50937bd5c748fdff7e2111b72a013ea9009c779324dd31b6b475b5e9a0788eb8c64c20c3527e96228a5e4b41b
cdc8db1197f978faa53c2f7d85e32339a563dc5e4dd19c2fcec93db7fa2e7359892886514acaf676bfb651bf0c56b076fefa8ce94c7f5a3b02df8b39
2e07adce53939194b2c514cbb42aaa043a3044a75cb810abd79fd45adf9a984d05a5d4eca36f312b707b1a7c8a6cb4958d4d268bb75e81a9afefd6af
d4d15fe1430d74b0c6e718a0e87d009a36028de9da294d7479ca0f6f42b7ac8d7561ed94adbd2ea33e54dd3a25d93d0fcd047682511e13ba55c42ac0
82fb78d5839257e17ab7ac0f5ccdd3e4dcdf9b2e9d52f8c8f0e3e2dfc5cb2f564b38ae591fe0e78fd0f5afda3eb474a21b626ccad864fe2d5f49153c
2e4d5fe1eb354dae72331205424ca74c185957f35737fa81aab90123aedfea941219364ac8da587ecc660aac94620f9830cb69ad8d2b6a287f0d9a77
4ad72b840209bcb8268195526c7396f44a5d2b305bd22cc7f55f8be944b190e9f3fbeae3dedcf3d2e9a194960f5b6b9da35277381493a699ac576a83
f53bbf732826256df8bcde1e085fb37cfbd2a68e7eb3b608658eb05e5f5d497ca5eb7884b780402e9eebbf56e3fa43c733be7dc7f7cc5ba2eb22ba25
126c4a67793d0fce8086d45c445b55292aba339d01258bebdf9c8130a6c6c6f5a7b4a7489a0764ac1b33f454bba0575eb5f9bec82201c9f1948bef56
6df0d8220739df6b4553a5e4485aa5b5be5e4b0d5afc3035cf91706cd183a0d6afef87d6960ca86372a6943e4fec2c450bb4e85cc1198744ce52aedf
350892b4e8ac77fdc6e224558e632691a35ca173669823de63c5c0394892f20aef9d9fb2384061b1282be9959e9b6f671145ce42ee7530e40b125bbc
6037f194e695bef8a1cc92800f8bdda794df5231fc9837ad2d0cacaeedde6a2525e7f55fd03b8f528896d521deaebc0ecfb5d643654573bdf21c0cb3
eeba7426a4162d3ba5c8064821d6fe1be5dee665cb8b4d22b225a35346923820192615b6c4dfa557a7fc362f651d85bdd68a94d8297b952cd975879c
2243c9809b269df2a0f5bad7a1642291a8e84d5edf23db938593dbe7e47d12cdea1ebb67076ede29773ec0b118e17ea58906921104600a221e21a7cf
89b5a523504a0f11d5bba2194d2d77af94aba24b5628be5438ab99a194741a2d5edf775a8c6400002000494441544b97dbcdc152d471bcd65e8f96ee
b977e68c56af6fd7abd0fc58ad7303255318e92e7ef44dab42af1497eafdfa6e1d4f692d231871741a8bf47b73dc97084f5bd3d1f75c4a6419fb488b
a058934a71931bef686d14827166b8a681a478daba66f275358cf521413894c589f00e29becd7b479e5f09581da072d6e67a70a91dd2bc8d561df286
d46b153179be88169fa87296c8113ab9df9c6bb72543ed39a7387cea2aa727097189f8978321c1aae36d3449d74c6f9d3925cd30e182e2dc8c8bc1d0
377a9d4b4cbeefccd421099acabae8909206313188cd25b4ac67b3c7f2c13541c992cacc8321ea9c43e457ec855973b3281a3aec433c5ad6b5eb504d
37d419cd2046f49c6f0cc374cae957ba11458af91ec59e2e07860ff6f4a00ed52967e634c9265bc568eb7d38d9eb5b63e820dfa376c8398fe55ecd18
449e0fc8decd16744a200e744a200e744a200e744a200e744a200e744a200e744a200e744a200e744a200e744a200e744a200e744a208ee13a259245
ed9f1195d22a180ff6c5509d725e10155ae20ac667984ea9947ab6fa7f2c22db292329e55fabff77168107e33344a7f4bdaeb92bca021988ef941baf
e92a1565415f4650cae06b9ad293801d21ba53fa5edb2b72533f03a188ed944aa927d6460fb02df78564a57c636df103db724788ec949955638f5894
749788eb9405beed3fac2d6048242a65b66f1bc11afb4054a76498de41b0c60e90a69498de01723a2557e48fd01ce92001119d5229f5c8da98cf1dc6
63810e4851cabfad2d20893d4d8975ef943502760f1a04bc9b29312c1cdb0f1ff672275d3b259635f0a1b57ec431c833b5d9e7da45bedf4729f57ef9
19fa634bf5dd2d117fed72caa327a93701295aebb7d68e3049833c5fa7f3b1f1f9af4a29963cec3d95f2a3b565a7382ac8c6f0462995da29ab117bfd
1c6fbf2e9df28081b9738995d41fecf7d4d762ca289c54edb1b5e3eaf1ae67286a5128612fa56ce2b911344d727ffe47c60ff6d5da12266914aeb57e
3f1768b576fec4e5ba3da53f1f45a184cd3b65e3c18dc46992e41fac519b6d0555cf75c2cddf03faab62b30f59716c404e576a39d4202c74ad733157
6b47014d95f2c053409fac2d8944b6ddc52b957980147a4d57a159a74c5973c37cdeee23d8b996e4924a0fe87bfaefefd69e4c7a54c56da994296b6e
3861fb817ab3d591b5d6bbf07f37e994586de86643c54fac2d07a1955262b5a11baf8a6bad8bedd051a9de29e1dffe4968049b428267c5b26347e110
51424aa95155e7b9b5259e60ba1bc954ed948254f2beb5a50f96d724a4687b19b8a48278ca8668ad5d7eec2d453bb7b6ec9c6a9d12b6a497a44ea6b5
3edc5af62a9d522975d7da082ec8ec64afad2d3ba696527eb6b6802537525a436bfdccdab863d83b65c870ef09f332de22b4d63f5c839e102347d1a7
524329a54e4588529bc5a0e7a5b5d343ad8e49f6bf732500adbf71c62d98edbeb143688dcf16aca16b1b6eb3de489916ba24a793a57cc72cc3087986
123a8d599af16679eed4e51129d7cdad945eb719e8826b0a4a3c6c4a89a00b91043b65893950d3c6e5544a045d0016583a25ca86c8a0c65a7a93e0c0
91788035f9c01a2ea5445e49193847c9a9cca36a1aa4f8121c7c2d19618728b629e14e1445f14033b30882f9ce3dad354b6e500464ec9ba4c1a7a743
3e5e2cad558179d5cf5c5382459dd27313400ed19dd2f55b52477cbfdaf62230f2fe9d23ee014ab9634213e74b7c1dd2fa60dcfee2b887ec4ee9ba11
00a6817309012178c4e5b6b5c541402ddbe712f2dc08e80867612bad7557ef1c94723fbcef7427ceb42e2523f1e44e09ef8d5892a3a038522506d2ba
64cf99e62825bc37fbc1d7a1ba92d4291109b43b44bef552951291403f11a9301938e7317d91e6ad88ee94c2a3ca0f8d2762c7e70ebcc495a290e0c8
90979da02b4529478f2a8f5e0b332096e2053a5c134a127419dfa6b5d1fad04f39ef955f9285c0446f32a179da1eab0e5dd7137b1daeefa67cdf778c
927688ed94c34f96a353fa717d3ff618aeef96b6c1e6ebdb63af8003d06b1c116353a6d67139324eef464d38923f0494cd84a279d7cbd750c9096e46
767ae4c174257f0856117311e84c7fd1b2871754ceef856f1944e01849046d4ad78907e54c6bcd96742bd02e27add342bbaea56890e1385e04e79cd9
e10ed129b9071fbe76197190e3c3755c1735eed9bb700ce169f2a9f91bad52b43c5878b13ed57e1b7895724f9d72af4ae9b98ec7eb7535a3e11ce8c0
a5e847524a4117a51dd2142df00d6442cc35cde9af68a9ad532977f6ea66558e50dbb4544adf75704c9a6f1d8b626a83218cbeefc66029256758bd04
467f9525f2a1d179821db214ab537a8a8e836d9ae525f779d9b4d649a605a959f275d3f74e6a55aeb05edf3b7b75b3ce9f4d81f619edd51d73cc48df
b7f1f8fcc5792d5794d275612353a1434aae7af1d4dad2801a05a8ae28e50e3b6593a9a01ae7f2e1bb86d2f3bb8e5b12fec6a294d2a73ac071b8544a
576f1f991aca156823b63478217ce7afe556ecae94204c2884ab45873c12179d5229f5ffedec9edf595bcab146982d71a9d1b4d3a24fbb544aad75cb
25a267d61666f6e6d0d8e237df13087ea194f27a8538e33403381d1a7b2d8db73ba5acf4433db4b6342290bba9ba42f7e25f472a4499c3c6abb3aa6b
91ceedf4333752e84d6a64d3b89812dad12bfc0655886523d436b51f68dfb91b4e776d4e7579bec73225d4c545c50d7787ec89efc7ae1504e181dd85
18c39e26cf6f7367a00db5494d950cc52b562eaf6c113a9fef3ba5d7b91ce80cad96153a6497cc6a14f5dfb4438670ad4230db421db218a394f39ff9
df41ffae2fef83e32fd40edce75ab4bf193458e79b2adde3eadc8f1ce78cf97bebfa8ceb1cd1d7626d709c40fa9febc60a7fa0e03dbbbe53b943be75
7da7f57d3bfe5ed5682fd712db73dfeb432215c2d39c51dd339546be6f03a916cf1a7aa86ef826ead7d43425ac4e6902630feee509e54e62cf15141a
d450e47cb3f9489abd501bbfff6920f9bef5d99c5436d67288cb1de10b1341cb20dea9fdf9d89772a44231b6170f45cb24aca14e692ee86f6b871c58
6318373a086b87244f4de835f9ae71508928bc9d72a2aaf88122e45d6995049560eb2411e7629f6f1d8d60a79c84bec61b774896f3c5bc791087f013
6ba0b3c63494b08e79c3da92498b0e19f1aa9e24d88f92880a5d13f4049f73f8b74d7c64ed0e39e7d6d9ea9054e85d6487344b4016f9815c7f96b787
e5bc5bafefd54586a62faa3342dec584b7cacbde654542a4be1d39852b29c8979ee81aeb5f62ce5dbcaeb9468734afe7a582581fb07947eab89b0e99
fb1d1f494ab9b880183b898d96e992b7ce45811a1fad1ddb344f3d9d83abbd5c6d5233e779d6720863d735b43393cba9cd512cae868bf8aecb765a2a
614a87fc40aaa806e990968de8fb9d7ddb597039c4139df85f5c4e78a6bf4fae733aaee199e0a09117ae6b96f8e7b8f66030c53ab2c8f5999cbfacd7
b70bac8a8c43f25ca4eb378c306774ec6763615bcd481774dbda01aec09582b915110bc3ee7107aab029e59583fe0cff0a45db800552d4d3a594334d
f36fd6e894574e10b85160c1be1a3385c8dfaabaf7a97a3282c5e853355e8957c2294de1a8d01fb93c5f32ded7f79e0f71a41a5e5b9820552299aa2b
65888d88eb1654f3aa7084fe49a960b6057b9c69cf4e990a838a7473ed153c80afb5d6de348435c87427b3991e43744a5370689aa6fbd68e38c455e0
524a9978c99bd68e0023a8a6b829a15a50c3a476c86f0b5b505c1d1dadf5add42934b2e19a2ae6622c1045283b5d0a629532f31532847fd9458a2af5
9c426ab16c4464a7ccb01dbb4ea5709132bf5b73115b64fd1ce76fc4715da25edf19e940eed12b661789adccda1cfa51373d24bde77f6baab5984e49
8d1c3b3a9da36f769900dfacab8e55ab8dfc99295c3e08ae68a19688e894a9f6546a0dc251a18eb95504f47b854a683de78efb77ca840e79ef88abfd
cc031871df9f193ae69501e2d6efa29472e5ae6459d4d76da0438df8d9dae1004b4f7f12f100670ff80211f5d6317df3ac5cbf53974e493ed337d60e
9b0f477955c752734aa664f0d46de1180764446f76c823d98e29a404dd66f0b2c635a5d2b45392426e1ad1785d87a9d5314d5cc0c28bb3b56af543aa
c7279666afef581b121d329ed659e25ad15229d121998950cc2133b73551ca5649a48e4aa87d476cd7ea4a196ab01974c83242ed37e27294aa9d121d
b229f77c27e30a296b45b5d737cdf8ff61ed58800ec94b4804466aeb9a9d32786074c83a84da7d9436aff2fa0e35cc840e599550db8e52cc9ebd5392
5f3484d7f601d5699629af841a4a6939ea179cef350652121b6a297eee92b55346bcb691d7bb1dcf3d678a0984e90a5ba7ec513e0ef8d15abbe21d2f
90ae96ad5201a6944e034c8416a249160916a5dc7af2d021fb306a9128aed7b7d74ec16bbb3b4f5d17b0f166eb4af1eb3b10467f013a657f7c1d50ea
6fc3a194e890f271a62aecbd94d60787527a0f804e2907dfef24f1372a524adf8d4ee890a000f159d7001bdf5c07f2acdfee4af6eb7b63cd8db55618
08f8b13d6f36696fb5924ee9fd225edd32f1fd66d27eaf1aaf6fe76b028058d83ba5c9526b6d045238735d8734bb32b78aad3397cc8457b778467885
e72aa5b343fa266901482157299d5f824aca67cf4a094035923b2555d202a01ac9af6f9ffc4f787d0f41e0f71353ee85f3f57d626d01232126109bb3
08fd904595c025623a25dbeb1bafee31f0fd7e93a0df10a36f200e744a200e744a200e744a30b3590fb215e8946046ccec093a259811d32931257430
1090014006e89407629484fc787d1f08df6f37e1f50d4018744a30e556afad055edf0761a404aa50ca836012a82e4a218baed091d3299d09de47a9d1
022e3ae89745077d2cad493857339e69adef5a5b0148ed5f58620ba4019b1288039d128823b753deb0b6fc7cada3c41d28863d3f25ec4a500a5edf40
1c259dd2375f89573828a2a864095ee1a006787d03719476cadbd6968082021043b58a637885835c385edfcef5c2bece0ac0162c45e87d1d106a0972
a83ad0f175560042b0744a2822e0a4fa9410d412a4c299c9d7ab964a29315962817c58063a33e462bc63edd8e8b4005ce9479c9d72da785da3638218
6ad8944e2fcff4ab46380041d89572825a8242aa8cbe431d2fd46101982a4f0979d713a3638210555edf331b9def546b8da922605175f23cf41a9fa6
e9be52ea89b5151c9eaa4a39b3a198f74c1a116b2b382cad22cf9f5a5b7ef1195345604913a59c7eaae58f699aae593b7e01c50417345ba3a3b5deca
ca6614f3adb5151c8e664a39b3615f1ace233a30d831cd3be514394fb93172073ba6cb12db980e8709f6e3d26ddd776cc7a4011238105d931144bea2
af51e7849d7910ba67c8a08e796eedb0f98e57fa311091b68546dba109f64b4835d139774c97d177888c0ef7586b3d44cd411087b84e39958dbc6f9b
7a31d656301422b3ae2d6abca4f2152b27c747742a40ea98efac1d01b4d6622afc833c44bebe5d28a5cc6bf9a663d715e0091a9f6192a66aad6f5187
7b6ded04bb6218a574e11810bdd65a3f737c140cc4d09d7289714722ba681feca6532e2197e47753c4749aa647982602201e8a6bfe7d9aa6134c1cfc
629762b9c6917bebc3344d4fb4d60842028765218a4b6ee0b97073886a3aa644fec25764f2773c9c638f567f5fb0be12ec0933ca524abd77f4754d42
f96e7e36e80f42e9e110966508a5941966dc0f7c64e6a5d6fa85b515000150da993711576282986f4114d339bc58ae89f5b7af78aeb57e656d0580b7
6f3e9aa6e95546ff4478270310cb0d32c5738989c47b81ce0ab620c7a419bdfcb1f1d12d208e15805826a2947ac5d099971867d32b781df70fcd873f
7338554a403ae246402c0b216be04ba1f5b98509817a8ba1be5c68cef049e4fc770908e7e904c4b2024aa95bd3347daa2ca03efe1f636d4cd3f4dfe0
a1ca8756cb3e2001ecf13b1a9e6aad91cb550810cb8678e2daa4714a426fbca55f46145c12ba5bf4f7a081b5570a164f0c00c4b2334a293387f5d7a1
1be1389cd162882f34faf88a39c771384450ba348c938802838de5f66a1514ac5601f4df8ede5e03f2dcf57bd2e2888bba3dc68aa46d0f4c4810d25a
cb0762d990c5ca89ebf3c3133abbc9cd39a7b47108e90da4b9e9c629395aacdf85fe921c71f43b9bdff4d662854db06f80f66018de80656aa496c9e3
68a8f784c25542950ac12fba8772adc2d35045520810cb8aac1278885f2eb9f0008fe014496538d1e9f592056e2096955825da1c360464751f4582af
9432dffdd3da112047241629fa6686cda4b36a7f58991dc19c6505561dfcf5c8b1722bb1fa93ac65b1d05cdf52284f464e1ab16affcff4c2011d8058
32b34eddbe93b4ed278b7fdfa17851719045f979715de73b09cc5f3af2fea4d542a031104b461c352e9e0f7d438443707e27e79134beafae47e23526
e378e1c6a46203cc402c9970bded77be96fbabb5a52326c1edfaec7bced98898ccf6402cf9d8fbdbfe74bd8186bd527828e85a5a8050b0c6402c2bb2
b3c062572811ca687604a5f1db7238b1a49a24563d12fa7b647da18cbd678c11f1b0fa44632f9ee38005efbc6f5087438925cdf3ac9d004bfe763869
4ab8c378ac6e480f177238a06692623a05e3ebb3befbbec438e21c46c1f2cf9a6b076e8e665946cdf3f82c951c98c5b739f430f9445fbce5b683f6f7
8e4e022f89f9bb4f221c716fa8681fd8e0306299d2213263f33e585b7e9d7bc80796a625bc8eab51621847f51c33e43f8db51a5df3d160c5912ccbd8
0e91151ba9b50ece77d2b0c737f7240e1a7aff1db82e51198f3696455e1bed85b5a8ebede3c4b3fd92943c9994571504388458a6cccb14c6463eb6b6
5ce5bb742b6776800586de173802a525f02e740df4c2126d0d9bbe1a21ecdf12ac7aef88670512506f7014cbd23b945c71c3da9280c94b68924d6c7c
e39ac444af26cc891e529f33e1920d2bae1b5aeb2711c992ef53fb5b41ec3d5988e4665f35394ead8d1eb6463c201ec459fee28c63c50765e58919ca
5fab18b2148d116d7a483fc77c477aaa3012122b80dec1c345fb775b16b948081df3423fcf6cffad11cf7c2d70f404d8bd58c67680adace529d0503e
c54afd7bf1e0562d5a355b308b87347625c8d928391569ae6e734e6fc1d7459b54b5384dece7aafd6331e9f1b2e6bc69c413031c3d018e6059c67480
a7d696428c954ae2729678a49b8e58384d757ba22d20aaebf2657d9ccc659937385f262d30737a99e2fed0d1f69a442e5aac8cc3c4bcf81ced9f1cfb
49a52a8ac2b462dba2e728473abb4efe4b3f7cc8a37b410b8b69d0f021d664b3ad92ffba18b4fd599316c78622212bbb9bbd5b96228472fa651de458
9a3db841d7bb9bacdc8bf63fb7760a6351f88cd509480e3090c9d11d3cb161156c50395425f0edfdbcd643baa0fbf24333ef476d7fdbdad9979346fd
e29eb565051c3d6e763b0ca789fa60da2e49824541c12d63ddbea584a094923b0c6ef51b31ac9649e53425689c9398df0243719b3d5b9641a18c0da7
6885f1a03b6a833f8e881b8ce1b5a3b67513a15c04b98bc60c511dedff9469d8fedcd1fedd3206c50821d2bfd9ecd2b2a41ffaa3b56301de9cf57154
59cca1a9057c1456b5c99de019b9ca5ec5327853e804f5a13027aed213e7b93186c00f9e93348ee8e009ae1f06e59065cf59a367b844188310747285
d2c31d91dd59965bc30bbc2deb42494baad523c2efc70be528f0aee2427bff628f96a5572863c226403e14745eb5701b2c4c5eb6a637e0e8f9c5aec4
72ab40d89e82aca54143b6267194104c764299b2828ed223b1ab6178e821c270a21e14c4dc3c09037e533ef0ec6c7314070f9c3a95a08cea5db2d58c
967d5e38de2c59e407383cbbb12c43491af066ac03a593bb29e0524e46a9072499d0ef8967685f62e9bb91db5aebaa39228fc89617b5034fb5d60875
2924f01cb166a01a915d0cc303791ecf2194fc0814ca894aba62b8588e2fcbffe12df75d5896110e86c3bf15b9615acac80986e24cf8accba30fc577
615952528250c283cf8b6cd5700830b0c804dfdb79764689292094059878caad52178111dc21d8ebda703119d28f42e821ab057ebf3212d6ef7f4095
c89d9795988918a623510303b42820aa4a6421ace5168e44c2f409da78c521e22ccd307d91a3d0953c60599616739b999879616a638e1c9c2ede55ce
e4be5b16436c9f50ae736ea28d571cc2b20c118acfc4f0a30cc6a1392cff0c36a21690273491a3d7e031d6d08b5576eca5a3685916b5a814e911a1f6
0cad3b8ef97d1484329e65f9638750de6b9d297f4f1cdeb20c11c8b88e00e844a4d7e0199940bda977a8e8c8c7e12dcb1073a1fe85d5794a1f7fb3b0
3883998ec04fa8fd5ea3397830239d8505390be5f9a28cb18250f202b14c60e5289a130f2c633831b40903a7410126246e2190f33cfbd385385e8763
a61ebfedf5c66a439df2728848739a5f95badc84d00b508c231ccb04e16334d30158964c381c45f3443b04132431970f260bf233258399fb1684b213
10cb4a186fe33c345a2c2543520fe06521902f56318fe83702c030bc01b46e793964379ef45b3d0bed0319cc230f9a6f84e75f30b02c3b60bc94b350
d2900b614807c28c30a80ae62c9288231d00886567287b0f423c76ce7251034dd1e0053918104b001a601c8068e7b1815802004004104b0000880062
09000011402c01002002882500004400b104008008209600001001c412000022805802004004104b000088006209000011402c0100200288650328a9
eb7fbffb1b0560c7402c2b62eaa72c4ac0fe974aa9f32e170200280699d22ba194fabfa669faf75747ff0f8c7822233600e301b1acc0c29a7402c104
603c300c67642e7c1f73c4d8cf01006400cb92891cf1838509c038c0b264a0c44a244f390a5601201c88650129c3ee0dbe433001900d86e19950bde7
6b8c87fcbeac2d0e009005c432835ace19cc610220170cc313504a3da8edc586971c0099402c23514a7d9aa6e963a373413001100686e111f4102f0c
c90190052ccb0d7a5a794aa97fac8d00802e402c3db4989f8ce0a652ea59e76b00e0f04c104b374aa9f7ade62723f80b319800f40773962b843a5710
837910cc8b516bfde3e8ed201158960b247ba1e1213f0c0f8ede00528158d2db7c0431c2fce5fed15abfc78b512687174ba5d4131ae68ec05f835c27
284429f5056d288b438b25059abfb176080656c721783c4dd31da5d4ada33784240eebe0195974cc705c6bfdcada0176010dc5cdad7cade9d83359b3
569be6f9d2fbd687fd9c6badbdd11a3472cb3548ee69adc558d8669988b571efecc13a1b71750f3d9c7f5a3b02a4de27fdb6efb4d64fac9d4c343c87
e183d6fa91f501de7314b1fe8d94526fa769fa9df152cfb4d677adad8d39dc307c2fc358aceeb159584abf53526576315b38d9aa9d83f840ff7d68ed
e1e38ce148afe77f507b6866a19c684aa2fb737b18b134f33f3b9befbb696d016b0bec4d85dfbcc5390c97d32cb5faadb1d68c55b8fc9ba6e9b6f5c1
30af16225995decfef21c492ac81afd68ec181b3c7e28eb5e597c5c315bfe89ccfa373b00d97b5d69f56c76f320cd55a9b11cb536b879fd073f5742d
c685c2dc7544b57bb1248ff76e436e5a3d443be023f585524e03dfff9be91c2e3e3bb6d5225b905682f8d6fac0d5cffe43a2f9d2dae9a7db886ad762
49a51f9c96c08e68f91049e7c3c6f5dd67b04cb6c490e31c4e84bf186fe43a1db5d66bafbc48762b963444e5ac912316589797c48453dd2cb4fe82d6
d2e21c1c82b9b6b844be18c98a2c5dcf9e625d7661976279c0b93c58978e79be00f773978ed29c5e0c46304b63616b0de939e11239f1f7ba3bb184d3
e3f084e6149794cc636f0df767fe2849afe712ff8a73a2598c3284e6603762394a328c5a602df14fb4d6d15eefdcfe921824ce9d7760ef73f062d985
58d29cdd28c9306ae10c9b3928ef626fbb2093d36b6b8b07c7b2c2223047dd87e1c592565060ce0e29dc2e495c8698351cd75aa7b475d212cf089017
a003438b25ad411d2a6b506590c28d480963c9b5d412cfc199d41743f10e0c2b9634d1cdbd0615ec8b1b9177936da92508e6611c217b6548b13c48b0
7916dcf3632343b17f318259da975a9ce30af89ddb339c581e29d83c13eef9b1a13182593b9d5d8b7338a8961e0eb8194a2c114309726921668d0513
59a71a338c584228e3419d713719491b72cff1dcda01866708b18450268310220f66c589c3028c5df5137b8e57748ef3c5e66fd607e5f231f6ca1817
43bcb7b678e895a64dbc584228b3c0bce5068b9c8ae78e84be5ce7b8be38077b3df0d2118489519e13f72effac0f86b9937a0ccfe753fc10375dc7a8
1dac2f5a2c2194a0362468a5197382d0396a5843a5e2502d465929e5b4142b3fd35517a788adee08a10460930785d97abe2d1c45f354c4a7d57f2f30
493da834efba3cef83c57fafcfcb6e03ebe74f16c3fce039a7693251065f1c01fd0f56ff9dc3b2a297b9e620522c219400d4476b9d54979cace3b585
9c24d6944929296ac0917da94be62551c3f0a3670ee2c4f13606e3fd86496206ea22462c69b2fae899833881588e8fc8c073f32276385772fefe1929
cc4dc4309cdea0a12a71201da4f11a1f112fbc8a868c992ffdaed4e5a8fcdc38c3ac4f09a1bb6509a1ac0602d3c7a76bfe03134349d362a942991b53
7a6d619dc6d43a6a4a57b1845002504415478749544322b99550fab9ab1eb8711c39b6ad83f4b7f89d44534ceece6e6209a104c04f6480356b291173
cec800f11312c024215b05e9c7f28714a76f17b1845002b0c9a6107106d3537ed8cda06e12c9228b96e6254fac1d0124086673b18450021045b3f94a
12ca98f345d71dda8204f7f6c6c7d6d7d955309b8a2584120059500dab28614eac3b1473bc7f523334f514cc666249e10710ca768c50a01f3820016b
45d71a56a9f39e53c705174dc41201e70024112360c5eba09552bef5dbad499abf4c49e7c64975b18450760396e5bee188431491f734c361d4a5ac4c
0bcb1242d981528f25e843ec1093e9f745d1bf04aa8a259262801558afbe4d749672065833c41722e95a9c54134b082500d5e01296a472ba8c25244a
a99ab7d24715b184507647fc5b1ad824d40267a9199e3194df5afed804ad75976c4cec6209a114010af08f4954ed24e6f9e8a7d69600159fefd8f9d3
6e9533cdd2256b63f6c120942270542f14414eff907a2fdca42cd8e06e9384153c33df52b3ac87a064197f043e3273aab5ee36efcd66594228012822
76c1466a4ce226244029533737b99e770ac042e52ae200002000494441548f11ca973d8572e24afe6b523a591b412fd8d6ef0a60a45adb4da8151266
8488c296526a861bc17caab5ce8af98c155c29a38b62cb920a9e7709120536dceb773b738858d1042badaae3ce083109d399b5d3cf1bca3b1955ea77
59abdcda69f354d2344c916549d98c6f5a3b00b0fb4ace1c97b86cd9dca4d4a069350cd55a5fe4d2a4e7fb77eb036ed886e6261b51a53aeb45645b96
4aa967090d09dad0cd53184172b8c7415621c5ae704b4998cb8209d1596439af1ddb78b2c8aa2e4e28a75c6f78eadc06688364cf313ce136e4dc88ca
fa23ad2d68a4f03673c9e4bb5eb19225248b25126388a56b58c516104b9b8436610dd50179e48825428404225d5832fa8d68f12f852248a21ca34789
35954e928307422916d121364aa99cfc837bf2ea5f81f248c6469054992ba4a9b4e5cb68fe77f4b03a55c47b9c939368cb124229971d5a95bbb6a652
daa3563b703ccf1962d9fc9c9c4479c36363a840173e486ef694d09805cd3dbfad48148c1bd6163e92d68433f1b86be317b2695926c65a81c60c6055
fe93118b7bb2c7b0a1c4b6e8e231ee61f54ab0b463085a9634c700a1944b0feb2095e4450b3b15caa4051c1d436bd8d79e4750d3826623289688a594
4dee9adc569040a4b2a7b5ed17503c65b4d17134efb7d67a88dc12de61381c3ab219e1818263277d0147effb4fb95ece6b1d21a986d3b284508a47fc
f03b335c685719de3356baddb3b60031586299397402ed38973efc261e5a5b36d853107a86509a0c3b526adc000757c45229f55fc0a1231bad754e28
4e53324726bb99aba45a3a2942f9729017e0a159afe0f99f8ede209219649e326be5cd5ef270669468c84e9e0bda72299698a714cf28f3597f595bb6
b9ddfa226b90f10ced329e74affccbccad4028c5f37284f9accc7e7426357f610a19f77e03423916ffa21f0c75a6e562b2ef882f6d9bfbc29db3728f
8a997648bd774a708bba558371310ca76245b02ee5f16d040f716eee80d1632a339e99f3111c74ad514a0d110570e90d47ce3c798c90f03573edf734
f23ca549b1962194cff72694b43289e31877ac1d02b156f0c0c294c1209eef5ca11cd6037c845549ad125ba42440e6385f2956503a2cccfe0c22943f
3285f2dd884269868a1942f9e100a53192a7604cdfa1b64c2ea14dc1fe5db0c4728260f6e47ca035df39b5e2872b546556b4d1fd260d15c989f3c8da
31062999f76fce75c0d731b64aa9bb26407fb15f17f49d996e6d6a0dc3afeccc1f668174c417a52a2c56f77204affe4c4ae5c5158fb5d639ebe245d1
703aeee2059a71beb3699a5e4dd3f44fab10aca0584ebf5664e4041a8378c45b5c94182379bd37318c801488e4d9e861506b2a0be6bd65ecf0082527
5083a73fe285a4f0b7bf31424c2185afe47865771d0e6486d2d3347db676e4e11d5dec4a2c270ccb6b205a4832d6395f6180921725d30a878b99544a
9969a22f29952963474c819795593063fae1dbde2bbd72ea8697743040481612066b42f490b4f025b0bbe13688c3e90d0f612c217ad0abd4333e0067
528592bc97ba50286f4b1493a5573653289f92871b427950922dcb353981a507e69ec484184cf352e286a5854ea90b104607668ac5f2f24070000591
f8d095ce492e1031f74a736a5fad1de920751ab0481e86fba0210adec2362f25b58b1194c2e1e892d73d33e8d0bdfc58dc4f89503e9ffb3084321d5a
2fff691d804e7f3f68df33f2790c099b65b90696a61c6bb28253ae8b9383ea4371963d11392d221d2a9bf127f365be969e2dbf9a585e9ee098e146c6
4a79656d6d4841707588262259e9daa711975b4aa175148cc4512adb30dc8759c247372ebe7c2b07348c6b2e94cbe13559f59c62f392db13ec5a335c
e1da5f2e86d60a429907fd26b142f97cd5e6f3b39fb2de5c24d52d4b173bf5a0377572303a67425cdc13591546281fd0dfd4e0dca9c06aac40cc745a
aa1518137d21d221da432c2f4fcee7bdec49b5b9164a47f5cab3b2e1a8fc9fd334fdd7bda7398e400da15cb2717c7111095dc57209f31ad41664cddf
9100ce7fd2acb33d60867b4fe0d12e83fae966edf3520b302098e2c4725d37bc1be495bc6cf80186ea77023f34e88771267e54ca7a86cf4944874f9f
d688d835ddef0bf3769e7844f901ad09178318cb32c44e86eb4026c396b8a8494678d0eed7cc0f21962e1a3938c0f1109f84b9158523a7dd59f2d543
876a614ac43a42149e8f7a3f400c3757a14c5dd38275e6ace0f4660aed6f4768d83f34521c8e61c5d285f190ae057421a42f1d5f0136ef6872dd6ac3
1e7f14a3776a5d653bd6e27998213b0dabb9dbdecc297f5db5e910b58a861d86d782de7a771771857b19ea1b2ff18b3dcecfd1c3f6aac34ab1c32400
6ee8cc149b101b6259818e9e7c787c5754584f1e6288121a25b4aac985a0f49d121b93c64cf7f5e7a3d2e881dfbd704e4c39437d48134c8865268dab
5e7a0b3d81726839e73f95460387f3aed3b4c85b8ef61495de1062194f4381bcddbb38d391a9289e878de92c187d8959c9b32b6f780de6ba3434c15d
4b286faf3ccd10ca8e509da9ebabac391cbc9993e14abe7f47b8cfdcffb33182b768cf9490a407d6964e402c3dcc19b82bad577f0a711c07630d2e1e
f41b0c177e6d2142c378d3297f4331a3aef4c1307c41c561f6e16a4c1f01e684b862b2b68756c771cd2186ce51e37c1cc0b2bc6a45720ae5b785e508
a1dc217359687aa0ef15dee167b23425585dde3942c678cb98f8e5921544ec1c562c8d55b0988be19ac85f0a24d6171f0863152e84b364d58b04d10c
4e0b954c2198922191827b2e6db87eb861788dfa2ea86a095c70e468edd1b70aa717d62f8a9c157022a3060e2396154412950141340cdee4a6a2d929
57abe8d0aadd0fc317663f87509e2f86d9104a104d695dfdd6493c564968cc7cecb9f5211e96495b44c7a0eed6b264b624912016b05262b9895c37fd
33e8fc16fdadf9b487321fbb134be6708e43acef05fd2810cddd672697c6aec4926b9e050e1bd092124710fa6a3b763167a994fac22194a5f34a0064
f6bb2f19cb002fa0b94ca4e46bc0d0962557213308249044ee8b1ffdb82ec35a96948ca054286fa0830169509ffc907a59a3ad351f8de1c4d258930c
ab6ee6441670de0091985adc992ff2ef94e3003033d430dccc4d4ed374c7da110f3c88603832cb94a0af33338c588eb60202004e94522653fe9fa987
44bfe743fc307c4ebe6bed88e7353a0c181d2a2b929c4bb3d3b2c55d22dab22cadcc0791047b244700f12c9423562c0bcbc99e6aadc5a4a307801b08
667b448a65e1d001c5bec0218060b6459c58ee2dc10000358160b6439483a74028cfd101c011c9e9f7148207121123960542798a1a37e0c86408e61d
4a61081210310c2f10ca97145201c0e1497d8e301a4ba3bb65a994ca75c6402801b8ca89b525006230d3e82a964aa957d334ddb4766c03a104600565
234f4af386f46ef1741b861764347fa7b5c67c0b001e301caf434fcb324728bf412801d8e4714a1361381e4717cb92e6299387df78030210478600a2
b4f306cdc53237bb3984128078729e333c63617a0cc373b29bdfb6b60000bce42cf92581051e9a8a656620ec19d67a0390c5f3c42f15d7b3da334d87
e158c70a405be019e7a39965a9947a646ddce669abeb03005c3ca7a8dfe3a1996509ab1280f6e494a3c073e746b2589ed08a040040c3670f62e9a6c9
303c67491584120020895673960fad2d615e07f70200aa91e95fd83d4d86e1180600d08f8c2930d4b07250ddb2544aa1d101180b24d376d062189e9a
4aedd4da020068c91db4b64df56138bce000f405617b3c88134bfc4800f002b1e4415475470000900ac41200b0067e0307104b00764c66a62ff80c1c
60ce12801da394fa314dd3b5943bc433e806962500fb264928811f8825003b050b4278815802b05f3e66dcd94b6b0bb8a0c59ca5992cbe6fedf080f9
120078407c252fe2963b2aa5ee5a1b01004964d6023fb7b6804baa8b65c6d245a48702a080dc2a8d5a6b24d00880146d00ec181aa999e4db37b7ee12
cf5d98560e1eac0800a0035aeb2f5aeb5b4608977fd334dd5b0ebb2194dbb4b22c8d79ffdddae1e7b9d6fa95772f000034466cc132bce90000926819
679914bf953b490d0000356866594eb02e010003d37a054f52d546c45c0200a4d0d4b29c605d020006a5c7daf0c7d696004aa9f7febd0000d086e696
e594b714eb86d6fa87b51500001ad1452c270cc7010083d13345db536b4b80ccc4000000c04237b1d45abf9da6e99bb52300a57b030080e6744dfe6b
d6ac5a1bc3dc574a25a57c0300000ebacd592ec918623f25cb1400009a20422ca73cc13cc9c8950900005948aac173dbda12e663664d64000048468c
586aadffc910cc3798c30400b440cc307c86b20d7db57684f9a0b546390a004035c489e5942f98e7a8210200a88548b19cf2b2ab5f80953e00801a48
72f05cc1ac05cf113ee35547e26000003762c57226d352fc8ad53e00004ec40ec3d728a5fe8929e7b906c372000007e22dcb195a1a99940b73fa352c
c76a1f004011c358964b723310c1ca0400e4328c65b984442f2963d1f4cbca441261004032438ae5f46b589ebae2c7708d44132b7f0000d10c390c5f
43d6e2356b471cf7b4d65f3a5e3e006000762196534110fb02882600c0cbb0c3f0358b20f677d6ce383ed3f01cb5ca010016bbb12cd730d4ec418261
00c025bb15cb1906d13cd35ac3da04e0e0ec5e2ca79f8269c4eeb3b52311c46902705c76336719c2386e48e8ee053eb609cd699abff75d6f0800d09c
4358966bb82c4de2b1d61ae209c0ce39a458ce30841bad79aeb57e656d05000ccfa1c57209832368cda9d6fa81b5150030248798b38cc1cc69d2bce6
19d321ef2fe638b1bc1280c18165e9a1c2107d0d2c4f00060296a587794510599b2fdd9f2a626d799abf2f24d2000061c0b24c4429651c387f343e2d
ac50003a03b12c4029656a95ff5dfb3c088607a03f104b46cc307a9aa63bcc87fd46b93b01001d815856842a4cde2f3903ac4a006400b16c8852cacc
3b7e4c3923c4120019c01bde10adf5a7d9c3befc9ba6e9b5e72a72ca6600002ab02bcb92c26e5e68ad9f593b0100a9cfd393699a7e20f7c14f76390c
5f84f760ad3600095084c75b1249381617ec7ace7255c80c99cf0170b010c8f959b9adb5fec7fee4b1d9f59ca5d6da0ccb4fe87fdf2c56ca3cb13e0c
c0813002698c094a20f33709e54b9a4b87503a388c375c2965de9cbf5b3b30540707c16141ce20963782c3850e29a5cc5bf3a6b5e327afe11c027b82
46516f02b774c3e441b0b6028bc3850ed11bf486b5e3277f2c86ea6f91d4028c887170cefd3820942734e486504672c838cb458df1136be72fcc90fd
fb221b102a3c029128a56e99ba500b810c257a99e7253f597b4010ace0891baaacc13c27e80acd3fbe0a4c29ad79a7b58663b30088e5820cd1347c98
a6e9193c88a0263425f4cae3a40c0191640262e920533467cc300725244031d40f5f39bcd73140249981580628144dc3375a7e896078b0090dad5f14
a6f9834856026219414eb6200f104f70099338ce60445319886502c6eb384dd397cc61918b7333cc42273f064a2913c3fb2cc129130396f136026299
0947625f0f1f484011da3130146af62cc321138379c9de8553b12d10cb42a81ef89f954ff3ce2c538380ca8484f1c9467c23071fb4d68f0edaccdd81
5832516188bec5e9344def4944b10aa311e4f47b526954e103436d01402c2bd0c8daf4f1bf4ed3f43f1a21d55a7ff17c0604a017df23fa6b298a4b60
450a036259110a24fe54a1e2632ea7743d9f8e3ea4a70887f9af9720ae3173918f8efedb480562d9087a38df371ca6e760429bfe9905d5fc7b342702
bda0ee92089affde12f4b2f281e5b30300b1ec0043b0bb244e17d7f2c9f3ef2597024cc35d5f1ec5079e7f4bb1024bf99fb5d6fff9d8b7702c20969d
514afddfd334fd7b876e84e370992fd564099aa6e921a54af3bd588020500ab71373cec1699afedbb92cee344df7a6693a3b6483ec97a78bb2c79789
a5c97963f2aa7ea45a514038b02c3b4022396d15865a54a904e3704a59a8a223111685f560650a069665434cf0f22c943185a18c25b2b03a4f60758a
e49c1c34b3f5f82035648b0aeb9d919509b1140ac4b211147bf97922a14c3dabb138b4d67717e2f918e2d985f345b671f3779dc3936d7e5bfa3def63
582e130cc31bb0aa2c59a5401485263d23a701e0e38cd6ea375941b318929f93c5098400b1accc6adeb1e99c14a5007bb6a3709bda9c512cecab9e4b
481773da104c41402c2bb2ca83794643added77477b1944f7ab0764dc42627a1f8d3aff4bf22fa0d80585665612164cd5376b8de478b2580bb12d211
da7f89a928baf80db0c2470010cb4a2c8572d4ac31ab7b98cb062f97135eef20aae794dd69b6089debdc57d75e659eb8367bb8873d01b1acc03a3e72
34ab66665d4ea3f43e560f7f0c59a5129452ff2cb2910f3befb7720c0edb8ff6024287eab00c247f3aea4dac2db611620069be6f59b6c1b7f65c3ceb
c263945300740262c90c593597ec2069ebebc5bf47f0aa7f5dfecf0e86aedf16ffde4bf295218158f2b3b46abeb53a692d96eb9927e1d62559954b5e
5b1f1a8f2b0980697103e800c49211ca24b3e4598bf33646b275b9b62a876f7fc7d2c95e19f80f0fc492972bab67b4d66bf11c9565cecab9a42be884
c382060d80583241c1de7b653deffa97c0f6df731ce2e9eafff7f2121e0a88251f7baebe37c2c3b94e65b716989159cf131f79e5553720967cecb603
bb3ccab4da4732c1f4778381b46d028058d6634f968d0b31f3969e39bc3d89a505e68ddb03b16460e7f3953e2479c58f18ac8d9ae28d8158f2e0eab8
488bd68e078e33b9accd3d81fed51888250fbbb62ca916b7645cc2b127b174bd0c40630e2b96662584598d42ff2d1583bd27687559ced27109e8a840
2c05f0dbd16e9896eb2d1f24f3ef3f9552c84aed07091cfa8221b7000e65595242555fc7bb9691422cc800e135b1f8da0c0c008d9ef4ea0f81ed891c
6d18be190bc99c060bd9ad3bb287553d9eb0a894effff0ac277fc86d1cec9dc3882525528d21270d962f68f8a6b5653006119c736bcb4fd6ab7a46c4
d7b736a10c45d7429f5ba714047e8e6459fe6e6de1c35b547f07095b47109c50fb8f3e0f5df2c27559949cc73f14081d626023bbd0b0095ba9acc408
84460dc35a4e1ba3a177d696fcf3c0db1e01c4d286ad13ce0c6c5d7eb4b6fc424c62dd8d6cf4d706b62e43a3a1d03da7826887080e219689cb116bac
b91dceba8cc8c83d92f3e4bbb5453814b9e1c555d1b280902803e2289665f49bd395612792a0453ae0447a70be4b6b2ded7e7c4e9e0b46f28cd3cb9d
238b55b04d401a4711cb58274549cd9c2d8bf4e628c3f1889012890fe156fbff315086f1cfd696ab045fcc0b30bc6604739657c90e228fb448df48cf
504471795b8873086ccc5bce7cb5b6082326f6715d22d7c786e3f10a70f26c03b15ce0280e95caf388cf7f96ea70a079b2605cdec4d34eb5d8cc212a
39103bf245550ba47cdb60f762d9d292d35ac7ce8b7d9766619250c6cc93bdb4b608416b1d651dd1b23f512f2c12f1cd17d5344db7ad2d61627faf3d
04f057e5089665ec1bf3a9b5258fd8cef9594a0d687a50a31c0a5a6be975ab63e79dbf4b187a1ad14eb176531d6b03fc5ec37004b18c0a058a9cf38a
394e4ae7fcb3e7d0cb8845e2b094eb85520dad758a13e7e356884e4de8659912d674c3da029a7104b18c19da707392727d342c6c6a0190488782ce2d
b85e280d480998bf43eddfccca5c5893c1f0ac156705616d1fac2d0e0e5a1e251a38787ec2ba6a87028653c390fea487b66ab807253c8e9d1fbb446b
adac8d42d15ae72c2cf8585b3449247fe404c96bad4b842cb64f21d428c0aec53236ae2e36142385c4e1e092378a39dfa0c9ab39e732cccc4d297ef8
bda640dc67d1640b62372f406afbef99239d54a7ce15122c52387902ecddb2ecfda62c99637ab848d6fa252591b0794998877d21907f5b1f8ae76ca0
e1f79a94e990357f2cdaff538ac56986b32609c6a2fd4b96bbbe645a2d55b2e0e2f01894d6fbcdff49439ead37f959e1102708cd036dadc8900a6ba9
8d8c18c797a5de5cb21047b598d8fa2689fde61cf548d32dadd9bb651933e4a91a8c4b01dcc30d63b985b217347fc99e49aa01ac2ff1d8c41b032d09
6dcee11d3c2d1242d030f69eb5432ebb2ade4673d2230966d5d1ce0670f278d8ad5846aed0689610822ccc11e2e4cef658e59204f3b1b5431e2f2b0a
65cc72dc1a290a77c19e2dcb981fbd69c7305e499a13929a3aeb71478ba63a945842f20beb76cd153791cb717bc4250fc19ec5727338d1cbcb4b969b
a435d666d8ad52b2d4a442f5dabbb3786145056a37e21bb53f8a870966cf6229ba1093b120e8a1ed1dd271d268d82daaf6b8d6fa115999bdadfcdb05
31b9396cceddc2c9e3e6c80e1e113564e841b9d141349f923553dde2939a259eaccceba541df999cf4b026231760c0c9e360976219e3dcc95c125705
7a686f91a559d36b6baca87bf490b69c82906ee5ff436d527b786e5e88375abda40a80583ad8ab6539ac47cfbcf9170f2e87f57b468e1b73ccebad13
f7764e689b8c199e53dbdf6612ced3c50bea5641320c4ece368e855ae20e76b982272291ed079ab31a0a5a8531ff5d5fdde329d5c836f7fe5e82b380
2cfc92ca8adf1acfe70559b5ffdd95e7f894dafe8bf4e5a131abcab092c766af62b97553b7e179ac4fe472d3207868eb10f18cdc1062058be1900e1e
08657dc8a35a1cb3376009e151d88a024070fa8a238a25b2afb481ab92e24da905de06674b0ce1e459b13bb1544a6d7502bc312b5321816ec9bc2770
1031af0a27cf8addcd596e39773007569f4ae566ef092ec13b245bbf139e95abec7118ee154a509f8a653182de5b9085888519a370b439cb98ac2ba0
8c92ace0416ad7273a1a5b0b3322a6b40ec51e87e1de1bc2b0a22ef470fd55f324f80d79093d2f9df36a8a635796252c8fee5415cae957ad6dc04768
9512a6b416ec6d181e12cb114b0b0c4343114ba9b50db6c1503b925d0dc3378614589150918db6e7e6546b5dadbef7d108fd7698f6f8c5611c3c10ca
7a98b2af8d4f292a37e60ef02ed48093e7174711cbad2c2ba08cdf5bb71f9641b21212c4e112ced4623762b9e1dc09ed0365edde2b2f23964132b151
4e04563cb19b394b7a689d3f2ce65dea119aef6ac0ae4af6f624f43be2f9f9c99e86e14ea1145c497178040c85afa15e0c1b58b0b1c111e62c43f331
20131a024b48b6c095dde8d084cae46e4c711d86dd8ba5f4acd50323c6c15221cb11b8cae1c572da8b582aa5e0b16b085995928af17fb4b6801c7c0b
370eefe499766459fa86dac8aa520771f925b10c9205df737478a6bd78c37d9e3c78f1eae06befdee0f72ec7f7dba26d773e674955ec003f123da7de
55282009e7020e4c75eddcb2249066aa0211e5869b02cb87874099dcc3afc73f8258cea0fc2d335286e3104a1ec87afcdb77b0a3b7f35e86e1f7ac2d
365f9552a1655d2011210fcf0d6b0b48861618788512ec442c4d212b7a707da10f330f8d358435c5acf414ab7bc82655861976d308616b81c1636bcb
c1d89583476bfd8444736bb2ff3b424d7820b17adae1d44f51edb10ccaa7e09a9f5c63a6b00e3f2adb5d0d9e19b21eb7e20191888189c60e9f77e6c5
686d0551443e1b13bd90b0028ed86de890b178c8ca3cb176fee21a0dcb0f1f16514ac388836f10ca7c2851f39650be34cf0e84f22abbb52cd744541e
340f2132d81452db437e748f6c0911bfcde1c383421ca9acc4ab0d27d04db2321193594635870f84320f63286c08e539599210ca0087b12cd750a884
cf0388376c0194d2eb0df36151702e03a5d48f8da427883f8ee43096e51a1a72dff02407be8f10a37c68aeeb94f19008114ac4ccc39335e913cac764
4d42282339ac584ebf9c40d70341ed26c4c89b1415f821cbdcf5224a052142896c0498cfce9bc38702a572d861b88bd0f011f365796ccc956df1efb4
d6fff1c667001158d73d616aa99c435b966bccf0d1e704a26139f2fda553e2f0f91fac2dc04920c01cce1b2620960e162b81d6e9aafea2097310df96
3fb054ae1e665e9dac775736f3db5874c107c43200055aaf9d4008644f84e6c72c6b1d9441f3e9ae00f313386ff881586eb07002dd5e7df26f01a560
8701ab6e78216bf28fd54167e7cda7d1ef4f2210cb48cc5b9a86e6cba4110864074df104989f924822394c45209689789c409f69821d806ad07cf972
c9ee37386fda01b1ccc4e10442203ba882a98bee0830bf815c066d815816e2700221901db041a9ef9675d167e70da2321a03b164c0e104fac331af04
40344aa95bd487e61ca170de740662c9c8da098440769003d58afa4a5f85f3460810cb0a2c9c40af11c80e625904983f84f3461e10cb8a68ad9f9168
7e212b131d1f38590598c3792390df8ede002d98ad0363619a40f6862518c0002ce6b74f3027291758960d2127d023b23261391c1c93e58a8412ce9b
0180583666e104ba4b13f9e080d03cf623386fc60162d909935c426b6daccc1708643f0e14606e5e98d7cdef7ff4f61809cc59760656c571a000f307
70de8c09c41280cad0fcf42d38f6c6066209406528af24d2f90d0ee62c01002002882500004400b104008008209600001001c4120000228058020040
04104b000088006209000011402c01002002882500004400b104008008209600001001c412000022805802004004104b000088006209000011402c01
002002882500004400b104008008209600001001c412000022805802004004104b000088006209000011402c010020028865039452ffbb52eafaee6f
14801d03b1ac8c52eadf4cd3f41f69ad7fecfa4601d83910cbfafc77d34fd1fc5f767e9f00ec1a8865459452dafc87cef09f60280ec0b8402c2ba194
7ab53ab211cdefbbb941000e06c4b2026441fee13ab252ea91b51100201e88651d4216e4dfd6160080782096cc28a53e6d1dd131440700080762c988
52eaee344df7238ee81ca20300e402b1e4e573ecd1622c5000801c20964c28a5be241e29c6020500080162c9000dbfefa41e4929f58fb51100201288
250fd1c3ef153711a80ec018402c0bc9187eaf817509c000402c0bc81d7eafb806eb1200f9402ccbc81d7eaf81750980702096993087fec0ba044038
10cb0c9452b72a84fec0ba04403010cb3cbe5638e6356b0b00400c10cb449452ef2b1e1bd625004281582640f38a0f2b9ee2a6b50500200288651aa1
d46b2c28a5de0abb67000ecf04b18c4729f5a2d1a97eb7b60000ba03b18ce7cf5627524a3db3360200ba02b18c4029d5ba8ced5fd61600405720961b
28a559dc64f100002000494441541ef408eba1584e0080102096db7cec74ded2041d0000462096013a673347903a008280587aa098caaed9cc114604
801c20967eaac75446803022008400b174d030a632e65a1e581b0100cd8158ba69165319412f0713006001c47205925900005c402c17506ca3b86416
4aa957d646004053209657a991a792833f845e17600619f3e502b124a45b6f78880e43cfd85e100062f90be9d61b622e8f4169b550500988659f4419
39d44c3a0ce4f04e29f508bf873c0e2f9654fb7b88a58548ae7108cc08e2efa33782440e2f968cb5bf5b50adfe0f9081d61a73964239b4580e189283
f9ac838004d0f238ba65395c480e86e2870109a0857158b11c78a50e86e2a029c6e164d2152aa5f4ea2f3ae7aab1941ddf77fd89b5a895d6dadab877
c83a931a80be89d65a09bf4427e6c1706d0ff0526b2d26a9492b288faa490ff85c6b5d65aa88e27697c264fedf383b9356b16df545324a7256c59d6a
ad452591f9cdda720c8615ca893abad67a8470a7a6180b486b5dd5f23659a01a386166b13443f15af3ea2f6a4e4329a5de1786bbdda797eb0d297dfd
70c37049e9d70a4080ba9b1719d6abc473b4f0883fb1b630605ee4d43e5c71c1dfa5ac5e3be29ca5a4f46bb92040ddcd45b4406531bb5ffb1c4bcb55
295545d4a66962b7d62898be46d26c0989b88f25969d6bea808634b0fe9a9ca3e230fcaeb5a5002a81522d983ec599548bc388a5849a3a9c200e6f9b
16cb581b9ca3caea32330f48ce9913e3489ba6e99bf5a14848c86a9740e91e637c24cb528429cf08e2f0b6b9d6608efa5aedc50d35e7eccc90df441c
68ad6f91782689260965132123a751370e2196a8637368fe6ce020f8a3f2399aad3433a2696d0c132b94a7d3347da0ffe6d275aefe2896e52eebd860
354f342d4615dce7588aca88553ebf51d88fa23f1372f588feabc88a3db3be2598dd8be5ce4b321c2e603b97162163cc7d6d5467e43989e1adadf848
adf5dd54c1ec394a3c8265b9e7920ca82b1e4f8b9031cebe76452c2995a074de69ad93a6234830538058d600951ac19216e1278c7d6e7dadd2471146
287363423f585bfc402cb9a1097771951ab919c4e290420baf2d4b9f730c61252f44382f10ca69b5465d2c7bb62cf7162ae403f39609b4186d1c6d44
933af4767c7f88f6daa5581e2c54084b1fd36831daa8720ea1153e4b42818662af96e52e4385c026e7311fa2a579b944796f2b2dad95385c3d4c71b5
dd89e54eb20a258178cb4b6245b0248a207615498da5b5e2c4d231b7ba5bf66859ee21ab502a5827fe93688bb160aa26e51cdc56d7105548f7caaec4
f2c05985f61c4b1a8dd63a2534286b9d71a23302256d77c4de2ccbdd641502d519c14ab39c27c2a65cacebdb33bb11cb16e9b8c010443fc005c3e4e8
20ea0ac97b6b2503061bec422ce96d7be8f91c6456ba2465fe36772d778b734c9ef5e187f13e4b632f96e5d005c8988093277dde322b1e3271de92fb
25de3d09ee51195e2c61515d82e0f45f44c55b16929d591c8cc91e2c4b04a08335d1f37a05738ad1c3e1821220ae6138e8c4d06279c40074b04d62ed
f02cb14c1ceeb34e916011421f46b72c8f1880ee050fd11562bde225e166b15e71eeb5e270f2746058b1ec5dbc4828082b214cf9026b23ff396a8b96
cf7a8558766064cb120e0d1b3c445769e1e8a9768ec0ba6b2cbee8c09062890ce85e1056729516d31298fa3808c389e55132a08372c832dbb4fc4ac2
cf5a9c03c86044cb12562548a1bae51799295c62e25e90c0506289658d2015b2fcb602c839acbe2def3b6a250dce68962596356e80e19e8da9616d6d
e43f07da7de70c2396108168e01177f3d8b9959793563783e7a13d23599658d618071e2207b4aac7e788619907d75a7f6ab8661cbf73638610cb0a39
01f70cc2873c041c316c4ec3c0901febbc076714cbf28db505803c6eacbf45162127d6392a447148b02c7d2f9fbd9df302f1628964198013f28edfab
d9a8ae7324e6c08cc167c116919821e98e52aad8cb9f98b1fe4eaff9da112c4b24cb00ac50c6a0a7355b95e31c1ba2c0b630c39c4729f54529a5a769
facbfa4098cfe67ba6163b2d18893de72d93df81ce995ad8ed239df33d8758c7225a2c0b8be103e0456bfdb6b6879cce51d58a2d81044793f3b474ae
dbd462ff4ec7f4fa18485435850196e67778b8106bdf3a7a36a45b9625c5f00fcb86450208e321d75aab9aed612cccdae7c8a12021710c211f43ad67
fa5a6d2b53ac581eb806380748eeb00f6abef4f6b8fcb2ea3d49b62c91862a1f8825d86277c64885a8862bfc666d1180996c96785d03816138086284
45a9cbd9816f14da348b8d79feaecc01d2e7d7fdeac1eabfa906ce299d677edeaf889de79cb7e8ef3aadb7bfdb2a5f8448b144607531c870b30faabe
f452e7521d965bb225c770ce6e881b8623b12f0b78d900c08c28b144625f00ae80797b4148b32c615502004422462cc9aa44625f0006c13cb3739079
e4df9750c0ba74245996b02a011880c532c5ef8941e6662efdcd42586b06c6b323422c615502201fe37c2591e42a43fdd7bcaedcda23102996e5776b
0b0040044aa9572492b59cafbf93688acef2df3dce928a900100044222d98abf9552e78124cd5d916059a20819002b1c2b579a4229d45a0ae5ccb54e
e7dda4ab58c2aa04c04bb7d2b994bd27d6883935ab727c7f9435de57fbc88b44c1ec6d59c2aa04c04d17b12403e6b3b5c3010964d0023659e369589d
5cf9529a6076134b58950004e9e5ec883160ce33d778bb6a13056991d437969e9625ac4a00fc340fa58bb5e4721d30aeda44119839cc5739e7e3a68b
58c2aa04401609850193adc325549be8a5b523cc1fc1bd8de86559c2aa04a08c6fcced175318f09cacc322b4d6c9155b2564236b2e96b02a0160814d
3c12926d079d3989dc4efc7cf76c643d2c4b58950004e810631995ff9486d02ce4d451ef5d97aba958c2aa04208a18b164118eceebb253c389bae6f7
6c6d59c2aa04609b6662999035e8d4da52484ec988dae56e4334134b589500442331437aad21f099b5254cb730a2969625ac4a0098e028e49518bf58
4b2c53735a767b9134114bca5709009045f7f84549d51bb7686559220b3a006373f85afed5c51259d00188a7776ab600357d0e4959897af93f5a5896
b02a0188476a41af9a5369a9e14bbb154b589500c49352006c2fbc1fe13eaa8aa584f59c072539d92a180af698c708aa4d0f8ce2e4a96d59765fcf79
500e3f19bf737a888bd4b9d4665413cbdeeb380f0eda7edff478194a0c946f4a4dcbf2f08ddb1131d9a5413c4aa958e7cea15f86bd86ed55c4125665
77300c1f9328b1e4c82909d2a96559c2aaecc848ab22c0155a3f37a9f18d879eb76417cbce299f0000f1a426a590500ba74724c005352ccb23c68901
301c19e51da29204a792983ba25bd03eab5826143d022086c34c27240c713f585b1a5269a961b400e66458e782dbb28c297a0400b0893534b8a7b9de
595bc2d4701ec6de7beab5b2a2b48e2a15bc09853dbce97933e08253adb5c889f8d8bad40b4e8ee2ac4aa8d9adac8d8dce5deb1a7ade7b0a9c962584
5206f084835492caea722e634e882ded6a554e5c9625cd632013ba0c445a633427f7d1da11a0b725d10a72707c8f395dad36c9b0faef71547b1cc5aa
9c182d4b08a51010633924b1a5156a2648796d6d09f339b83702a5546cb6a1d42a905528164b948c00911c3e11438058c768b51866ad756a2d9c1c6b
74f95d53a5f1a1b5c3e6548a01c0615962691d88012fd572aa2ef8c819eae608a652ea51a4657a2ec959c9219648c32607c9792cbbd57bde0b1c7384
11dc48fd8211cc58470d3987feb67638d05a8b7ac116892512668843f25253e40b7090588eb63a94a42359304d340c89e62763391a87dee2ef2dedd3
91c6d5b944e75e9137bc64ce0254416c5c624e5f3982373ca15dbe69ad9bd69ee9f47c9f69ad458e42b22d4b696f44004ff8ce69febcd1cbaa65e28a
c75285722ab12c6155ca43b22506cbd2863cc25121383ddb22250e3413b1abce966459960951f700e4d22d15574386486768e63149ac6f333b118d48
aa1184722a18866369a33cba2f0763e608954163539e89887230197f8c879a84f3b9f58138ccbd3c1d4924677eb3b66c502945132847ec1c7266866d
c4effe425cea43adf5ab659fa32985078e785ae35dffd428eca92ac962894e2c13e19d1162b922a5a200099368a8ffedfb374b75f0c0b12313e1ce9d
4fa971960770ee443f47474928229da4394b04a18be54cf8f521203d1fe9bfed614875f0a0d3cb04e53c0622b1a85fb7c813e39f30a562e6d537cb3f
ebc37ce7bc4ee7fce1386fd7e58fd1c370aaaf83b21102913e4c438ce555240cc11df5b21e241a431fb4d68facad0118ce79de73bd788a5862ae5228
3b14cbe64bfb5a911ae02d28d9ef9ae440720e0de9d9d7a386e114160064227a4e8bd271a5b2e7daf329f3fe4fad2d0c64867215d1e39cdcc4ce5916
674506d5484edada989c39b73d3b12a36b6f6badf7fcd2188ea8613886e072c17ce53890951d95cbb1763b60189ecea66589702100d88816ca5a43f0
19129d7bb46cb1c93afc1ee7e464d3b28455299ad739b5535a02cbf217d203d129a4e9776b871b964c419494273ad78458cb12d98564338050e63c4c
bb0cc24e1ca125d5f1e6426bddfc791f695e362896c82e040ac90996dfab5323259e10953005e2154b94b815cf8701ae3179c5d708492352514a258d
004c2a346b23e88e572c915d483cd24386c02ffe4a688bd7d61620829058a2c4ad60a45b1fc87bfa93d4795be9f3d047c629968e359c401623384172
fad008530ba97c4cf87c17c70e88c3193a847021f1dc1ec0b2cce943624bf9e640d6f5d784afdea0badd3daf39f677632b32364a6e4fcbb2c4f0493e
7b7500ecb0946f8a504ebd851284b1c472e7eb72f780f8a16a66f28c5d911163fad8da0244610dc3310497cd08ab5b9452ff64380893f3234a26f539
92f2bb6218eee78a65991a0f0680879c488addf4bd8c956fb965654143ae5896b02ac5f35c7ad0768653e3823dad071fd5aa9c60590671cd5902a10c
b2bae5bdb5e54028a5527fa397d61620924bb1544afd5bfc44a21925062f3ab9ed82aae9c81af347cae9b4d688691e84a565f99f1ebd318423beb447
6ef991bd640427c7560a7b7a49ec9e0bb1544afd67476f08e90c128377d82138cdd52639b65036622c7e8353670846b14072bce07bf104a73ab56e5b
5b8068e0e01980112c90dcb0b33da464534aa55ad4e748c3361eff82374e3cef06b9ce943464bb8186df0f53ee476b8d5cb103f22f78e364d323d57f
2a05f904f6b0c42f75f88d7c9583721194ae9432ce836b476f0c817cd35a8b4f6c92db7f460f44574a7d490d95da53e962ae7b192d285d7c58ca1119
4128899c17edd085c92851466a4c299c3a03732196986c16c9f9081799115b38337a51ae94a4be132d0fc473b662a444e3976bc353ebf782ea884ff0
3b15e413187908bee75ae8ad87c423ada3bf0c1d4280ac2c0611cadc6b1cd6b143f394a960f8ed20631d7d57d67196082392c1280f5756513badf590
2b7d68c8983a4ff961afc3efdce5add3af52db49ebe87b83e4bf021924c16f6e04c5cb11c3d548183e5b3b36186dba21f1f9cf8ad620a1fc6eed8840
5a8a362422ed8b78ab92e22ab342cd0615caeb47104a22a56cc94d6a9b68c837922594bdb12ccb09d6655706b12a73fbc7bb1182ecd764deef3dad75
cefc6657a87ed2df29d710d367739342e79cab162ecbd270626d012db821bd954b4a8f1c4828df8d289453e67cb26923a5d4a7f51ca6115edaae3984
b2374ecb728275d983f311d60c1fc9aaccbcd721565d8510feec779bf30e8965f6242c4867e74e9dbd3b3a2ed9432d21f2faff69eda88069afccb636
6becdfb7ac35ef1b86cfc966c5d7a8de09a7d26f8396f7e5e60f182a24edc842393574c215b697093bfa4853004d2c61af6579f9010cc7abb373a7ce
502252709f3706c9661f458ea327812b534e4c1a53bdf2a9d7b25c80d50775116f75d1f03b977b7daf3e1e08e52fc8d15323d9c9cb4a73f3d5e7fb37
2dcbe9672732f302f7ad1da098015276bd2a5869318ad3aa647e7ec810a15872d2d079f0f60526cbb2faef1025961386e3b5109d2c2377d5cacc20d3
0b25c3cd5d0be50c83c3e724e488c934c68cd56bacf98be3b698674d114b78c779111f6252f88214bfac51296592c7fc6eed88e31042b924b1bd4c8a
c14721915c1ddbd7d79a7bbd7d448be5f42b20f990b556b8d953c66c1703dc5f497580c309e51a8a8e58e724fd87846d57f3b733496239f1ce611c99
ea9ebb1218a65cc43a3c1896ddedce9903e2482e85abb546098a42840b65a910bc142c94af4a84d258cb10cae3926c59cec0e19387e4e12943e13aaf
c7b3377b9f5600f549b62c17884ffa2010b131951c153e250aa5b1260b85f20c4209a612cb7262082d391a521f3aa652c8a29c1e4c29c144cf2d83b6
fc567236f37028a59ea2d0d936828592633ae5b530a1fc27b7e4c5023872c0154a86e11750a1b377d60eb0e4a9b4d63071b34c426986a9d9392e3931
7180744f2542f90d8e1ce0a2582ca75f495d21986ebe49ab9c49ab5658161848888e302b4c48247303cc671e8f9e8b12d4a368ce724de18a885d226d
f8cd1927dbfbde0ad7ad2f19620dbb54687ef80905a9bb962d9ed2b2c44f1256e2e4c22a961304738d98b5dfdccb557b0a25736297c7a396e6ed09f5
a72f05531edfa6697a3652dbb30cc397d0905cdc1c5d075e0b12ca57ccebfa9b878d99c80be3b5a7e13687509ed2dc248432019aebfe41fda9646ed8
7cf76f4adefb83ac53d1b05b963354f2f2a85e7211c33aae8a7a2b9a7989c97a31239587d6ce7c30e4cea451b9896086a29e5413cbe9c0719812e629
99c267d65417ca4a0239233a259e641ae784102998ecc3f025147b77b4953e5def77517a7418a1a492a9ffd0757faf2094b769c80da1cca043f29c75
3623115415cb890a9f91a5756eeddc1f27bde2f3163186ec19edb9e30e4d7aaf59d4e99affae20eee710c97228cc2c5628cfe81950f39ff9ff8cb042
91625975186e9d6cdfe9ddbad4c5564afd1fd334fd87d60e1e8ae7f76858fd8cfe4a9754c6f0416bfda8c1790e41e4c285b39878db8439749189b19b
8ae5b45fc74fd31fb7511b5e8a0e257abd4b45a1ccbf6f55b0044b4108103391215ac946428cd12471797073b19c7658a2a2c50f4bc3a1b78dacb351
782a6d75d49e88b12a73fbfed6b1258a65f5394b178b79cc53c7eea1a8f9a3cecbf816f37a10ca69fa48198e1484b21e347ad9e2f9c6fe10c395d8ee
62595eb98081c38b38859286ba4fb0fa2919f3c27d31f2323a8944869e1585f884ac4b58960e4c78d1a05666728810ad4279429eeb1fb3d5489de623
84320b33a7f671d196c83fc9438b39e90fd616c174b72c97a0dc2e60c63c8c4f906e2d9d90d5b7a028f28046531fad1d70f0c473f0a592a00e5d42bb
462536d769a9a8f9ce8361782466e29e1a0b39320117bf2f86ea104d267c62170b3de72f691a4ef4549c48cb720dd2be814a4405531f910c11149b00
830b9196e51a337ca237d06b6b2700f9dc59589b2297d80dc4c7bdb7e3106239636abd9068225f26e0e62386e857c8cde530b7e3ee62608712cb9979
4e938473a8f003209e37f4b0bf38f84f557affcb39e25d08e710739631d022fdf73b4ed401fa70d825954cd53fd7bcd45a0ff922da8d582e21e17c55
2981ec51998b4e99951dff704ce6d3fcd62dfaf315bb92c2ee1d186b2a25909e313578ee8e1403bb4bb15c430fe50be10fa30466417c2b210724fd6e
0f1aa6778ba159590d0954b22e97884cc7e6e21062e94229651ec047071650238caf464b6b46a386674c25707319e6012fa5e10211f1d31d87154b1f
0b6be601e56fdccb1ce83b4a38b1abace1b444f64527f11c76fe2d858671cea2e35e21969560ae6d9dc319ad8bfe22a95d6ad329efe7ee87e6a175dc
cc8815cc2143872433e7a0ec2494ef16f54fee1e4d28a79f6165ef4d290c0a2b7bdca8f6d377a5d4aeb3b41be756a35a5a77688a4c1cb02c99a854a3
3b06248888801ec0bf1a9c0a56260348a4b15328c4a2a5507e5b542e845046a0b57e450fe00d6abf5a7cdf7b4ecd85957962ed64829e2951402c0ba0
44be356a74fb784e02790be55df3a09226b72a67b5faa341c84d7766d1acd496d20ae2412c733199ce1b85542ceb5f230b38238b042d2f6b1cff4809
3ae6b6e4144e69ebf431679948c3b9af53ad3532e134a4624ce1216b995358d797022bb1b86e3d2710cb04c89aac1d9202874d672889c69fcc5721ea
c16f4d6e289d24470f86e11198a114cd41d514ca9770d8c8c0049a57c868758d86e5e2057359488ffe8ad7c4d32869b8f2b74b60596ea094fa527915
0f2c49e1541851dc931c03eb714eb12d47f41cdf092ccb01301600fda8b584f21496e418d0f0f91ee3c57e1e305f2667a0f890b5b560593aa89c3ce0
d07357a3432b75b852ff8974fcf82c3f4e2bcf778e9ae72c0562b9a2f2b0fb36e223c787b9bebdb8b5d03e213bba586218bea0e2b07b76de40287700
05b67315d0bb4373a2e2e910f7d8625d7f34b02cf92d852587c97b785418fb8e98e99990d5c765e985ceb1e0b9a4851887b72c29a5570da13c8150ee
9f8595795678b3d746b030398a8fc5ae9d97b662edd09665a5a4a628dc7f50981c83dd2dcc08abaf288c28d2aa14e7083d7259891a8e1cd1f173a03e
4cc3f2ae4211296659f1c1093196e29ca1871c86d370875328cfc98103a13c388b617909d7e8652e99b92e7894602e3274c57026d1197a38cbb2c26a
0c5193d0400e0ca3972e719805e9e5ce2971c692a1d7832f39945856c83178a8b2a8201d8620f6d75aeba665167ae6e2942a94d39186e1cc1d601e76
43284110b20c4b96f7fdd121272667029168240be57414b16416cad758ae08522047c8d38246fbd8325b1109fc8d8641e1afa50be574846138b350c2
db0db2290d2dea252815eb860f95e07ad79625b350de80508212283631dbc2ec3597b82a19f1942100ff291d6fa84a00bbb52c393bd6084304300e85
16a6c8450f34afea123f33afffc51437b3f60cc62ec592313c08e9d440150a4b57b025e205f1ec4e2ca9de3047194d2c5b0455299c0b44d85a637635
6749b5423884f21d8412d486bce4b9f37f3592bf8000bb114b2a519bbc5ac0016ae28066d04b392b4467802591bb6217c370a594e9709fad1de9602e
0874a1c02179b207e7c9080c2f968cc9572194a01b25fd18d11a6dd8c3301c420986879c352739f7314a598ad1195a2c993ac93b082590000da773d6
655fa38cffa022c30ec32935fd1fd68e34e0cc01e2c88d13c670bc2e438a259343077194402c990e1f14c8abc8a862597ad15899034453601020d94b
2586134b8ea58c18ae8011c8cdb48efe5d87a11c3cb49e1642090e41ee341147b95a60338c65c9144f89f5b460289452660ef26bea35c328e06724cb
b254284f20946034a8cae169ea6523f6929f21c492c2844a7887256160543293e45e23ab1430217e18ce30fc86e71b0c4fae771cc3713e46b02c8b86
df104ab007281c28393b5187ca90bb45b45852dab5126ef7bd0300f8c87cf17fb4b6802ca45b967f595be2794793e300ec89d7a9f742357f402162e7
2c73037267305703f64ace0a363c0fe588b42cc9a903a104c04d722a375897e588b42c0bd77e3fd75a97861a01201a5897ed116759167aefce2194e0
20243b2fe1192f439c65596255e2cd098e04accbb688b22c0be7555e5a5b00d83739d62556f56422cab2845509401a19cf0c56b46522c6b22cb42a11
7c0e8e4aaa67bc28c5e19111635916589548a50f0e4dc6b383da531988b02c4bac4a082500d3f3c426f8ddda023611615916589578430290f70ca156
4f22ddc5323713f404a70e0097642c0f86a3271109c3f02ca144a8100057480d3887a32711099665d605c0aa04e02a19cfd253ad358a9b45d2d5b254
4ae5967a805509804d6afab637d616e0a5ab6509ab12005e529f293c4bf174b32c95528fac8d71c0aa048009a46e8ba79b6509ab12007ea8144b4a85
0178c523194d2c4f33cb820270183014af439761786e1d70082500fc2013511cbde62cffb0b66c935c06148083f221f1b691303b822ec3f0cc213896
6701108152eaee344d9f53da0a43f16d9a5b9605437008250011e059a9438f6178ce10fc9db50500c00659a32080c852b86b905908806452e72d5f58
5bc0159a8a6541203a00208dd4e9ae87d6167085d69665ce7c656a6253000e8fd63a37ef02f0d0d41b8ed29d00b403c1e9bc0c31670900a80fd68987
692696993f446aca2900c02f529d3c10cb00cd86e14aa97fa669ba69ed0880610100f990819294b312cf9c9f966289f94a001a83794b3e24cf599e59
5b0000a0134dc4327375c0336b0b000074a29565993c718c383100daa394421a440fadc432673d3800a03d104b0f52e72c315f09000fa9cf12c4d283
54b14432520078f8917894fbd616708148b144e17700d8c0dc3f13d5c51213c600803dd0c2b2845802d00f58964c4814cb6fd6160000e84c0bb14c9d
30c67c2500401c121d3c183600c0041677f0513d910616f203d0173c833c20f92f00004400b10400fcffeddd3b6f5d4dfbdff7993f1eb84880900202
d8400a91f00b10d5a513f50a4475e948764e25aa73e59b2c5389ec5c1820f90a4416869b0024cb20854820450a03e40692224841b248603bc504235d
4bf7d29e7d5887995973f87e0042cfb3b76e711fd7faad6b66ae41071c2c818231cfd99fa0074b366e07508ad0c972d3b965351a680048526a97e17d
17fd03588dcb704f421f2cfbbe51cc09039024067880b2912c3d49ed6079efdc02600cfa537a42cd12003ae0321c003ae06009001d70b004800e3858
02851ab8d491e6db4b70b004ca35e460c95ce725421f2c79e181e91c0df8cd7c6797205902e5da18f0ccae9c5bf053e88365df7993ac360026648c61
aef312410f96c61856e40013d05a6ff1bafb95da6538c912f06348bd122b50b304caf465c0b3ba736ec16fa91d2c59f40f4ce794d77e39b6c2050aa3
b5b63b143cf77d567cf756e3321c28cf31efa97f3192e59352eaad73c7129cdd8071fa5ecd35f8eead162359b2220088eb6cc06f3b716ec11f6224cb
3da5d477e78ee5de303116f0436b7dda65649c54b95ef064698ce9bb7ceac0b905c0d0efdf913d10cac17048e28448718067cfb905c0687307cecbd6
bff7caabbb5ef0cb70c5f4212069b6efa53186b185353858024007b12ec3af9d5b002023b10e9617ce2d2b0c6c870f00c144b90c57fd2fc5af8d310c
f4004846aa074bea960092c2da7000e820e6c1b2d784583a3d034849cccbf0be6da32e8d31ace601908468074b45d221a9ac0000200049444154dd12
40c6a859024007b10f965f9d5b56d05a73190e2009512fc31597e20032c56538007430c5c1b2d73a712ec501a4608acbf0de3bcf71290e606ad10f96
8aba25800c4d55b37c706e5941f6110180c94c952cb9140790954992e590dd1b592b0e604a534e1dba746e598d3d42004c6692cbf006033d007291d5
a474ad75afed2900c097a993a5dd3ae2bb73c70aa44b00539834591a63ae9c1bd760450f8029a47019de77a0e7dcb90500029bf432bcd177a04729f5
d118c3e83880685239583e29a5de3a77ac40ed12404ca98c86ef38b7aca1b5de5dfd3700c09f2492a51a76294eba04104d4af32cb79d5bd660641c40
2cc9244b45ba0490b0d456f07c746e5983f66d0062482a592ad2258044a5b8367c48edf2c9b911003c4a2e59aa81e95229f566489f4c00e822d5ae43
6f9c5bd6ebd5791d00fa48f260290971e6dcb186d6fa78f5df00806192bc0c6f30d8032015a937ffbd766e59436b4ddd128077491f2c8d317bce8deb
6d68ad8f623f5600654bfa325cfdbda4b1770f4b2ec701f894fcc1520d9f4ac401138037b96c5836642a111b9c01f0268b83a54c25ea3dd8a394dad7
5a6f39b702404f595c8637b81c073095acf60d1fb26e5c319d088007591d2c8d31b661c69d73c77a763a111b9c01182cabcbf0c6d0cb71a5d4a13186
411f00bde57ab0dc1cd138635b122a0074965bcdf227191d3f71eee8e631f6e30590bf2c936543066e369c3b3a60841c401f591f2cd5b8fa25074c00
9d6579193e67d07422c59422003d647fb094c19a33e78e6e36386002e822fbcbf0866c5af6d6b9a39b5763cce6040f1b40268a3958aa91f54b0e9800
5629a166f9dbc8011b2ec9012c55d4c1520c6ae72638600258a8b883a54c58ffecdcd1ddc6c8cb7900052a3159da03e6d58811f29f386002682b6a80
679ed6da1e343f3977f4f346d22a808a15992c1bb23be48373473fcf5aebdd091e3e8084147db054bf0e983b765a9073473f37ece703d4ade8cbf0b6
314d375a988b0954aaf864d99083dcd884f973a45cfa6902a84835074be5ef80a9a48e79ecdc0aa058d55c86b779ba24575c9603f5a82a59363c26cc
e6b27ccfb9074051aa4c960dadf5bd52ea9d73c730a44ca0605526cb864c2bba74ee18a64999a7133c150081559d2c1b3258f39773c7381f8d31ec55
0e148283a590553a37ce1de3b15c12280007cb9691fb91af620793b6386802f9e260b940c08e43f6a0b923fb0601c848d5033ccb48c7f5bb25778f61
e7763e32dd08c80f07cb258c31bb239b08aff35d0e9a34e80032c0657807111b011f1a63387802092259762097e5a33aaf77742e69d3fe1c45f87d00
3a2259f61070b47ca591bb5602f08064d9839dfa2307ae936c1e34002f48962368aded14a0b7817fcd832ccb04302192e508c6982da5d476e05fc35a
732001244b4f422d97a45e09a48164e9896d9a2107b6c3229e10803f70b0f4ccce939483a68f09ed3e1a1403f080836520c6982b39686e8f38e851af
041241cd322259dab8dff53752af04d2c1c172025a6b3b8a7ebb6eda11074b201d5c864fc0b668b3d38e5ab5cd59752f0290190e961393dae656abbe
792d8fc8d7de40003ce0323c03d2fbd2d63b9fe4cf2b1a080300629273d18152ea9304fa237600c13caec2332057ea9bb244dc8e0b5db53ab9d99f7b
dbd14d9a22010030980d9076f296d6faa57dae915d91f66cd5d8187340a8c422542c332641f274c5cc4d3b547425154e0e000080dfa402d9fc6c2c78
65665295bc72ee019620581644f6a1385e7280686b02e72d43ea0050b60e01b2cd6e6b7cc0b90143112c0ba5b5de916ae6871ecff081399c00901f19
c1b2c171b763806cd8264fa7c69863e71e6000826525b4d6f6a071d4e360d3366b0da9df3af70200a2901e4ccd229a77037f2755490443b0ac901c98
8e57cccdec83617500f0ac357cbdeb61ff605b1c38b67b7838f700be3fbb044bc89694c73d87cdd799352bd899f80d002e99b2d404c8a1d5c7455e65
2ad4290b37111bc1120eadf58104cdb157c9cbcc5a554e4227806249786ce63dfabc789f77295549468e30298225d68a1034db7e573a257872b50d20
69012b8fab1024912482257a8b1c34db5e2574de4ae8bc77fe06007836b7e2dac79cc7210892c802c112a349d03c083cccd3c56c2e78720006b09604
c7ddd64facaae332af326a439044760896f04e86858e3cad3af7c906cf7b2a9e405da413c64e42c171de4c16db5c30fd07b923582238a9061c48d89c
6208a98f07099ef7543d81f4c9f1a51d1a7706f6eb8de95a566cd31718c521586212095735bb7890e0f9f387930310861c27da3f534fb719622673d2
afa846a206044b24431a02dbcae6a742de95592b80daf0f9440514b56b55189b9fad4c03e322afad216dbeeba812c112494b68615068afad10fad4aa
8652e140f2649385cdb9b098da3c46dfec77f64286b40991802058223badadce721c46f7c99eccfe4fa5d4ffac947a9130fac2a2240c258b5cda3f3b
12184bbfb05b672621924a24b006c112459013e29104ced41708a56026c1b4f96982a9a2529a2fa91caa56206c42e2660515445faea5d50f73228101
0896285aabbab997c14ad19c354155b5c26aa3bdb8a9ea8aaa2c46d96cddb4dbfadf4d085404c128a8420201102c519dd62e1a7b052d1402b0980d90
ffc17ef59552ff835c60dadb8e8c31570bff0b0083fd132f1d6a23c35b4fadea90921642ef95529f65ebb4191f0c202bf63bfc5529b56d8cd1ad9f2d
63ccff688cf957c6984d7bbb4c9bb9d05a1bf9b9900b4e002351b1443564fed9456b0ea65dd5b9d7a50fa5cce16cf60aa6ca094c6326f31f6f7d551b
b5d6b6c7e45fad9bec1ccb03e65702c3102c5134a942dccecd57f37ae22074025e357bfe3701324ac093f9afb77373b14f8c31c7ce5f06b014c11245
92fe97e773cfedd2187310f3f94ab0dd4d788f62600ad7a9eed9bfe462f4412e4669e505ac41b04451ec5ca905fd2daf8d317b293ecf56f06cf63aae
bd5f20f2376bed36955c70ec4a4622eee72a9876facc2e0113588e6089222c0994f624b095fb5ca905e17387d6499848334c5dcd3ef94b463f0898c0
12044b646dc9415fd53837aab56bca6e45dbeac18fbbf676a23504c6beb4d6f70bbe4f3660eed00713f81bc112595a324cd5784f25613d59acb0d50a
a15b54438bd034abbf6deda8f444f8194f6b6ddb147d5bf00f253bdd06888d6089ec68ad4f95525f163cee99540f681312900ccdffef4aa97fbee0b7
bc124c477b9040b870bb4df959b46de99931e6c8b9155e49dbb29b25ff2617b5a81ec112d99040b3eca43ab38d909d5b118cd6fa76c962a34101bf35
94df98ffff4ab63adc71fee33fef8f35fc7fd7daae727ee8d8fb7eeb6b020d6d712292cfeae392df48c047d50896c8820cdbfe58f2580995135932ef
acb13dd5f0eb82a6d721440b736b9e0fa172026bc2e5833166d50510502cb67444f2b4d67b2b42e5eb9a0a16029293e7eb92dff02855368c201d0f96
85ca4b42e534e4a2e9e3925ffe4e6bfdc23691a811c112499355dfdf573cc65de6544e6e5578bc91f7100368adaf16b4d16acc6237fcc79f64f5fc99
73c72f76aef113e112b52158225952ed5ad44aa871c644f9e9c97bb0ece46a9d53b9ec4f2a95abb60825542640e653ce963c928dd6c22ba00a044b24
49e62f2d5ba8d0600830111d162bdcc87b8a0ea4adcdb24aa592f636f49a4cc7aacfff5b59e8065481608954ad3b105f32049e9c55554bebcab9050e
59a8b6a85762dba9730b26638cb95a51b5b43ec8c502503c822592237d2a17b5146abb706ec1d4d605c7779c5c3b591bc0a9562669dd31e91b557bd4
806089a4c8817751f3f33f70624d4fc7f7e41b8b199693e0bdeea2eaceb90529e8f2f9a7d28ce2112c911a2a5a79eb127a788f97eb326f988baa7c7d
92a90e40b1089648cdda6a25159bec112c17907ead6c87593e56f3a368044b24a3474b9a45db08221f1bb41f5a686fd18d280ec1124523582225848d
fc750dfdbcd7aeaeaf09af5dde36180e47c908964849e7132615aff4f47c4f78ff5ceb16ed34a8d8a7a9cf679a056c2816c1b212765b3dadb5dd5ecc
acf9b9c864d52ec349e9212c0ed4f742890bab24f53926f1fea15804cb82d980a8b5beb74152b646ec5211b1bb7d3c4bf84c39bcedd3b626397fd5fe
0244c402a884c8c2abae15e7a0ec63b1bd80ed6e3fad9f2bdbca8a3e9a8881605928e987f76c9b528f7886768fe794f7b9654bc74448537b0cd7f77b
f689aa6552fa7efebded1a66c36253409022c277e9aef1a1f5f34976737a94bff79278e10019235816480e18ebb684ebcaee9692ea567c5f98043fbd
ae4dede7a47cc112ddc0ed49d97d2a015aebe301d54a2f9f7fb9f07f1c5040d890c281e10205be112ccbe47b982c56756448e3e75b86c42737e47d7b
726e41dffeac6fed9c68e7564423c7c5de5340c6ee1c26d39c5e468e48356ea85ec227826561e440e7e36033852107db0daa5fd3b173b806ce2d4bb5
0a3ea5219fff7da9982132a9d4df0cf8ad3e367838f5dc4cffdcb905188860599e104d965f63eccd2dbf63e6dcb19eaddc50018b4c42e590d63777c6
18de2fd7d0b0fd17e1322e09958f037fa98f63a9f74538322f1f188d60599ebe73ddba8879c019ba08e4ad4c4867583c0209f243fb29327cbb8031e6
7e4435eb2f86c5e39051a1a1a1d2becfa95e04301c0e2f08960509344fc65697a29db08c313658be3a7774b321ad92988c1e885d2c252b4f87b65699
c5fc3c65684ce8d8975eb55c5c052295e121c3df8d33e79661425cecbfe3d8091f08966509112ca7b88a1dfb3b6fa8def8272d857e8cfc87196e5b41
a6835c2fff1b6bbd958b2bf61df7a8e9093cb257ab9d52e4e5f32fd5ed4be78ef1a85a6234826521a4ed8eefaddecea6980b678cb9f270d0dc97561a
b4231a49aa942f1ea6595cca7b8b158c317b23aaf68defd2db90eae5489e7a022bdfa1cd18132204ee3bb7003d112ccbe1fb20e3edea7a0839680e59
c833ef07732f876955697e785881fa10e844582a1f4392efa47a49f57e003b2c2cd33e7cf4043e09745175e8dc3212ad873016c1b21cbe17eda47070
d9f150b951adb997f75430d76b054a1f551a251729bcee3dc850a7afd0d054ef09981dc896882f23e752b65d865ab023f3951f9c3bc6215862148265
0102cca7ba4e61c8527623d9f2142e9584a4a682c91cb43952a179f1182895bc77ec4f3c8084069f15a92660deb367b4cb2ecc696d89e8ab47e46584
4abdef7fff0317e01883605906df43d6c92cb008102e959c34becb49b6ea9d7ba43a792527d41bcf4d97ed5486ad81db15224cb85472d1d0ec195d75
ff4b993ffc249fff310b731689112a9beab6afd5e60daa96188c609939a93cf85cb473925af36a1b4c8c319b01867c94bc76cf7292bdaaa1922361f2
424ea6b63af9c9f94be3d93655844a0f245cbe0ff4cfff259ffd975a42a654e6efe5f3ff6344ebac55bec69c532cf3e17d5e7c87e8878c4a102cf3e7
f3e0354bb879af92797abeafccdb3eb52a394f254d629779634fad301972f5a7bd38a11f9e4752957ae36941db221bad90d954f28b180e950ba96309
cf4d653ee4b6b7efa51f6f6c5e479a58c483a10896f9f3793049bec7a05c997f74eef0cf5631ce5b27da273939253f6c2e27d2a32648b6e68d85a8cc
b4bdca4995ed050390cafd56e08babc607998fdc0e9a59040d3bea2015f997d685d45f9ea7792c623b1f68b908884e2adb3ef6216f102c3188fd12f0
ca654a16a07cf7f4e8afa57f5e36ecd075a061dcae5e657f67dbd4fa2af6b0af84dc5dd91f7e2fc289739528f3c9ba92615ddf73e6e69d4c15a265ca
c6fdc4eff94c3eff57d2d83d2a790df6e43b30e571c0fa9cc282c7917b982ff27eaaa08c7c112c3366ab081ee757bec9713e9c0cd7dd4e7c825dc556
105e2404347fb635f359e7e776dae7b5d9fad377f37b5f5e535ca0537ab06c4815f1dcb9230daff2796f3efb4fadcf7be35e3efbed9180e673af2434
6e061eba1e23a90b2af5f70e59bee6484efe19477e089699f27c65fa75a23941de78aede623d1b1a7653ad66d4122c1bd2a3925d53e27990cf7f9217
e3d236ccc7c5f6ab2c9c043a638e65be7c5d25cf720f954ab681b4f39b42ec44813f34f3283719224b87ad9ac9e73fc4fed1f8db838ceeec243ec2e3
ebfcb041cf5ff445b0cc97af03475113b4ed047639c1bef7dc7ea376330265fa5a01930b2cbf2e65614eea81f22799ef79eddc314cf28b3a91168265
866c1f364f2b7c2fa798741f830d3f3284f3862ace28677242ddca2c50563d7c37778115a2ff6b0d5e65518ece749f7b5f81f0033b35a10f82659e7c
1d308abf1295162d079c647bb10b8eb6e5849add67442ebc68f0fcf705d68e7cfe3f53c5efe4443efb9b29acf41e4a36ba38f1f4cfd17a089d112c33
232d667cb4d6f85adbae287327d96dcf3ddf7277dd0a93bba9edbed49584ca9b3c1e6d5c320f79b3153243355bcfcd6b2b4cea925641cb73f1f13e33
1c8ece0896f9f1f1057f2861c1ce1836384980d272a2fd5a5935c79e6c0e5b27d3bd5cc3648350d99d84ccadd645566dd345ec45e5c75665b2e4963a
3ece192ce2416704cbfcf818926058638e0dda4d35a7b5f8a1a461f3bbb920b9253b751441fa39122a07908bac83d6677f5b86504bbad0ba94c567ba
55952f727ef93c8f0b79386fa013fa5866c4534526b986beb990f0b297c02e1fabbcb67643c9767e581f133609afa679746b979b83849b95abd66e40
17742ff89b4ca17a76eee82fcb8d341017c132233e9a204b45021ec94977477609d909bc4b4eb39bc96df367cd07fa488dd097e122edefddaf9acffe
4ee0e0396bede6735b4bd5d107adb51d12ff36f29f62271eac45b0cc88d67aec9b7558d2f027ea3671a86c7c24dc20175aeba791adeaec861ab41ec2
4accb1cc845c6d8ef140a84429a47a3f75a8b46e64281ec8c1d8cfea5b9992052c45b0ccc7d80302273f1421c17db1cfa57a0a244daaeb633b00d07a
082b31149e010f8b76ce726c740dccd35a5f25bc788aef1992e769210f8b78b01415cb3c8ca936be72b24309b4d6b789afc8ff22d55420591208bf8e
7c7c9c53b01415cb0c8c5cb4c3821d644f4265c8d5f63eddd93e897cea9032adf5fd8815fc76cefe8e732baaa7a858a66fe4a29d3b422572272b5973
0995d607396903291b3312f68e453c58866099be31db68b1600759f3d01e652aefe4b103499206f26316f2707ec1420c85274c1a0fff18f808594880
6cc902031bcc36327f176d43fb2d163a205563a65ab1e10616a16299b6a157842cd841b60a0a954a9ec393ecce04a4e870e863f2d05f1905a26299b0
1157929f6bd9271a659100f658e8dbba6d8c61781cc919b19087453c7050b14cd488dd3cae0995c851e1a1d27a64c1031235f47cc3221e380896e91a
fa45676802d991f9c42587ca065b402239b290e76ce0e31ab3c0140562283c412316ed9c1863d85a0e59f1b0b3548ee82f8be468ad5f86cc6d66110f
daa858a66948456346a8446e2a0d958afdc591a841235e54e1d146c5324103af1a59b083acc8c9e8bcf277edd218c34919c918b8cb15bb4de1372a96
89d15aef0d08952cd841560895bfed6badf9ee2225432e743ec8142e806099a0214311543c900d19022654feed93548980c9494bac210b79380fe127
86c21332b0ddca5763cca9732b90e667dc86cabf9c3ba0649e348dd491842153b258c40345c532397daff866844ae4426b7d4aa85ce9addd5f5c761e
02a6d6bb02c9221e282a9669197085f8d118c3101a92a7b5b6ad75f679a73a617f712441e6ff7eeaf15858c4032a96a918b068e79250898cb040a53b
7b1ce0bb8d14f49df3ff817df141b04c47df2fb05d4d6ae885871cd8ae0532ffea92376c255bad7ccffecb48812ce439e9f950180eaf1cc1327f7f49
c07cd15ab39d2392263d1bb72540e14f76379e4dd95e0f989cccf7ed3bb4cd79a87204cb74ec8d3cd9dae1b36fad90c9fead4892ad82d800653b1af0
0efd64a7b568b678440a6c98b4732bedb94429f53ca059fa86eca8854ab1782731728578e471f5eccc0e4d301f132992cfbbfd6cbeabf00db2dfcd1d
16e9606af23db4d3aabe787828b4cdaa1c15cbc4d8938cddf35b2a18764eda47bbd26ec4a37c6bf762964ae63dbb232025f279b79fc9f795bd31b6a3
032bbf31293b47bf5599f4112a0f0995205826ce561a6dfb8656d0fc3a62c8dc56857eb44226070024c1ce2bac6471cf997c971941c024ec5c7c992e
653c8c8cddc94592663a071a0c85674c82e1b187fe80f6e0b047f504299061b9a7017be6a7ecc12e82e03b8629c89cfb0b0fdf295bd43866630eac42
c53263b208e2a055cdfc2cf3b6fab293b39fa59279c1ce1f98920c8fdbcfe061016fc4eff641844ac46417d0c84e4eb67af47d44a8b4a308db729ed9
2454621d2a9685f2b408c80edbd13a0293d25adf0e58999a824386061193cca1bf18b918ce16278e6cef59e71ea003826525e480733ae2047d621715
39b70211c8e7f73693e1f16b630cedbe10854c89ba1879f1659ba09f5255870f04cb4a69ad0f2468f63d5133c70693d15a9f7a5abd1a027b7c230a19
91bae8b98f77db9d1cc7594406ef089618b308e8557a643264826812ed7df999ef01421ad96b928200a22158c221bb265c480fccae68c48ea864a5eb
f7895ff54bd9a61208c2f69a1c3857fe52c2e493730f1010c1122bb516011df518367f9090c99ec708ce6e3f376248708cffc918f3af7987e19bed35
29d5c93e5395b8b84712683784955a3b016d4a4ba3f71d76026a3762bfa5113b42928532ef47eeb53fc47fe28d852fb602df6a5cfead43a87c954537
6fa415d016a112292058a217d921a5bd13d0e19ade9976a5e2a384cc2b7a642204f95c6eca8916c882f49a7ce9d16bf24efaa2363d258f592c86d410
2c318aedd32757ca3664be59b325dfa75623f65342267c939658f673f87ff1e22245b67556ab71f9cd9a30f92afd509b2d1377996284d4112ce18d0c
9bb77702fa28f32d17f9d20a99f4c7843752c1f9b7bca248859d0ea4b5be9730f963cdc2c8f99d6e68b28fac102c118c9def235bd935d5cc9325f3e0
fe9280f92293d601206b764446e698db30f9b8a23d96bdf8fed8aa4a1eb0921b392358228a058b80b6172c02b24342df5a2193dd4b006443c2e49584
c9e725bbe1cc2fbad961d10d4a42b0c424ec15f99a454036647e6f85cc5dde2900a9913079da0a938b5a5fddcd0d6fb3e806c5225822090b16019db5
1e970d99371232ef65df6800988c9d1bde0a93f3bbe1cc162cba61781b5520582239326c7e34b708a819366ff7c8bca747268058ec1cf0567ba0f9dd
702ee77a4ab2e8065522582279b208a83d6cde2c027ad7ea91794bfb2200bead685cbe68d10dc3dba81ec112d959b008e85a26c93fd3881dc058d2b8
fc69ae71f9cf4537ad20c9a21b60018225b2268b80f6e61601edb41bb1f30e0358471a97dfb71a97bf958bd63f16ddacf96780ea112c5194058b8094
044c1ab103f883342ebf6d352edf9c5b74b3c7a21ba01f82258ab56011d0ad9c4468c40e546aaed7a46d5c7ecfa21bc01f8225aad15a046487b44e65
85e7138dd881b2cdf59abc524a9db6aa92472cba01fc2158a25a365c4a85e24a86c48e69c40e9443bed336341eb5462f7659740384f30f5e5be0d722
20a514733081ccc93417bb80ef5816dbf0bd062222580200b26787bb655e359d20800931140e00c81ef3248134102c010000e005c1120000005e102c
010000e005c1120000005e102c010000e005c1120000005e102c010000e005c1120000005e102c010000e005c1120000005e102c010000e005c11200
00005e102c010000e005c1120000005e102c010000e005c1120000005e102c010000e005c1120000005e102c010000e005c1120000005e102c010000
e005c1120000005e102c010000e005c1120000005e102c010000e005c11245d05aff1f5aeb57adf526ef280000d32058226b5aeb2dadf57f564afd77
4aa9ff4629f57f6bad7778570100888f60896c69ad8f95528f4aa97fd67a0eff504afdd05a1ff0ce020010d73f78bd911b19eebe574abd5df1d0cf6d
e5d21873e4dc03000082a06289ac68adf79452cf6b4265e38bd6fad6b91500000441b04436b4d6ff5e29f5bde7e3fda0b57e726e050000de31148ee4
d9053a4aa9ff4d29f55f0f7cac6fb5d62f4aa92d63cc8b732f0000f0828a2592a6b53e92053a4343a5913f37ec10ba845400001000c112c9d25adb05
3adf463e3e3df7ff1fb5d6bbcedf020000a3112c911cbb9a5b6b6d2b8def023db61b59040400003c22582229d29bf24784c7f49d5e970000f8c5e21d
244386be43552917b1bd2e378d31a70bee0300003d51b1c4e4220c7daff24daaa4000060248225261571e87b95bf089700008cc750382633c1d0f72a
365c2a630c0113008081a85822ba8987be57a1720900c008044b44250dcfa71efa5e85700900c0400c85231aadf5adddbb3b8357dc86cb2763cc8573
0f0000588a6089e06c4b1fa5d4936cab988b73997349b80400a02386c21194ec70f39c59a86c9cb3fd230000dd112c118cd6da361eff9ef92b7c43b8
0400a01b86c2114462ad84c6b2e172db18f394f7d30000202c2a96f04a6bbd95682ba1b1ee65ae280000588260096fb4d6074aa9c7425f513b47f4de
b9150000fc46b084175a6bbb7afabcf057f3adb44c0200000b102c319aedf9a894daafe495fc20211a0000cc215862b0d67ccab795bd8afbb2831000
006821586290c2e75376f18d36440000fc896089de2a994fd9c50d2bc50100f81b7d2cd14b61fd297db0afc756fe4f030080f1a858a213994ff942a8
74bc65310f0000bf102cb196cc257ccc74bfef18f665ce2900005523586225adf5b19d4bb8eaefe0a773adf50e2f0500a066044b2c25cdc0ff5a763f
1c344f070a2217d6007a2058c261573a4bd3f30fce9d586583f996405998e602f443b0c41f6438f7b9c2a6e7be30df1228c7954c73a1f303d011c112
bf4920fac12b321a2722a000c618db4eec4eda8a01e88060899fb4d6a7343df7eaaaa0e702d4ec56a6b9f09d063a2058a259a4f38557c2ab77ec270e
14a19937fd89692ec07a04cb8ac9229d1716e904f38d2171206fc618bb90f1419e04d35c80350896956a2dd2a1e979580c9f01f96b7f8ff94e032bb0
577885b4d67b4aa9efb5bf0e91fc1c1237c69c56f16c8132b57bfaf29dee40766cdb911f5be5dd941f5b01be8af1fa4901a5f9fdbbadbbb6a4f3c94c
1e4fc3becf7614efde18435fe281b43126cb078e61a4e12f4dcfe3db9621354410e9737e628ca181762564da507b84a7baefb404b5cd56586b82dbd0
91af3363cce8b9e832f7d5164c3e3977fa63bb031c1338d7a362591159a47864af610000200049444154c37cca695ccc5d3103c88b1d02df6f3de22b
095559496cc4ca2e1aed1d2ce5399c46eeb76ccf9d375aebe6ff7b09c525628e6525d84967721f58518a3eec2e4e5a6b937b7701dbca4c9e47ee43c7
f3732b73edfc90e5e7c986495b35b69f2509c6536fe2f1453ed7f77621ac736fc5180a2f9c7ce09f58a4938457630c07a0084a180a971368dba53126
bb8b93529e875afc5cac37c69817e7d6442d790e9331c6e865bf5b86deed68cf3be7cef43c1863b2ab608740c5b26032799a95dfe9d828a06a830864
a86fdebe5448b20965720c9ab79f7105f3ceb9e5d7906c4e4e127aac8b5e4ffbb9399200fc239350a9a4826d3fd78b3ef3552158164a4e3e37b5bf0e
09fa421f3c74b0eae4742e4382395447563d8f2f199e88172ddcd8cfe9394895fd8dad1c3b77c6f7fbf594beca171228bf25f0d886baa97dda13c1b2
40320cc8f68ce9a26a8975d685463b0af14316e4a5ac4be0b227e25cf6e25ef67a67f59db643f7763a821d865ef4a3947a6f17a7d8e93bce7fecd7ad
04ca7b195ddb0ffcfb62395f32ea5005826561643f5bda09a5ed53cd071d74d275a1dd874286df9a61c4a4bf172b5acdbc2ba94a658cb17d1c8f644e
f847e72ff8732a813297e1ee3ebed7baa887605910b9ea0bd9c70bfe50b5844f3789ce59ecdb89e2bb5c1ca76ce1bc40dbe3d0b9a50012a6cf023d93
b181d2be175f6df85d54799dff914a6ccc290017ce2d15205816a0b5e77789577da57a5bfb3c1c7867e72c3e155025f9247348537d1ecb9aa2bfcdbd
35d40a2984fdbb050172d7eee0b3a292fc07a9c4fe9e021021647eaa714e3dc13273f2a165e5779ea85ac237dbdbef39a193d9d0397a1bf23c525ca0
b46a3e283b31f975273b1c3521d2eb9c62697bb5eddce1577505048265c6645ed563edaf43c6360aae70609c8791fffd6322f32e5785b02e7e2438ef
72d573e23bedc7652b4c2eab107b21ff7ec870595dfb218265a6e4e0453ba1fc51e1c0223e1a6e97d2f6e47b4acfa343d58ceff470771228a3bedf12
2e0f9d3b3008c13243d24e28e73e5ff81b150e2ce26bc8ef7ce25056caf398b76c018f92ef34f3a7fbfb6c2b9453fd7263cc4584f64a552058668676
4245a2c281793ee7924dd953cff7f34825b0adab28f39deec72eca4961815095abb87d235866847642c5a26a893ff85ea420c3c9d1ab41019ec77926
7347dfb2b55f6787013e2743e5d2a83f6904cb0c483ba127da09158d608979ab865b87b899a885cfb573cb383709ac7aef1284a85aae772d43d0a908
ba50a81604cbc4c901f449da88a05c6fd98d0773429c70a7a8c894f23cfafa50630fc39e68b956208265c2a487db233d2aab41d512bf05aae4d80b98
a81522993be77b51c4c6943bf4f418ba6511cf0a090d81c3238265a2647ece8fda5f87ca7c605e16e69c38b78cb73f41753cc4b0f0a78917f37409cb
5c2c2ec70aec42112c13448fcaaa51e1405ba8a1c2d855cbd34041e27cc2ad1fbb0cc76f70b1b8140b650a45b04c0c3d2aabb75fc05ecff0c418f312
a86ab9117b483ce06296d4e7e951b5445508960991033d3d2ac18908bf19638e0355fbf663eec52d55cb9973c778519f474bd7f981b4884355089689
d05adb83d47eedaf037e627538e6859a2211bb254ea8cf76d2554b3a3ea02604cb04488fca0fb5bf0ef8ed1df3b2d0262bab2f03bc289f22572dedbc
ba33e78ef1a658f8d667453373a7510d82e584a4f1f90b3d2ab100150efcc11873106828396ae831c61c057a1e294f2161381cd520584e4416683cd1
a3124b7c597c332a17a22a37c5672dc4f3483abc310a815a102c2720bb313c132ab10a2722cc33c6d88bd143e78e91627fd6e4797c74ee18ff3c521e
72e6fb8c2a102c239303f863554f1a43b13a1c0ed991c7770ba2e8814c765df11d92a33d8f01bbc6102c5105826544b23290c6e7e88a795958485a10
f90c9793849e002139e545902cd04415089691c810cdf72a9e2cbc61381ccb780e97932d20f41d9227ea69d909df67d480601981eca6735efc134508
ac0ec7523e43d994a1c773b84c79e72a82258a47b00c4c6b7dca6e3a1881fe77584942d9d7557f2707f23c7cccb94c39bc112c513c826540b245236d
6330c606c3675847b64b1cbbca7af2cf99ccb97cefdc510ee659a27804cb40d8a2111e112cb196ac527e13685ff16864779e37819aa84f8e0b45948e
601980844aae4ce10b272274628c7931c6d83986d739bf62f23cb6026dff08202082a567ecfb8d00f83ca11763ccde80a1f1be7d198393ed1ffb0e8d
3f39b7a4850b45148d60e991844af6fd86770c9fa12f3b346e8cd14aa9bb9c5f3c3b342ecfa36b159660094c8860e981ddf75b6bfd42a844409c8c30
88316657aa7e2be75e0ed849262aa9c26ee7fe3c3215a237e807d9de3825211e4f8acf332882e5483654ca1532fb7e232482250693aadfe68a763e59
ccc9b47b8ccbf358363c1ead3a3b621421f9a92d3608d9fecb76144e6b6d9452df9cbfe4c7a3fdf7b5d657b2335dece7698b4207765d843ccf50fda6
9be7792bbf2fe55eaba3112c4790ab906742252248763711e4c3b6f39161e5f980799ad9f3b86f3d8f7605f3c2f9cb58c8ee50645be269adef25f4fc
feb14148fa2fc71a85b3dbd77e9f7f1cf2d82ec64e05923077d504e5d6f37c9630192bec7f90dff7bce0b9dad0795a42e8d4c618e7467478e17e85ca
47e70e209c3776b52cafef7ab2db55e88d094ea4a977b664fb431b36b3be7091e771252bc963fd4e1b766e9c3bbaf938d5907de6a36cdbb66aeddcba
84d6fa2860b535a43b99c292252a9603c801855089d8a85ac22ba9fc65ffb992e751d53cb6118e321e65eb5b91ce31542a999749b0acc5c8ab54600c
e65902e9c8f5fb98f3a847df2a6fce1d11b27d9f08963d102a31312a324019a60ca539cf43ed1b2c73ed10f0203b5065e91fb93ef0d8ece4df802bc6
802e0896403ab2ac58da79da5a6bbbaadeceb5bccf61de76332c3c605eeaa984cb971c829aacddd8caa017eb4a04cb0e089548043bf000182db76ad8
d0854e129ab3a95acac2a4ac43a562287c3d42250060012ef4800508962b48cb12422592c1d68e0080943114be44a43e780080cc70811786bcae5b73
3f0d3b27f49d34c49f1fcabf9755d43f8792d9d6735a04cb05089548d86ec62b1d815214bd255f28d29cfd407ede0dfc351b0ba621fcf1ffb5d6ce7f
a4949ad926fab22140b62bae73c050f81c422500600d362be840b6527c9adb42f1db88503986dd9ef28b52eac7dc568a17b21a1b9e102c5b08950080
0e180a5f40f61fbf6d05c9f388fb8d0fb56f77d26b05cd22f6eb9e12c152d8ab16422532c0090d981e154b61e7453655495b0d5c304c9d1b5bd57c96
907945c8ec8f60f977a8dc77ee0000e0cff3c566c67b6d7b615f03095d4676a34bbd2a39d4a756c83cc8f329c4577db0245402007aa8b65a2943dd2f
3257f293f317c6b1abbd2f9552874aa98fc618bdecc7deaf94faaa94ba8ef8f4cf9b3999ce3df843d5c192500900e8c9c77494ac3a3bb402e50f8fd5
dabbb900b9698c3930c65cac6b1764ef37c69c1a63f69aff5e29f5462975e6fc65fff60998ab551b2c0995008001aa99e72c43def71e03a5ad4abe97
30b8ebb3dfa4ddbed118732421735b29f5e0fc25bff619225facca6049a804000c54c5568ed225e5d9536ba026506ec6e82169f7dc36c6ec48153374
c0b443e4f72cf2f95b75c1925089ccb11a159888af1d7752de1946aa944f1ebba49cc50a94f3a48ab9237332437a278b7caaefdaa16a0b96844a1480
1d2380e9ec95fcdadbb99452a5f4b5cafbd00e4f3bb7462673327584eae50d43e315054b42250060241f15a93be7960448b5ed87c747622b95492d70
91eae5a573875fe7b587cb2a8225a1120030866cfb37c55684c149a8bcf1f87b6629542a17b12bcf099761151f2c099500000f7c0d832735bf5286bf
7d864a2bc950d988182eabdc83bce86049a804007852dcc20c59c9ec3be8be1a63ae9c5b1323e17216f85165d5afd497628325a11200e091af9d6652
0a1b1701b6a7cca97178e88b85b7d2b6a92a45064b42250af6c49b0bc4e5b38d4c2aad86b4d67b01b665543975aeb0fd2e23ecd6e3ab6d53368a0b96
727540a844a90896407cbe82e5ab73cb7442cd83ccea1815639191d63ae939a7be15152c2554567775000008cad7c29d24aa79b260a78a1d843a3a09
fcef17ddff745e31c19250894a5439191c988a84305f6d8652f9fe060b3a29ef2ab4c2e9f2bbbca82ac417112c09950080407c86b0544257b0452b39
6e6b68b77e0cddb8bea6ed1eb30f9632778150895ab0a5231097cf40c01ce974311ae449d6c1523adb7f73ee000a2557d600229006d7be86315f6515
720a420ecde65a9923587a926db0945079eedc01942bc93d868182f9dc962fa5d18690c7922c779bc9746e6892b20c96844a548a6a6577556ea506ef
4a9c5f195ace2ba0b978f720bb602913600995a811f32bbb23586214cfabc15562c132e490fc464d0b55baaaa9229a55b0940fabefcdf2815c305403
c4e3b5f29658b0087d915a55dfc60eaaaa8466132c651235a11235a36209c4e373b794d48245e8fdbcbf38b7d4edaaa6679f45b09450f9e8dc01d463
c68a70200e191ddbf0f8cb920a1691fa36866e3a9e0d634c55af45f2c1925009fc44b51288c7f7dece2956ac8e9d5bfcfa22e7ef9c8468c3f4d5b9a5
7049074badf5262754e027be07e961ce6b81e4bcf3c9e333bb4ba87fe56f32e7f3dab9c3af6c8680032d387aa8ad5aa932a8583e791e8e00724588e9
6727a7078ba4f8ae56a6fcddb5adfb5e9d5bfd7997d190b8eff75d65dc2c7e946483a5d69a5009fcf24af3dede387660a81a86c17f92b996a157707f
91ded3c992d6523eabd4d676adf3e2930c965a6b7b127debdc01d4895009442001c8e745891d0a4d7a1a8b5cb41e3a77f8759e78b8f47d8cfd98e2f4
8758920b96122a43ee630ae4a6aa561563456cceccbcd7f25453ad6c33c65c440a97a1170cf516607474bbf611a6a482a5d6fa82500938089609a2fd
5359b4d67b9e77da5111fa457a23e1f263e05ff397d6fa5e16484dca3e06adf58bc7d151db124ed75ca96c24132c6582efbe730750b73b024c6f2cdc
c110beab9549ae065f452a6d6f6c485af1d7c6b2e1fd596b1d62b14c27322cffecb15279628c611b599144b094f2389dfa0117f32bfb8b510da96a8b
b6d2c9f409dfa365d9542bdbec85ac84a413e74ebfbe69ad4dccb99776918e5429cf9d3b87b12beadf1863921be29fd2e4c1523e547f39770050b99e
9c2656658b0f8ce23d18c8d072b6242cd9eae543e0e7702e01f32ad410b91d11b5bf4329f5c36395f2b33166931125d7a4c152ae127d5d3900a5c96e
282d11312a9654920b21732b7d572b4357fba290eae54ea480f94986c86dc87cb245a72141d3eef66387d9e5df3012287d8e881eca5c4ae6be2f615f
9cc5f784fec5bffa46fd70ee00d0f85ae3ae0d63c98924b41386bfca6017930458b4b35de245a104bdab0a17d9da21efddd45b47a562928aa5ec1f4a
a804562054f617b1d51015cb02c8542cdfa1f2b2d49106a960da80a5a53d51c85d7b5270267328370995dd450f96ecff0d7472c9cb34489415e1b5f7
a92b48888bb72a2e08ed1c52095c36647e2d28645e4a98b423ba47cca1ec6f8a8ae53ddbad016bb168679818c1b2f42a4d15a41b89ef73d1758d952d
3bbad20a99db995d18cf5af326edcf0161729ca8732cd95507e8e44126cca3a740f3e5e6d94555ac3ccf984cc77a0cf00cde3364fa2719a5dc933ea1
a1bf9bebbcca1cd10b461dc2f947ac5fc4ae3a4067542b079013588c131727a4fc8518aeaeb25ab98e54ff2e161dd7644ef4ae8c34ec78da05672623
a3f67b7acb7b125f948aa5ecaa430374a003194e424fd236e67b84d7ed23d58e7c05fc9c14b9121ce82bf81c4bd9b6895009745344ffbb89ecc5f8b5
84caec3995330f8a5d090ef415b46219b1820094e20d13c787b10d913d0da5adc2fcca8cc994acfd00cf80ef2d2082cdb19406e8844aa0bb334e4ec3
c85cadd0a15231bf325f52e808112a4ff8defe22dfc366bee496fc2899f37895fb36970dc9373b323f744b76fbb23ff6e2f6aaf61ec4412a963289fe
d9b903c02accd11a485ac7fc15e15731bf32535aeb9700ed8566c6982de7d6b25eb7262436616a73e442dccf296e8728cf73b315169be73cf433630b
0547cead1508152c630c490125a9f620e443a436432cacca94d6fa4af6a2f62dc990b44cc0d7a18fe0d34912799ed51e2fbc2fde915e95844aa01ff6
9d1e482a0d31da0cdd39b72079b26d638890719d59a8dc49216c85263d4a8b7f9e29f31a2ce955090cc2dcca71a2ac0697c6cac888848cf3408ff8c0
b9256d9b957c768b9e9a90036fc152e6388598180d148d21f0d1629de0695c9f9f50f3610f33bc18bcaf613b529903cdb6ab13f2122c65a821c6c479
a0345f7947878bb81afc8eaa725e649e5d88cfc6758eab9bede7d7eee76de785caee34c592e7f9b1f4e799aad1c132f0500350b259ed6d293c8855ed
a55a991119410b31cfeed518136bea4510765ea85dc96e17962cfa91e099fd7c625bb9ecf03caf9dff10a38d5a152e6d859e02b470006a90d58ad2d4
c46c6bc66af07c04de98a39a765301cfefc96d3210aaab04abc287b92754028364b5a23451b1aa9597ce2d48928ca0850a952735f53095a91fb58ca8
30cfdda3c1c1525680d356081826b715a52962181cbf4985ed3ed02b62ab6c35b604ab2248b3e9815f8382252bc08151beb210641c59301863b4e481
934e36428da0bdb23f3cd05def6029f35758010e0cf3c0821d2f62558f78af322073e4428da0ed38b70058aa57b00c3c7f05a841d62b4a5320d5ca18
d3705e736c2b531bd9ed2dd4ce4b87ecdf0ff4d3b762c9901030dc574e525ec4aa56b2cd66e2245486daeded840b0ba0bfcec13260b359a006d70c81
8fa7b53e8a751ce2fd4a9b9c934285cacb4a17eb00a3fda3cb3f10b0d92c500b5681fbf12dd2ef39716e413202572aed0a70beafc0406b2b96b2651a
8b7580e13eb30a7c3cad75ac0ae22bd5aa74050e9533568003e3acac584a5fb01be70e005d9dd1087d3cadb55d99fb25d2af2354262ad40e29c286ca
2de75600bdac0c964aa9ffc5b9054057b6b5103b3af8116b1105fbb7274a6bfd14707eed2b6d85003f96064badf57f544afd4be70e005d31a4e6812c
d80955a59ac785406202ee59ddb0a1728be92a801f0be7586aadff0da11218e52327aaf1a4776eac053becdf9e1879ff9f0995403e9c60295f645644
02c37d651b406f62063d5602274476797b0cf888ecb4874d4225e097132c9552ffab730b80ae2e99a3e787ac028f3504cefeed099116772177797b60
a10e10c6cf6069e7b0d82fb2d6faff534afdb7bcd6c0200ff4bff343aa55b15681b37f7b42a49d50c81677b64f250b758040fea1b5b6ab2df7798181
515e3959f9218b354256abe6b17f7b02222cd251d2fe8b055a4040ff241596192f32300ac36afedc47fc5d87ecdf3e3dadf541e0453a4aa63b102a81
c09a7643bb8127490325db667e9e1fb2ff7394bdc0651578acfe985822f04e3a8d8f2ca803e2f839c752aed80f79cd81de3e52f1f24316eb7c8af4eb
ecd40586c027647753d25a9bc0a1d2b6137a43a84c4a72538664eb6a78f27b55b85cb99ff1c2029d1d72c2f2435601c75aaca3d865655a7211f123f0
83b8a39d5092424e77188a8b4c8ffe683724f34fae737f524004270ca3fa21f3eb42ae029ec7bcca8948079297081711763e255528ac25bdbb635ed4
16cfe96329c3430fb5bf30c00a36541e2fbf1b5d49a83c8ff8829d7141300da94a875ea0f32a739e691f95b0c4869e1975f2cc0996ea57b8dc61a538
b0d025a1d28f0942e51dab82e3b31521a95286ae4a3743df54a3d33779b06c55cf632d16ac8636c62c7cae917a8a0139b9a401ba1f1384ca193badc4
17b14ff267f6790f431658f96617cf6d4ef89c6cb0bd71eef0cc18a3633fb7142cac58aa5f2fc88bf4e67b75ee04ea43a8f4648250f9ca629db8ec7b
2c812474a89cc9aa6f42655e366401575452a5bc8f112a6bb634582ac225d020547a22f3ec6287ca2d5606c7d11af68ef11edb053abcb7e1dd05fa0d
5fe4223338696d752f737cdf25f2ba166b65b054122ea564cd9c4bd48850e9890c8bc65cfd6ded103cc26b55821e234c9f6aaa942cd08923e4e29673
392e78271739b75239ff3145a0acb53fe6da60d990f949ac16474d08959e48e88831d7ae6d9b851ce149308855093aa44a195de82e0afb36fcd963c4
d02026173676fac595fc5b462e7242efe8b44e95f3ba972ede5946b65c8bb53b0630155a0a79203de2ee27580448a80c4c6b6d57d87f8bf4eb1ea45b
09261071115669aa5c34d8b962d9903e9727ce1d4039be122ac7d35aef451a1a6d7b255486d55a98132354daf7f33da1725a327253ea74b8cb80abb7
df4a05f54a2eb2abd0bb62d99093c677e70e206f8734d01e4f567cc6decd82853a014db09a9fef6242a405e173414fe98fa94e81da2a7565a7191e97
d2dd6070b054d30e7301217c64efef7126ec7f4b9fca40260894cc6d4e5421fdad175eb04c1c2cdbb21f71e93d14de669fbcac180fd58e0088a1193e
25548e200124f4967d8bdc112afdb3f3eae4641b2b543ec86a6f4265a25a5d62ae337be877f2d9d28b426562b2fffc8f0a960dd9ecffab730790be99
0c9f32276f0459f51db3aad53893e30ffcbc8f9bcdcada888b357ecfa3641a431e64adc576e23dae1fe47365c3e46e469fadeccf45a386c29d7f8ca1
71e4e58e5032ce04c3a46d0b87b4d09f6d20ad94ba8abc6fb20d25f6847fefdc836cc8f0f85502ad7d9454528f86140ab4d64f09ec1b3ee95697be78
0d960d5a1320a2ffa294fa67037e1ded84469093c9fd8407e2f70492f122b70c6a10280b25c5a5e388e7ff3b59f4327a1a93e7efc283541eefe5cf9f
ffbb968a7c9060a9febe02fee1dc014cef337b0b0f37f18523fd0c479293ffd5043b9110282b241d6476e567e867ee4142da55c863b75c301fcb3cc7
f6c8eb4cc2e1ad3c8e7ba64f2d172c5836a85e2221afb2c51f078401261ef656329ff2c8b9159dc83eedb1b7d454120af6f8de0175081e2cd5dfd5cb
5be65e6242cca71c2891ef2f43df0348b5e862a2f7ee4e02250b72808a785915be8e3d21c88454568e630a2784cafe6ca0d45abfc89496a942e583ac
ea24547664f75bb60b116465f7f709debbb30c57e202f0244ac5729eb426893dbf0775a2d2d55342230ccc85edc88649a5d4e984c755e64f02f8294a
c5729e4cbe7f9f780f2ce48d4a574f89542855ab5136a172053bccddaa4cde4c142a2fe5bddae4bb06404d152cd59fc3e387ce9dc03827ac1ceece2e
ca4924502aa952d2287b09bb00c7be57ad61ee29da3dbdcaf6a7f6c2ed80f70a40db2443e18b24b0e214f963d5770f89756c6071d502322de13491e6
d3f47e05b05632c1b2417b220c742ddb8c6105092a1709cd71666e5e4bab8fde17e7ce69d89d4ca84a02e82cb960d92060a2878f3e765e2899d6fa34
a1b0d2a87e4b46099247f2934a3bb63b099354fe910cf9aefc9ce2c4f13e6dc906cb0601132b307cbac2c43d0c57b9b473f356dc5fac09b6bceb6af0
1ecbc0503282628f057b01e70bcf64b71c1b466f191d092ff960d92060620ead681690b6331713eee1bdca830c7bd7b15feeaf0acb9e04c914df8f33
d96799616e04d7baa8da4be862f755b6370dba55646db209960d16f9548f2ae59c842b938de203a584c803f949b547ef4caa929c401145e217baabdc
c9451743ee0364172c1b6c135925e652fe392f6f8a7d9ffb282e50b6aa90f6e793f317d243551251c977e42293ef47572c62eb21db60d9900ff12d3b
f914adda39790db9f23f4ea4edcc3ad9074ab970dd950099c36bdeb89420c95c494455c1b9f8ab31e6d4b9158e7f38b764464e5e3f578a25baf215c3
55db9732e1451eab6475012061bdf9c9293cb6112431b9886b2066f2e71443eb6cbad151f615cb4518262f42155787f6b36a57296656919c975cdb20
09e64dd57127e3e0d83693cfc8154372488154299f3c9f6b475d2cc9b1f42040d0657e7f474506cb36aa98d979287d3b46093d4772f0cbf5e2e741e6
1c0d6edd21afc39673c72ff633b0e9dcaad4bf90dbff85fcb79b855e40ce649eda05d548a44802dc8da787166cb4431ee7958fe384ddc6d4b9118ee2
836523c11d47f0a72277609115db7b89b5d818e35a7ac26db57e725bf1999aeb56cb132a91489e9c4f7f78789cd1aa803e46320996dd54132cdb3268
cf529bac776069ad14de2d2840c23f0224b2e771f87b926d78b5d6572356acbfa7c1fa7a5506cb36adf571066d5b72673f648baef4b258ecd1da4a2c
f7851e48c35d6b1790eadb67212f1e17ea4cd63e6e44b8a4e55d07d507cb364266108b426512ed6864ee4d131a9b1f867531a5d726744a6593f99548
865c643f7b7a3cdb537ebeb5d643c20fc1b20382e5125aeb235981c9b0e630ffaf52eabf9afb2fff1fd9f9e3df75f917572ceed89c6bfdd0feff9bcc
a345a19aede76e194ec714e4bcf8cdd3af9e740ad4c042125b097740b0ec402a5ba70416008962fe26820b30aa3775d5f2bee779fdc41873ecdc8a3f
fc132fc77ab6f46d5be0c88ab037f6c325d5030048819d2f766e8729ed109ffcbcd8766bb21a16f0c1f70aee47a9824e425adb9df0c9f08b8ae54899
ee9002a04ecd70fa29ab5bd157e0750867c698c94226fca16239922de3db95cdb69a2915cd6dd939000052b32117c13f16543617cd670662f9d2fa4c
5ef179cc1715cbc06415dd81ecb4c28a630039b893aa260b15f09b84bdc7095e910759e7c01ce20c102c2750c8967e00ea3193fd9bb3ddc8007e8c6c
30eed3acb5608d16400921582642c2e69e844d569f0348d9ab543459215b9909ab965d35f3882f089cd32058268ec089c2ddcd3dbd27f9196bbed7a9
92ffcf08411854342ba2b53e902e0439b996b0c9f48ec008961993fe9a6c3388d81e24fcd955c52fcd9f25ae3296563d9bf21d6bc2ea16f3a5d7ba93
cd1058795ea84cc3651b55f74008960593e0d9deae908a2716994938fcfdc35682c3cced2bcf36a1bf70022f945c78dd163212602f860e38f68d47b0
447beb44aa32e599b5f69ebea7829486ca471baea59ac909bc10b65d956d1754d053bab46d049d5bd109c112dec8c9b26dd95edf63d970f4dfcbbcd3
7f5ef93bf8daec1d6dffe4649dbfd6bceadd4456df86f42021934516052860787c9efd7ceed2e2a81f8225b2c00e47bff782263c564c6bbd27a173af
c085483309992caec89c4c09b92aa81a4f05b3078225922527d1d38a86e45f5b7dd938b9a233a914ed1554e17c95f96e7c0f32274581d3423e9bef99
4eb41ec1124909bc176d4aae5b211e9f3caa000013d049444154926116782753530e0aa8f2db90b9c7707919e473799c7135f390b65aab112c312919
32392d7c88fb9266bd48818c021c657c527f904a2655a34264ba13dd67aae9cb112c119d1c482e0a5d0d4b884436e4bb782027f6dce66c5e4bc8a4e2
5f182938340bd8529d4fbccd7cf7c5089688a2d030f920fdf9181641115a27f4e3cce6367f35c69c3ab7a238092d60bb36c6ec39b7826089700a0c93
97b26d1d57a9a886cc7bcea5a23993f9980c955764c2799b6fa898bbfec9b90518c1563cb4d6575a6b7bc5f29871a8b40b064e64b843cb0fbb32a03a
76c71c63cca6fd0ed813a952ea2ce1d7c056597fd8e38f34ed8607f27a2efa79d25a5f2ce8611c959d7a648cd96d8ed5b6822dc7f0d0a8582e40c512
5e14b0f3c2ab2c223ae50a14e84602c569e2dbc53e4815938bc281a450d0c583316627b1c77e117071e8095b95baa8586230db3b4f6bfd22079d1c43
e565ab22b92995194225d091548a7612af66dad0fb28d5369a5c87f54e6b9dd4c245696cbeeddc81600896e8456bbda3b5be9730799ed94a525bb9f8
c8d036e09fbd2833c61c4d301cd9c7b9044c16dc8593dcf42739ce7f74ee4010044b7422f3686c98fc91f8b0d7bc139960ada5b2421b202002bb4abb
3537f3a32cac49c5be04cc7b59090f8f52ac0ccbb1ffceb903de112cb1949d3fd51aeacea581f96cae2ac9f03630311932df4a3064da8be467099849
cd0d4c4cdf4056cbc2292adf0b102ce16855276f3219eabe6ecd95dca22a09a46b2e647e4e64b8fc9dac267f21607ab191da5ccb00ee984ab518c112
3fc9dcc9a78caa9367ad216e567c0219b2dbe2b586cb539893b941c0f4e6839c5392986a20efa7cff99f2c045b8260593959d9ddcc9d4c7da78d93d6
10f71143dc40399a3999b2bafc72e22746c0fcd3d0eae35b996a3069089370ebb382fa9562c67204cb0acd35313f4ffc15688749fa85018593d5e507
52c57c3ff17c4c02a61fe753b57b92dff9ec715ad725db87ae46b0ac48d32a48be649f127ee6ed616ec2245029bb35636b3ee69455cc26602633b41b
99af6a5f13309f42864c299e9c06289e9c485f4cacc0ce3b15902beddbc417e2dcc9ee180c6f03584a76fbb99af878667be2eed672bc92d7fcc6b9c3
af3b394fd99ffbaeafadd67acbbe17f2b317f073f19e3de8bb2158164cae08531eea9e4998e4cb0aa097d6bcb929fbea5ed652c1eab1ad6369ceec9c
fe4a9ffb200c8517a8b52027d550f9b5d51a885009a037998bb933f130f97e455b45a6b68b526897cd42d1b29fa67f04cb82241e28ef5af32699f80c
c09bd6629fc3895ed5f3d217f8c88afd14b7e9f4692643de9ab994c3112c0ba0b53e4e3450da03d0a17c49ab998f04601ac6988b090366b3c0a7d851
98f6369dadd7f9c1f98b79b96bedd6c6289a07ccb1cc58c27328ab9ad80e204d131f23ed4575755bfe69adf75a0b6952ec8d6ca74d5cb0435b3804cb
0c251c284f680f042035131e33eda8cd0ecdb4ff26d305eccf56eb4f9f0bb0ec6b7edfac2eb77f52e4888b609991482d1ffa7a95ea24c30700923661
c0ac66f53840b0cc40a27d2819ee069025adb51da2de8ffcd8b908471558bc9330d93de05ef6f14e25545ecaeaee1d4225801cb55691c75c78d22cee
a96ede25ea42c53251135d51afc2500e80e248a3f5a7c817efccbd44b1a85826a6d58b32955079424f2f00a59246eb365c7e8ef8146d887db4fb593b
f70099a362998889ae9a57a9b2550680ba4d305a3493ea25538b50042a96099003d97322a1b269684ea804501d199d7913718719dbebf159fa3f02d9
23584ec8b60f4a68d8fb924009007f0c8f7f8df8727cd75a5f39b7029961287c22b2dadb6753d8a1589403004bc834a5db88c76b5b29dd62681cb9a2
6219596b71ced4a1f28e453900b09a542f7722ee3fbe2143e3bbce3d4006a8584632c155ef323436078001e4387e1f710fec3363cc91732b90302a96
11c8a4ece78943a51d5e794f6373001846aa97766feb93482fe11799360564838a65605a6b5ba5fc30f1c3a075100078a4b5de92ea658c6e1eccbb44
36089681c8fede3f267e18761e25f374002090c8c5838fc6985be75620210c850720bb294c192a9b616f4225000424c7d9586d896eec0250e7562021
542c3d9a6062f7225f8d316c13060011451e1ab75bed1e3bb7020920587a22ad216e267c08acf606808945ec514c0f622489a1700fb4d6c71387cacf
acf60680e949cfcbb3080f645fe6770249a16239d2c43be83cc8410c009010990b791ee111cda40512900482e540329ff229d27c9a456c95927d6501
205111bb83102e910c86c20790f994cf13854a3b97f20da11200d2668cb1235a6fa45347486fb5d64f7c1c900282654f5aeba309e7539e30971200f2
21bbf5d811ae59e0074db84412180aef41fa537e99e057dbab5d1b2839680040a622355367581c93225876a4b5b643cf9f26f8d5ec9e03008588742e
215c62320c85772057995384caaf844a0028873166cff6a00cfc84ded28a0853a162b986cc599962279df732f11b005018adf585ed4519f859d1441d
d151b15c61a250399355df844a00289404bed095cb7dd9c00388868ae5121385ca6b192601005420d29ccb4363cc85732b1000c172818942a56d25c4
9525005426d26a71a657210a82e59c894225bbe80040c5229c7b6cdbba2dfa202334e658b6c8bedfb143e57b422500546f27f00e3d76a738568a2338
82a590a18877ce1d616d33340100904a62e8f672ef64a30f201882e5df6d1f42cf6f697b9595dfeca40300f8490a0d87815f8d2f5a6b16892298eae7
584a2b86bf9c3bc2619e0b0060a9483d2edf701e420855074bb96afbeedc110ea11200b05684c53cb4b74310d5064badb5dd47f5d1b9231c422500a0
9348e728fa5bc2bb9a83e58bac928b81500900e8456b7da494fa16f85563481c5e55b97847763a88152a15a11200d09731c6aee07e08fcc2b14a1c5e
55172ce50a30f4f6596ddb844a00c04007815f38bb9f78e83647a8485543e113ccabfc688ca1212d006030e93df925e02b3833c66c39b70203d4162c
636ed7c8a4680080175aebd0276bce59f0a29aa170b9e28b152a2ff98202003cfa1af8c564ae25bca8a262a9b5b67bb0fe70ee08e3c118b313fb3902
00ca1661d4edab2c180206ab2558de47da079cb642008020b4d67621cf79c857d718a39d1b811e8a1f0a9755e03142a5b547a80400842053ac5e43be
b872ce04867f864aae586aad379552cfce1d619c1963f84202008289d034fdd518b3e9dc0a74547ac5f2d8b9258c0742250020b40873203764c81d18
a4d860293d2b43f6fd6a63237f00402c97817f0fc11283955cb18c55adb4abe89e9c5b01000823743bbb0fd24d05e8adc86029d5ca7de70eff66b466
0000c4243bbacd02ff4a46e23048a915cb58d54abe7800802984ae5ab26e008314172c23562bedee3af7cead000084771bf837d8453cbbceadc01a25
562c635d657135070098840c8707ed69c9a81c86283158c658097e46237400c0c4ae02ff7a82257a2b2a58c6eabd45cf4a004002420f87bf657538fa
2aad621923589e38b70000105f8c79feccb3442fc56ce9288b761e9d3bfc7bc330380020055aebd027f16b630c43e2e8aca48a658c6ae525a1120090
90bbc00f85a170f45252b08c714545337400404a420f87db79969bceadc01245044b19067fe7dce1d71d7d2b01008989318a46d5129d9552b18c51ad
0cbdcb0100007d855e19ae58c0833e08961d196308960000002b94122c3f38b7f8157a72340000a9a26289ceb20f9691f632a55a0900481173ff9194
122a9631261587de360b0080de22b5c00b3d2a888294102cb79c5bfcbaa377250000c07a542cd78bb1e20e0000207b2504cbd0257ae6af0000921469
9d01d059493bef8442b0040000e820eb6019e34acd18f3e4dc080000000715cbd5e85f09004819db2d2229b907cbd0154b5683030052b6c9bb839450
b15c8df995008094c558bc33736e01962058020090af18154bd61aa03382e56af4b00400a4ec1def0e5242b004002043117b5852644167044b0000f2
142b5832148ece08960000e4896089e4102c0100c88cd67a33c296c63f1963180a4767044b0000f2b317e911b351087a21580200909f83488f986a25
7ac93d58f28107005425e63038e759f445c57235f6600500a4e628d6e3617e25faca3a5846f8c0b3072b002035b186c12f9d5b8035a85802009009ad
b55db4f336d2a3bd706e01d6282158865cb116ab471800005dc41a067f65181c4394102cef9d5bfc618e25002009b28563ac453ba7ce2d40072504cb
17e7167f36623f19000096385e7c73100c836390128265d0527dc44dfe01005848e656c6aa565e1b63d8c61183641f2c23cc0121580200a61673689a
61700c56caaa7016f000008aa4b53e8ab812fc8e453b18a3946019f24b106be80100008731e6d418a3ed8f52eabd52ea4c293573fea21f31e771a240
f6839afdb392799037ce1dfe7ce40a0e00902aadf58e344e1fd3e7d2562b19a5c32845044bf5eb4bf5127015f7893186ab380040567a06cef7c69890
2dfc50819276de09d91a21d6f6590000786383a231e6c818b3d50ca7cb90fa67d9b2f1557ed725a1123e9454b164381c00006042c5542c25f4859acc
aca85a020000ac56d250b80a3c1cbeefdc02000080df08963d482f310000002c5054b0942da8ae9d3bfc21580200002c515ac55205de8aeaadd69ab9
960000000b14b32abc4d6b6d5b26bc73eef06366db364cf4d4000000925562c55205de928aaa250000c00245562c55f8aae5ab3166d3b9150000a062
a5562c55e0aae586d69a2d1e0100005a8aad58aa5f554bdb34fd8373873fdbb2121d0000a07a25572c5584f6402157a003000064a5e860291bea5f3a
77f8f389853c000000bf143d14ae7e0d87db4536cfce1d7ebd31c6bcc47e6e00000029297d285c49e03b74eef02be856920000003928be62d908dc7e
c83a34c61030010040b56a0a9676b79c47e70ebf58250e0000aa55fc507843025fe821f15be7160000804a54132cd5af706987aaaf9d3bfcb1db3d5e
c57e5e00000029a86628bc4d6b6d17f46c3877f873628c61671e00005095aa2a962d7bce2d7efda5b50efd3b000000925265b034c6d8b9905f9d3bfc
fa2e0b86000000aa50e5507843e6437e72eef08be6e90000a00a55074bf52b5cdad5e26f9d3bfc7935c66cc67e5e000000b1d53ac7b26dc7b9c5af0d
09af00000045ab3e58ca30f57be70ebfde122e010040e9aa0f96ea57b8bc8fd03cdd864b1aa803008062112c85344f0fbd52fc03e1120000948a60d9
628c39b5cdcd9d3bfc225c0200802255bf2a7c11adb5ad5eee2fb8cba79931863e970000a018542c1730c61c28a52edd7bbc62410f0000280a15cb15
22552e5f95525b3451070000b92358ae21f3213facfe5b5e6c1b63a8600200806c3114be86316637c2b0b8f5a8b5de756e050000c804c1b28348732e
ad1badf591732b0000400608961d49b83c8bf0abbec9dc4e000080ac30c7b227adf5b152eaaf08bfeac118137a1f730000006f08960368ad6df5f23c
c2afb22bc67758d403000072c050f800b2fde3c708bf6a4316f530ef120000248f8ae5085a6bbb73ce63a45f77272bd40100009244b01c496bbda994
bab73be944f8750c8d0300806431143e92dd3147f6fcbe8bf0eb181a070000c922587a22c3d427917e9d6d49742fd552000080243014ee99ec9e7313
f1577e36c65c39b70200004446b00c402a894f32741d033d2f0100c0e4180a0f40e65d6e469a7769bdd35a1be9af09000030092a9681c9429b6f117f
e54c568ebf38f70000000444b08c6082a171ebcc18c3ea710000100dc13222adb55d64f329f2af7d6f8cb9776e050000f08c6019d904abc6ad07a5d4
2ec3e30000202416ef44668cb935c668097bb1bc534a3d6bad2f6a799d0100407c04cb89487ba0cf917ffb3eabc7010040280c8527c0eea22355c5d8
987f090000bca162998089aa97d60fadf58bd69ae6ea000060342a9689d15adf2aa53e4cf0a85e65810f154c00003008c132415241bc8ddcf7b2c10a
720000300843e109b25543d912f2648247d7ac2067881c0000f442c53271b26bcfed448b7b1443e40000a02b2a9689b343d2b2b8e7bd84bcd8365a8b
7cf66a7f3f0000c07204cb4cb486c70f277ac436607ea70f2600005886a1f04cc92e3afb133ffa4b630c21130000fc44b0cc5802f32f1b7625f901f3
300100a81bc1b20009054cebd018c39ee40000548860599089fb5fcea31f2600009561f14e415a0b7ca65a41ded6f4c3b48b7d8e9c7b01004071a858
162cb10aa635534aed31171300803251b12c5862154cebadf4c4b455cc5badf596f337000040b6a8585624c10a66e35a5695331f1300808c51b1ac48
ab82b92d8b6b52f1a9351ff34a56b9030080cc50b1ac58626d8a16b9534a1d31271300803c50b1ac58b30fb93146db5d74127c253eb4e664beb09524
000069a362893f48783bcfe055b9966ae693730f00009804c1120bc9429f2b59c99d3abbe2fdd4feb000080080e9102cb196d6da6ed1b89fd12b45d0
04006002044b749670bba22e2e2568b21008008040089618446b6d2b825f327ef5ec2e401754350100f087608951a465d195ace0ce5d13362f581404
00407f044b782343e51709f7c51cea5ac2f315d54d00009623582208adf5ae84cc1c56950ff520734e2f98bb090000c112115412324f8c31c7cead00
0054846089a8b4d65b12324b9893d9f68661720040edd8d21151d94531c6985dd946f28dadf449dfc99c9d112a0100a0628984c890f97186d5cc6d56
91030040c5120931c6dc36d54ca9681e4a0ba0945d122a0100f8858a25b2a1b53e908a664a8b80deb3221c00805f0896c8960c9d1f29a53e4df41cee
6c85d5b91500804a112c510cd905682f6255f3a31dbe776e0500a052044b14ad15368f3cef08f4608cd9716e0500a062044b5449b69f3c92d0b931e0
35f86c8cb9726e0500a062044ba04502a75d24b4bba2c23933c66c39b7020050398225d091844e5be1b44dde2f78dd000068514afdff0a065c21cc50
e5d10000000049454e44ae426082
END FIGURE
